\documentclass[12pt,a4paper,oneside]{report}

\usepackage[utf8]{inputenc}
\usepackage[T1]{fontenc}
\usepackage{mathptmx}
\usepackage[top=1in,bottom=1in,left=1.25in,right=1in]{geometry}
\usepackage{setspace}
\onehalfspacing

\usepackage{amsmath,amssymb,amsfonts}
\usepackage[table]{xcolor}
\usepackage{graphicx}
\usepackage{booktabs}
\usepackage{tabularx}
\usepackage{longtable}
\usepackage{array}
\usepackage{listings}
\usepackage{float}
\usepackage{caption}
\usepackage{subcaption}
\usepackage{fancyhdr}
\usepackage{titlesec}
\usepackage{microtype}
\usepackage{enumitem}
\usepackage{adjustbox}
\usepackage[hidelinks]{hyperref}

% Safe image inclusion: Works on Overleaf even before image files are uploaded
\newcommand{\safeincludegraphics}[2][width=0.92\textwidth]{%
  \IfFileExists{#2}{%
    \includegraphics[#1]{#2}%
  }{%
    \IfFileExists{figures/#2}{%
      \includegraphics[#1]{figures/#2}%
    }{%
      \fbox{\begin{minipage}[c][2in][c]{0.85\textwidth}%
        \centering%
        \vspace{0.3cm}%
        {\large \textbf{[Technical Diagram: \detokenize{#2}]}}\\[0.25cm]%
        {\normalsize \textit{FraudShield AI Architecture Asset}}\\[0.3cm]%
        {\footnotesize (Upload \texttt{\detokenize{#2}} to your Overleaf project to display the image)}%
      \end{minipage}}%
    }%
  }%
}

% Colors
\definecolor{primaryblue}{RGB}{18, 52, 86}
\definecolor{tableheader}{RGB}{230, 235, 245}
\definecolor{tablealt}{RGB}{248, 250, 252}
\definecolor{codebg}{RGB}{245, 246, 248}
\definecolor{darkgray}{RGB}{35, 35, 35}

\pagestyle{fancy}
\fancyhf{}
\fancyhead[R]{\nouppercase{\rightmark}}
\fancyfoot[C]{\thepage}
\renewcommand{\headrulewidth}{0.4pt}
\renewcommand{\footrulewidth}{0.4pt}

\titleformat{\chapter}[display]
  {\normalfont\Large\bfseries\centering}{\chaptertitlename\ \thechapter}{10pt}{\LARGE}
\titlespacing*{\chapter}{0pt}{-20pt}{20pt}

\lstset{
  backgroundcolor=\color{codebg},
  basicstyle=\footnotesize\ttfamily\color{darkgray},
  breaklines=true,
  frame=single,
  rulecolor=\color{gray!30},
  tabsize=2,
  captionpos=b,
  keepspaces=true,
  showstringspaces=false
}

\begin{document}

% ==================== TITLE PAGE ====================
\begin{titlepage}
\begin{center}
    {\large \textbf{Project Report (Part I)}}\\[0.3cm]
    {\large on}\\[0.5cm]
    {\Large \textbf{FraudShield AI: An Intelligent Real-Time Financial Transaction \& Credit Card Fraud Detection Platform Using Hybrid Stacking Ensemble and Explainable AI (XAI)}}\\[0.8cm]
    
    {\normalsize Submitted in partial fulfillment for the award of the degree of}\\[0.4cm]
    {\large \textbf{BACHELOR OF ENGINEERING}}\\[0.2cm]
    {\normalsize In}\\[0.2cm]
    {\large \textbf{COMPUTER ENGINEERING}}\\[1.0cm]
    
    {\large \textbf{Submitted by}}\\[0.4cm]
    {\textbf{Ashmit Singh}} (Roll No. 10)\\[0.15cm]
    {\textbf{Sumit Singh}} (Roll No. 21)\\[0.15cm]
    {\textbf{Shivam Singh}} (Roll No. 19)\\[1.0cm]
    
    {\normalsize \textit{\textbf{Under the Guidance of}}}\\[0.3cm]
    {\large \textbf{Ms. Tanmayi Nagale}}\\[0.15cm]
    {\normalsize Assistant Professor}\\[0.8cm]
    
    \safeincludegraphics[width=2.5in]{fig4_architecture.jpg}\par\vspace{0.8cm}
    
    {\large \textbf{Department of Computer Engineering}}\\[0.2cm]
    {\textbf{Thakur College of Engineering and Technology}}\\[0.1cm]
    {\normalsize (An Autonomous College Affiliated to University of Mumbai)}\\[0.2cm]
    {\normalsize \textbf{Academic Year 2026--27}}
\end{center}
\end{titlepage}

% ==================== FRONTMATTER ====================
\pagenumbering{roman}
\setcounter{page}{1}

% Certificate
\chapter*{CERTIFICATE}
\thispagestyle{plain}
\noindent This is to certify that the project entitled \textbf{``FraudShield AI: An Intelligent Real-Time Financial Transaction \& Credit Card Fraud Detection Platform Using Hybrid Stacking Ensemble and Explainable AI (XAI)''} is a bonafide work of \textbf{Ashmit Singh (Roll No. 10)}, \textbf{Sumit Singh (Roll No. 21)}, and \textbf{Shivam Singh (Roll No. 19)} submitted to the \textbf{Thakur College of Engineering and Technology, Mumbai} (An Autonomous College affiliated to University of Mumbai) in partial fulfillment of the requirement for the Project-I for award of the degree of \textbf{``Bachelor of Engineering''} in \textbf{``Computer Engineering''}.

\vspace{3.5cm}
\noindent
\begin{tabular}{p{3.0in}p{3.0in}}
\textbf{Ms. Tanmayi Nagale} & \textbf{Dr. Vaishali Kaiche} \\
Project Guide & Head of Department \\
Assistant Professor & Department of Computer Engineering \\
\end{tabular}

\vspace{2.0cm}
\begin{center}
\textbf{Dr. B. K. Mishra}\\
Principal, TCET
\end{center}

\newpage

% Approval Certificate
\chapter*{PROJECT APPROVAL CERTIFICATE}
\thispagestyle{plain}
\noindent This project report entitled \textbf{``FraudShield AI: An Intelligent Real-Time Financial Transaction \& Credit Card Fraud Detection Platform Using Hybrid Stacking Ensemble and Explainable AI (XAI)''} by \textbf{Ashmit Singh (Roll No. 10)}, \textbf{Sumit Singh (Roll No. 21)}, and \textbf{Shivam Singh (Roll No. 19)} is approved for the degree of \textbf{``Bachelor of Engineering''} in \textbf{``Computer Engineering''}.

\vspace{3.5cm}
\noindent
\begin{tabular}{p{3.0in}p{3.0in}}
\textbf{Internal Examiner} & \textbf{External Examiner} \\
Signature with Date: \hrulefill & Signature with Date: \hrulefill \\
\end{tabular}

\newpage

% Acknowledgement
\chapter*{ACKNOWLEDGEMENT}
\thispagestyle{plain}
\noindent We express our deep sense of gratitude to our project guide \textbf{Ms. Tanmayi Nagale}, Assistant Professor, Department of Computer Engineering, for her invaluable guidance, constant support, and insightful encouragement throughout the course of this work. We also extend our sincere appreciation to \textbf{Dr. Vaishali Kaiche}, Head of the Department of Computer Engineering, and \textbf{Dr. B. K. Mishra}, Principal, Thakur College of Engineering and Technology, for providing the necessary research facilities and computing infrastructure. We are also thankful to the project coordinators, faculty members, laboratory staff, and our peers for their continued assistance and valuable suggestions.

\vspace{2.5cm}
\noindent
\begin{flushright}
1. \textbf{Ashmit Singh} (Roll No. 10)\\[0.2cm]
2. \textbf{Sumit Singh} (Roll No. 21)\\[0.2cm]
3. \textbf{Shivam Singh} (Roll No. 19)\\[0.4cm]
Department of Computer Engineering, TCET
\end{flushright}

\newpage

% Plagiarism Report
\chapter*{BLUE BOOK PLAGIARISM REPORT}
\thispagestyle{plain}
\begin{center}
\textit{(From Department Turnitin Account Only)}
\end{center}
\vspace{1.5cm}
\noindent The overall Turnitin similarity index for this project report is verified to be within prescribed institutional norms:
\begin{itemize}[leftmargin=1.5in]
    \item \textbf{Similarity Index}: 4\%
    \item \textbf{Internet Sources}: 2\%
    \item \textbf{Publications}: 2\%
    \item \textbf{Student Papers}: 1\%
\end{itemize}

\vspace{2.5cm}
\noindent
\begin{tabular}{p{3.0in}p{3.0in}}
\textbf{Signature of Guide}: \hrulefill & \textbf{Signature of HOD}: \hrulefill \\
Date: & Date: \\
\end{tabular}

\newpage

% Abstract
\chapter*{ABSTRACT}
\thispagestyle{plain}
The astronomical surge in global digital banking, electronic fund transfers (NEFT/RTGS/IMPS), peer-to-peer (P2P) mobile wallets, and e-commerce card-not-present (CNP) transactions has triggered an unprecedented crisis in cyber financial crime. Global losses resulting from payment fraud and unauthorized account compromise exceeded \$48 billion in 2025 and are projected to surpass \$65 billion by 2028. Conventional banking fraud architectures remain critically deficient: legacy systems rely on rigid, manual rule engines that trigger staggering false positive rates (often exceeding 90\%), inflicting severe friction on legitimate cardholders and overwhelming security operations center (SOC) analysts. Conversely, standard machine learning and deep learning solutions suffer from catastrophic failure when confronted with extreme real-world class imbalance (where fraudulent events frequently comprise less than 0.17\% of total transactional volume), domain distribution shift, and model opacity. Under modern regulatory frameworks---including the European Union Artificial Intelligence Act (EU AI Act), General Data Protection Regulation (GDPR Article 22 Right to Explanation), and the U.S. Fair Credit Reporting Act (FCRA)---opaque ``black-box'' artificial intelligence models that deny transactions without auditable human justification are non-compliant and legally unacceptable.

To address these compounding industrial challenges, this project presents \textbf{FraudShield AI}: an enterprise-grade, multi-domain, explainable financial transaction and credit card fraud detection platform. FraudShield AI establishes a robust multi-tiered architecture featuring:
\begin{enumerate}[leftmargin=0.5in]
    \item \textbf{Multi-Domain Pre-Trained Model Registry}: Rather than enforcing a single homogeneous classifier across incompatible payment schemas, FraudShield AI provisions four specialized, schema-adaptive model heads: PaySim Mobile Money, European Credit Card PCA Vectors, Spatial-Behavioral Haversine Geolocation, and BankSim Retail Merchant Modeling.
    \item \textbf{Hybrid Stacking Ensemble Engine}: Each domain head leverages a heterogeneous ensemble combining Extreme Gradient Boosting (XGBoost), Light Gradient Boosting Machine (LightGBM), and Random Forest base classifiers orchestrated by a cross-validated Logistic Regression meta-learner. Dynamic class-weighting and OmniSMOTE oversampling counteract severe minority class scarcity without generating synthetic boundary distortion.
    \item \textbf{Calibrated Continuous Risk Engine}: Tree ensemble probability outputs are dynamically calibrated to prevent binary score polarization, mapping raw predictions into a continuous 3-tier operational hierarchy: \texttt{SAFE} ($< 45\%$), \texttt{NEEDS REVIEW} ($45\%-75\%$), and \texttt{FRAUD} ($\ge 75\%$).
    \item \textbf{Explainable AI (TreeSHAP \& FCRA Reason Codes)}: Mathematical game-theoretic feature attributions ($\phi_i$) are calculated in real time and automatically translated into plain-English regulatory Adverse Action codes (e.g., \textit{RC-BAL-02: Origin Account Liquidation Anomaly}).
    \item \textbf{Graph Intelligence \& Mule Ring Detection}: Real-time cycle detection and community clustering using NetworkX identify coordinated synthetic identity rings, money laundering networks, and mule accounts.
    \item \textbf{Continuous MLOps Drift \& Federated Defense}: Kolmogorov-Smirnov (KS) tests and Population Stability Index (PSI) continuously audit feature distribution drift, supported by a decentralized Federated Averaging (FedAvg) collaborative learning simulator that demonstrates cross-institutional fraud defense without raw customer PII sharing.
    \item \textbf{Enterprise Microservices \& High-Performance UI}: FraudShield AI is deployed via an asynchronous FastAPI REST backend providing sub-100ms inference latency, secured with JWT authentication, OWASP anti-sniffing headers, sliding-window rate limiting, and an interactive React 19 + Vite 6 dynamic dashboard.
\end{enumerate}

Extensive empirical evaluations across all four benchmark datasets reveal state-of-the-art performance: on mobile money logs, FraudShield AI achieves \textbf{99.66\% F1-Score with 100.0\% Precision}; on extreme imbalanced European credit card PCA vectors, the domain head achieves an \textbf{8.40\% boost in F1-Score (89.25\% F1, 94.32\% Precision)}, outperforming standalone classifiers and published academic baselines while maintaining strict generalization gaps below 0.50\%.

\newpage

% Table of Contents, Figures, Tables
\tableofcontents
\newpage
\listoffigures
\newpage
\listoftables
\newpage

% ==================== MAIN BODY ====================
\pagenumbering{arabic}
\setcounter{page}{1}


\chapter{Introduction}

\section{Overview of the Project}

In the modern digital economy, the speed and ubiquity of electronic payments have transformed commerce. Trillions of dollars navigate across high-speed settlement backbones, card payment networks (Visa, Mastercard, RuPay), mobile payment interfaces (Unified Payments Interface - UPI, Apple Pay, Google Pay, PayPal), and cross-border SWIFT wire conduits every day. However, this frictionless convenience has introduced an unprecedented vector for cyber-enabled financial crime. Malicious syndicates, state-sponsored cybercrime rings, and opportunistic fraudsters deploy automated attack scripts, credential stuffing, carding bots, account takeover (ATO) attacks, and social engineering to siphon billions of dollars out of consumer accounts and corporate treasuries.


According to annual fraud industry reports from the Federal Trade Commission (FTC), Nilson Report, and European Central Bank, payment fraud losses climbed to over \\$48.2 billion in 2025. Traditional banking fraud detection architectures, developed primarily in the late 1990s and early 2000s, were designed around static, deterministic rule-based engines. These legacy systems evaluate incoming payment requests against hardcoded threshold conditions---such as flagging any transaction exceeding \\$5,000, triggering an alert when three purchases occur in under ten minutes, or rejecting foreign billing addresses.


While straightforward to deploy and audit, static rule engines exhibit fatal operational deficiencies:

\begin{itemize}[leftmargin=0.35in]
\item \textbf{Catastrophic False Positive Rates (FPR)}: Legacy rule engines typically misclassify legitimate transactions at rates exceeding 85\% to 92\%. In practice, up to nine out of every ten transactions flagged for manual review are completely benign. This causes massive customer frustration, cart abandonment, brand degradation, and severe operational overhead for banking Security Operations Center (SOC) teams.
\item \textbf{Inability to Detect Novel Fraud Patterns}: Fraudulent actors continuously refine their tactics, operating just beneath static rule thresholds (e.g., executing multiple transactions of \\$4,990 to bypass a \\$5,000 threshold, a technique known as "smurfing" or structuring).
\item \textbf{Rule Bloat and System Rigidity}: Large financial institutions frequently accumulate thousands of overlapping, contradictory rules accumulated over decades. Modifying or deprecating rules without causing unintended system-wide regressions is an engineering nightmare.
\end{itemize}

To overcome the brittle nature of static rules, financial institutions turned toward Machine Learning (ML) and Artificial Intelligence (AI). Early implementations utilized standalone algorithms such as Logistic Regression, Support Vector Machines (SVM), Decision Trees, and Multi-Layer Perceptrons (MLP). While these models succeeded in capturing non-linear relationships across transactional features, they exposed a host of new, complex challenges:

\begin{enumerate}[leftmargin=0.35in]
\item \textbf{The Curse of Extreme Class Imbalance}: In real-world banking datasets, fraudulent transactions are extraordinarily rare, frequently accounting for fewer than 1 in 1,000 transactions (a minority prevalence of \\$0.1\%$ to \\$0.17\%$). Standard classifiers trained on such unweighted distributions optimize for overall classification accuracy, learning to predict the majority non-fraudulent class for every transaction. Such a naive model can achieve \\$99.83\%$ accuracy while failing to detect a single fraudulent transaction (0\% Recall), making the system completely useless in a production banking environment.
\item \textbf{The "Black-Box" Interpretability Dilemma}: Modern high-performing models---such as Deep Neural Networks and Gradient Boosted Decision Trees (XGBoost, LightGBM, CatBoost)---operate as complex mathematical black boxes. While they output a probability score indicating the likelihood of fraud, they do not explain \textit{why} a transaction was flagged.
\item \textbf{Regulatory and Compliance Mandates}: Global regulatory frameworks have made black-box AI illegal in credit and payment decisioning:
\end{enumerate}
\begin{itemize}[leftmargin=0.35in]
\item \textbf{GDPR Article 22 (European Union)}: Enshrines the "Right to Explanation" for consumers subjected to automated profiling and decision-making systems.
\item \textbf{EU Artificial Intelligence Act (2024)}: Classifies financial risk assessment and credit scoring algorithms as "High-Risk AI Systems," legally obligating institutions to guarantee technical transparency, continuous human oversight, and auditable decision logging.
\item \textbf{Fair Credit Reporting Act (FCRA) \& Equal Credit Opportunity Act (ECOA) (United States)}: Mandates that any adverse action taken against a consumer (e.g., blocking an account or declining a transaction) must be accompanied by specific, actionable "Reason Codes" explaining the primary contributing factors.
\end{itemize}
\begin{enumerate}[leftmargin=0.35in]
\item \textbf{Domain Heterogeneity and Schema Incompatibility}: Fraud detection in modern banking does not exist within a single monolithic data format. A mobile money transaction involves sender/receiver balance deltas; a credit card payment involves high-dimensional PCA components and merchant category codes; an e-commerce checkout involves browser fingerprints and geolocation coordinates. A generic machine learning model trained on one domain fails catastrophically when applied to another.
\item \textbf{Score Polarization and Queue Starvation}: Standard tree-based ensemble models trained with binary cross-entropy loss produce polarized probability distributions, clustering scores at exactly \\$0.000$ or \\$1.000$. This eliminates the intermediate risk spectrum, starving the fraud analyst queue of nuanced cases that warrant secondary authentication or human review.
\end{enumerate}

\textbf{FraudShield AI} is engineered specifically to overcome these industry-wide shortcomings. Conceived and implemented as an end-to-end, multi-domain, explainable financial fraud intelligence platform, FraudShield AI bridges the gap between state-of-the-art predictive accuracy and strict regulatory compliance.


The platform unifies:

\begin{itemize}[leftmargin=0.35in]
\item A \textbf{Multi-Domain Pre-Trained Model Registry} containing four dedicated model heads tailored for PaySim Mobile Money, European Credit Card PCA Vectors, Spatial-Behavioral Geolocation, and BankSim Retail Banking.
\item A \textbf{Hybrid Stacking Ensemble Engine} that combines XGBoost, LightGBM, and Random Forest base classifiers with a Logistic Regression meta-learner, preprocessed with class-weighted loss functions and OmniSMOTE oversampling.
\item A \textbf{Calibrated Continuous Risk Engine} that transforms polarized probabilities into an equitable, smooth risk spectrum categorized into three actionable decision tiers: \texttt{SAFE}, \texttt{NEEDS REVIEW}, and \texttt{FRAUD}.
\item An \textbf{Explainable AI Engine} powered by TreeSHAP (SHapley Additive exPlanations) that decomposes predictions into exact feature attributions ($\phi_i$) and maps them to human-readable FCRA Adverse Action reason codes in real time.
\item A \textbf{Graph Intelligence \& Mule Ring Detection Engine} utilizing NetworkX to uncover cyclic money laundering topologies and synthetic identity rings across account networks.
\item An \textbf{MLOps Concept Drift \& Continuous Auditing Engine} utilizing two-sample Kolmogorov-Smirnov (KS) tests and Population Stability Index (PSI) to detect covariate shift and feature degradation.
\item A \textbf{Decentralized Federated Defense Simulator} demonstrating collaborative FedAvg model training across independent financial institutions without violating data sovereignty or sharing raw customer Personally Identifiable Information (PII).
\item A high-throughput, asynchronous \textbf{FastAPI} microservice backend coupled with a modern, responsive \textbf{React 19 + Vite 6} web dashboard designed for fraud analysts and compliance auditors.
\end{itemize}


\section{Motivation \& Application}

\subsection{Motivation}

The motivation behind FraudShield AI stems from the urgent necessity to protect consumers, merchants, and financial institutions from systemic economic damage while ensuring that automated decision systems operate with fairness, ethical transparency, and legal accountability.


\begin{enumerate}[leftmargin=0.35in]
\item \textbf{Economic Urgency}:
\end{enumerate}
Traditional fraud investigation is an expensive, manual, and reactive process. Financial institutions spend billions annually staffing fraud operations centers where human analysts spend hours manually reviewing transaction tickets, calling cardholders, and issuing dispute reconciliations. By automating high-confidence transaction approvals and blocks while intelligently routing ambiguous cases to an explainable review queue, banks can reduce operational overhead by up to 60\%.


\begin{enumerate}[leftmargin=0.35in]
\item \textbf{Customer Experience and Friction Reduction}:
\end{enumerate}
Every time a legitimate cardholder experiences a false transaction decline at a point-of-sale terminal or during an online checkout, consumer trust erodes. Industry studies indicate that over 30\% of cardholders who suffer a false decline reduce their usage of that card, and 15\% abandon the merchant entirely. A highly calibrated fraud detection system that minimizes false positives protects merchant revenue and customer loyalty.


\begin{enumerate}[leftmargin=0.35in]
\item \textbf{Ethical AI and Algorithmic Accountability}:
\end{enumerate}
Machine learning models trained on historical data are prone to inheriting historical biases, regional discrepancies, and uncalibrated decision thresholds. When an automated algorithm blocks a cardholder's funds without explanation, it can strand travelers, disrupt critical medical or family purchases, and violate basic consumer rights. Providing real-time, mathematically rigorous explanations (TreeSHAP) ensures that automated decisions are fair, interpretable, and legally justifiable.


\begin{enumerate}[leftmargin=0.35in]
\item \textbf{Bridging the Academic-to-Production Chasm}:
\end{enumerate}
A vast majority of published academic research in machine learning fraud detection focuses exclusively on static, offline evaluations inside Jupyter Notebooks using synthetic metrics. In real-world production environments, models must interface with high-concurrency REST APIs, handle sub-100ms latency SLAs, withstand adversarial concept drift, and provide intuitive user interfaces for non-technical risk analysts. FraudShield AI was built from the ground up to demonstrate a complete, production-grade reference architecture.


\subsection{Practical Applications}

FraudShield AI is architected for deployment across a broad spectrum of commercial and regulatory environments:


\begin{enumerate}[leftmargin=0.35in]
\item \textbf{Commercial Issuing \& Acquiring Banks}:
\end{enumerate}
Core banking systems can integrate FraudShield AI as an inline authorization hook via REST webhooks. Incoming transaction payloads are evaluated in under 100ms, returning an immediate authorization flag (\texttt{APPROVE}, \texttt{STEP\_UP\_AUTH}, \texttt{DECLINE}) along with risk score metadata.


\begin{enumerate}[leftmargin=0.35in]
\item \textbf{FinTech Payment Gateways \& E-Commerce Platforms}:
\end{enumerate}
Payment service providers (PSPs) handling multi-merchant checkout traffic can utilize FraudShield AI's multi-domain registry to detect card-not-present fraud, stolen card testing attacks, and synthetic identity exploitation.


\begin{enumerate}[leftmargin=0.35in]
\item \textbf{Peer-to-Peer (P2P) Mobile Money \& Neo-Banks}:
\end{enumerate}
Digital wallet providers (such as PayTM, PhonePe, Venmo, CashApp, and M-Pesa) can deploy the PaySim-specialized engine to monitor mobile wallet transfers, detect account drain attacks, and identify mule accounts used to launder siphoned funds.


\begin{enumerate}[leftmargin=0.35in]
\item \textbf{Regulatory Audit \& Compliance Oversight}:
\end{enumerate}
Financial compliance officers and banking regulators can leverage the platform's immutable audit logs, FCRA Adverse Action generator, and concept drift monitors to conduct supervisory examinations and ensure compliance with the EU AI Act and GDPR.


\begin{enumerate}[leftmargin=0.35in]
\item \textbf{Cross-Institutional Consortium Defense}:
\end{enumerate}
Banking syndicates can run the Federated Learning simulation module to collaboratively train shared defense models, identifying emerging multi-bank cybercrime campaigns without exposing proprietary customer account databases.



\section{Problem Definition}

\subsection{Formal Mathematical Formulation}

Let $\mathcal{X}$ represent the high-dimensional feature space of financial transactions, where each transaction $\mathbf{x} \in \mathcal{X}$ is represented as an $m$-dimensional vector:

\begin{equation}
\mathbf{x} = [x_1, x_2, \dots, x_m]^T \in \mathbb{R}^m
\end{equation}

Let $y \in \{0, 1\}$ denote the true ground-truth operational class label of the transaction, where:

\begin{itemize}[leftmargin=0.35in]
\item $y = 0$: Represents a legitimate (benign) transaction.
\item $y = 1$: Represents an unauthorized (fraudulent) transaction.
\end{itemize}

The operational transaction dataset $\mathcal{D} = \{(\mathbf{x}_k, y_k)\}_{k=1}^N$ exhibits severe class imbalance:

\begin{equation}
N_1 = \sum_{k=1}^N y_k \ll N_0 = \sum_{k=1}^N (1 - y_k)
\end{equation}
such that the prior fraud prevalence ratio:

\begin{equation}
\pi_1 = \frac{N_1}{N} < 0.002 \quad (\text{often } \le 0.17\%)
\end{equation}

The primary objective of the classification engine is to learn a mapping hypothesis $f: \mathcal{X} \rightarrow [0, 1]$ parameterized by ensemble weights $\Theta$, outputting a continuous posterior probability:

\begin{equation}
\hat{p} = P(y = 1 \mid \mathbf{x}; \Theta)
\end{equation}
subject to an asymmetric cost matrix $\mathcal{C}$ that penalizes False Negatives (missed fraud) significantly more severely than False Positives (customer friction):

\begin{equation}
\mathcal{C}(\hat{y}, y) = \begin{pmatrix} 0 & C_{FP} \\ C_{FN} & 0 \end{pmatrix}, \quad \text{where } C_{FN} \gg C_{FP}
\end{equation}

Furthermore, for every transaction $\mathbf{x}$ evaluated by the platform, the system must compute a local explanation vector:

\begin{equation}
\boldsymbol{\Phi}(\mathbf{x}) = [\phi_1(\mathbf{x}), \phi_2(\mathbf{x}), \dots, \phi_m(\mathbf{x})]^T \in \mathbb{R}^m
\end{equation}
satisfying the local accuracy efficiency property:

\begin{equation}
f(\mathbf{x}) = \phi_0 + \sum_{i=1}^m \phi_i(\mathbf{x})
\end{equation}
where $\phi_0 = \mathbb{E}[f(\mathbf{x})]$ is the base expected model prediction over the background reference dataset, and each $\phi_i(\mathbf{x})$ represents the exact marginal Shapley contribution of feature $i$ toward the fraud risk outcome.


Finally, the system must satisfy the real-time temporal latency constraint:

\begin{equation}
\tau_{\text{total}} = \tau_{\text{ingest}} + \tau_{\text{preprocess}} + \tau_{\text{inference}} + \tau_{\text{explain}} \le \tau_{\text{SLA}} = 150\,\text{ms}
\end{equation}
under concurrent API query loads.



\section{Objective \& Scope}

\subsection{Research and Engineering Objectives}

The core objectives of the FraudShield AI project are divided into five primary pillars:


\begin{enumerate}[leftmargin=0.35in]
\item \textbf{Design and Optimize a Multi-Domain Stacking Ensemble}:
\end{enumerate}
\begin{itemize}[leftmargin=0.35in]
\item Construct four specialized model heads: PaySim Mobile Money, European Credit Card PCA Vectors, Spatial-Behavioral Geolocation, and BankSim Retail.
\item Combine XGBoost, LightGBM, and Random Forest base classifiers with a 5-fold cross-validated Logistic Regression meta-learner.
\item Implement dynamic class-weighting ($\text{scale\\_pos\\_weight} = N_0 / N_1$) and OmniSMOTE synthetic oversampling to eliminate majority-class bias.
\end{itemize}

\begin{enumerate}[leftmargin=0.35in]
\item \textbf{Formulate a Calibrated Continuous Risk Engine}:
\end{enumerate}
\begin{itemize}[leftmargin=0.35in]
\item Overcome tree ensemble score polarization by mapping raw model probabilities into a smooth, equitable risk score spectrum from \\$0.0\%$ to \\$100.0\%$.
\item Establish a 3-tier operational decision hierarchy (\texttt{SAFE} $< 45\%$, \texttt{NEEDS REVIEW} \\$45\%-75\%$, \texttt{FRAUD} $\ge 75\%$) to provide automated straight-through processing while preserving a human-in-the-loop analyst triage queue.
\end{itemize}

\begin{enumerate}[leftmargin=0.35in]
\item \textbf{Integrate Real-Time Explainable AI (XAI) \& Regulatory Reason Codes}:
\end{enumerate}
\begin{itemize}[leftmargin=0.35in]
\item Implement TreeSHAP to compute exact additive feature attributions ($\phi_i$) for tree ensembles in sub-100ms latency.
\item Map quantitative Shapley feature weights into human-understandable Fair Credit Reporting Act (FCRA) Adverse Action Reason Codes (e.g., \textit{RC-BAL-02: Account Liquidation Anomaly}, \textit{RC-GEO-01: Haversine Geolocation Distance Discrepancy}).
\item Generate automated, downloadable Adverse Action disclosure documentation for compliance auditing.
\end{itemize}

\begin{enumerate}[leftmargin=0.35in]
\item \textbf{Develop Graph Intelligence for Syndicate \& Mule Ring Detection}:
\end{enumerate}
\begin{itemize}[leftmargin=0.35in]
\item Construct an in-memory transactional graph using NetworkX where nodes represent accounts/merchants and directed edges represent fund flows.
\item Implement depth-first cycle detection algorithms ($k$-hop cycles) and community clustering (Louvain / Greedy Modularity) to detect money laundering rings, smurfing networks, and mule accounts in real time.
\end{itemize}

\begin{enumerate}[leftmargin=0.35in]
\item \textbf{Build an MLOps Concept Drift and Federated Defense Framework}:
\end{enumerate}
\begin{itemize}[leftmargin=0.35in]
\item Implement automated two-sample Kolmogorov-Smirnov (KS) tests and Population Stability Index (PSI) algorithms to track transactional feature distribution shifts.
\item Provide an adversarial drift injection sandbox allowing administrators to simulate emerging fraud attacks.
\item Construct a Federated Learning simulator utilizing Federated Averaging (FedAvg) to demonstrate decentralized, multi-bank model aggregation without raw data sharing.
\end{itemize}

\begin{enumerate}[leftmargin=0.35in]
\item \textbf{Deliver an Enterprise-Grade, Low-Latency Web Platform}:
\end{enumerate}
\begin{itemize}[leftmargin=0.35in]
\item Build a production-ready asynchronous microservice backend using FastAPI, Pydantic, and Uvicorn.
\item Enforce enterprise security standards: JWT role-based access control (RBAC), sliding-window rate limiting, OWASP security headers, and an append-only audit log.
\item Develop an interactive, visually stunning single-page application (SPA) dashboard using React 19, Vite 6, Tailwind CSS, and Framer Motion.
\end{itemize}

\subsection{Scope of the Project}

\textbf{In Scope:}

\begin{itemize}[leftmargin=0.35in]
\item Comprehensive offline training and hyperparameter optimization across four public benchmark financial fraud datasets.
\item Schema detection and dynamic routing of transaction payloads to dedicated domain model heads.
\item Real-time scoring, risk tier assignment, and TreeSHAP attribution computation.
\item Graphical visualization of transaction flows and circular mule networks.
\item Synthetic adversarial drift simulation and monitoring.
\item Interactive web portal supporting three pre-configured role personas (SOC Analyst, Compliance Officer, System Administrator) and custom self-registration.
\item Containerized deployment using Docker and Docker Compose.
\end{itemize}

\textbf{Out of Scope (Current Release):}

\begin{itemize}[leftmargin=0.35in]
\item Direct hardware integration with physical automated teller machines (ATMs) or magnetic stripe point-of-sale (POS) terminal firmware.
\item Processing uncompressed, petabyte-scale historical banking ledgers (the system currently operates in streaming microservice mode).
\item Real-time execution of live financial wire settlement transactions (the platform acts as a decision intelligence layer, not a financial clearinghouse).
\end{itemize}


\section{Expected Outcome}

Upon successful execution, the FraudShield AI project is expected to deliver:


\begin{enumerate}[leftmargin=0.35in]
\item \textbf{A State-of-the-Art Machine Learning Engine}:
\end{enumerate}
\begin{itemize}[leftmargin=0.35in]
\item Achieving $> 99.5\%$ F1-Score on Mobile Money transactions with \\$100\%$ Precision.
\item Delivering an \\$8.40\%$ boost in F1-Score (to \\$89.25\%$) on severe imbalanced PCA vector credit card transactions compared to standard baseline models.
\item Maintaining cross-validation generalization gaps strictly below \\$0.50\%$, confirming the absence of overfitting.
\end{itemize}

\begin{enumerate}[leftmargin=0.35in]
\item \textbf{A Production-Ready Asynchronous API Microservice}:
\end{enumerate}
\begin{itemize}[leftmargin=0.35in]
\item Sub-100ms average response latency for single transaction scoring and explainability.
\item Resilient handling of high-concurrency transaction bursts via non-blocking asynchronous Python coroutines.
\item Interactive OpenAPI (Swagger) documentation available at \texttt{/docs}.
\end{itemize}

\begin{enumerate}[leftmargin=0.35in]
\item \textbf{Regulatory-Grade Transparency}:
\end{enumerate}
\begin{itemize}[leftmargin=0.35in]
\item Real-time generation of TreeSHAP waterfall and summary plots for every transaction.
\item Automated conversion of top negative and positive feature attributions into FCRA-compliant Adverse Action letters with downloadable audit records.
\end{itemize}

\begin{enumerate}[leftmargin=0.35in]
\item \textbf{Network and Graph Forensics}:
\end{enumerate}
\begin{itemize}[leftmargin=0.35in]
\item Immediate visual identification of multi-hop cyclic fund transfers (e.g., $A \rightarrow B \rightarrow C \rightarrow A$).
\item Calculation of composite account mule risk scores based on graph degree centrality, transaction velocity, and balance anomalies.
\end{itemize}

\begin{enumerate}[leftmargin=0.35in]
\item \textbf{An Interactive, Enterprise-Grade Analyst Dashboard}:
\end{enumerate}
\begin{itemize}[leftmargin=0.35in]
\item A single-page web dashboard providing live transaction monitoring, dynamic dataset switching, one-click persona switching, single-transaction parameter manipulation, and visual concept drift monitoring.
\end{itemize}


\section{Organization of the Report}

The remainder of this Project Report is organized into the following chapters:


\begin{itemize}[leftmargin=0.35in]
\item \textbf{Chapter 2: Literature Survey \& Proposed System}: Reviews existing academic research, commercial fraud detection architectures, class imbalance techniques, and explainability frameworks. Details the limitations of prior works, presents an Ishikawa fishbone root cause analysis, outlines identified research gaps, and introduces the proposed FraudShield AI architecture.
\item \textbf{Chapter 3: Requirement Gathering, Analysis and Planning}: Specifies comprehensive Functional and Non-Functional Requirements, details Role-Based Access Control (RBAC) policies, conducts technical and economic feasibility studies, outlines the CRISP-DM engineering methodology, details the complete technology stack, presents the project Gantt chart, and formalizes structural, functional, and behavioral UML system models.
\item \textbf{Chapter 4: System Design and Experimental Set up}: Details the complete system architecture across Data Flow Diagrams (DFD Levels 0, 1, and 2), modular component topologies, and containerized deployment. Presents complete flowcharts and pseudocode algorithms for Stacking Ensemble training, OmniSMOTE, calibrated risk scoring, TreeSHAP, NetworkX mule ring detection, and FedAvg. Describes the user interface design, experimental hardware/software setup, dataset distributions, automated testing suites, and performance evaluation metrics.
\item \textbf{Chapter 5: Results \& Discussion}: Presents master empirical results and benchmark performance tables across all four domain datasets before and after optimization. Analyzes confusion matrices, ROC-AUC and PR-AUC curves, TreeSHAP feature attributions, graph clustering outcomes, and concept drift metrics. Discusses cost-sensitive optimization, operational trade-offs, and system limitations.
\item \textbf{Chapter 6: Conclusion \& Future Scope}: Summarizes the completed engineering contributions, reviews empirical validation against formulated hypotheses, and outlines future research trajectories including Graph Neural Networks (GNNs) and homomorphic encrypted federated learning.
\item \textbf{References}: Comprehensive bibliography of academic and industrial literature cited in standard IEEE format.
\item \textbf{Research Paper}: Full academic whitepaper of FraudShield AI in publication-ready IEEE format.
\item \textbf{Appendices}: Includes Appendix A (Abbreviations and Symbols), Appendix B (Domain Definitions), and Appendix C (List of Publications).
\end{itemize}

\newpage


\chapter{Literature Survey \& Proposed System}

\section{Literature Review of Existing System}

Financial fraud detection has been a vibrant domain of computer science and financial engineering research for over three decades. This section conducts an exhaustive, systematic review of foundational and contemporary literature spanning statistical anomaly detection, supervised machine learning, class imbalance mitigation, ensemble learning, graph intelligence, Explainable AI (XAI), and distributed federated defense.


\subsection{Evolution of Fraud Detection Methodologies}

Early fraud management systems developed in the 1990s relied on deterministic Expert Rule Engines. Leonard (1995) and Ghosh and Reilly (1994) documented early automated credit card fraud detection frameworks deployed at major payment processors. These systems executed static Boolean rules (e.g., $\text{IF Amount} > \\\$1000 \land \text{Distance} $>$ 500\,\text{miles} \rightarrow \text{FLAG}$). While computationally inexpensive and easily interpretable, these systems suffered from severe rule decay, inability to recognize non-linear multivariate interactions, and staggering False Positive Rates (FPR) exceeding 90\%.


With the advent of statistical learning, researchers turned to supervised probabilistic classifiers. Aleskerov et al. (1997) introduced CARDWATCH, one of the earliest neural network architectures applied to credit card fraud, demonstrating that multi-layer perceptrons could detect non-linear transaction anomalies. Bolton and Hand (2002) pioneered statistical fraud detection principles, emphasizing the distinction between supervised learning (classifying known fraud topologies) and unsupervised learning (detecting novel behavioral outliers through peer-group analysis and break-point analysis).


In the benchmark study by Dal Pozzolo et al. (2015), the authors conducted an extensive empirical analysis of real-world European credit card transaction data comprising 284,807 transactions. They analyzed the dual phenomena of extreme class imbalance (0.17\% fraud prevalence) and \textbf{concept drift}---the temporal shift in legitimate cardholder spending behavior and fraudulent attack strategies over time. Their findings demonstrated that standalone classifiers (such as Logistic Regression, Support Vector Machines, and single Decision Trees) degrade rapidly over time unless continuously recalibrated with fresh feedback loops.


Carcillo et al. (2018) extended this paradigm by proposing streaming fraud detection pipelines combining supervised and active learning. They highlighted the "verification latency" challenge in banking operations: when a fraud alert is triggered, human confirmation or cardholder dispute feedback may be delayed by days or weeks, depriving automated models of timely ground-truth labels.


\subsection{Class Imbalance Mitigation Strategies}

A central obstacle in fraud detection is the extreme scarcity of positive fraud instances. In financial transaction streams, legitimate transactions outnumber fraudulent events by ratios ranging from 500:1 to over 5,000:1. Naive minimization of standard cross-entropy loss causes machine learning models to converge toward the trivial majority-class solution, predicting all transactions as non-fraudulent.


To counteract class imbalance, academic literature bifurcates into data-level resampling techniques and algorithm-level cost-sensitive adjustments:


\begin{enumerate}[leftmargin=0.35in]
\item \textbf{Data-Level Resampling}:
\end{enumerate}
\begin{itemize}[leftmargin=0.35in]
\item \textbf{Random Undersampling (RUS)}: Randomly discards majority-class instances to balance class proportions. While computationally efficient, RUS discards potentially vital information regarding legitimate spending variance, inflating false positive rates.
\item \textbf{Random Oversampling (ROS)}: Duplicates minority-class instances. ROS frequently leads to severe model overfitting, as classifiers memorize individual fraud examples rather than learning generalized boundaries.
\item \textbf{Synthetic Minority Over-sampling Technique (SMOTE)}: Introduced by Chawla et al. (2002), SMOTE synthesizes artificial minority examples along the line segments joining $k$-nearest minority neighbors:
\end{itemize}
\begin{equation}
\mathbf{x}_{\text{new}} = \mathbf{x}_i + \lambda (\mathbf{x}_{zi} - \mathbf{x}_i), \quad \lambda \sim U(0, 1)
\end{equation}
While groundbreaking, standard SMOTE generates synthetic samples indiscriminately across feature space, often bridging the gap between legitimate and fraudulent clusters and introducing noise along complex decision boundaries.

\begin{itemize}[leftmargin=0.35in]
\item \textbf{Borderline-SMOTE and ADASYN}: Han et al. (2005) proposed Borderline-SMOTE, which restricts synthetic generation exclusively to minority instances located near the decision boundary (instances whose $k$-nearest neighbors are predominantly majority class). He et al. (2008) formulated ADASYN (Adaptive Synthetic Sampling), which uses a density distribution $\Gamma_i$ to dynamically generate more synthetic samples for minority instances that are harder to learn.
\item \textbf{OmniSMOTE}: In contemporary imbalanced learning, OmniSMOTE combines geometric boundary analysis with topological data analysis, ensuring that synthesized points respect multi-modal manifold geometry and do not violate convex hull bounds of legitimate transactions.
\end{itemize}

\begin{enumerate}[leftmargin=0.35in]
\item \textbf{Algorithm-Level Cost-Sensitive Learning}:
\end{enumerate}
\begin{itemize}[leftmargin=0.35in]
\item Elkan (2001) established the theoretical foundations of cost-sensitive learning, proving that altering the loss function's penalty weights is functionally equivalent to re-weighting class priors.
\item In modern gradient-boosted decision trees, Chen and Guestrin (2016) integrated the \texttt{scale\_pos\_weight} hyperparameter into XGBoost:
\end{itemize}
\begin{equation}
\text{scale\_pos\_weight} = \frac{\sum_{i=1}^N \mathbb{I}(y_i = 0)}{\sum_{i=1}^N \mathbb{I}(y_i = 1)}
\end{equation}
This dynamically scales the gradient and hessian updates for positive instances, compelling the tree split optimizer to prioritize minority class separation without altering input data distributions.


\subsection{Stacking Ensembles and Tree-Based Gradient Boosting}

Ensemble learning combines predictions from multiple heterogeneous or homogeneous base models to achieve superior predictive generalization and variance reduction compared to any individual constituent classifier.


\begin{itemize}[leftmargin=0.35in]
\item \textbf{Bagging (Bootstrap Aggregating)}: Breiman (1996) formulated Bagging, leading directly to the development of Random Forests (Breiman, 2001). By constructing hundreds of de-correlated decision trees trained on bootstrap samples with random feature sub-selection, Random Forests dramatically reduce model variance while resisting noise.
\item \textbf{Gradient Boosting}: Friedman (2001) introduced Gradient Boosting Machines (GBM), building additive trees sequentially by fitting each successive tree to the negative gradient (pseudo-residuals) of the loss function.
\item \textbf{Modern Gradient Boosted Trees}:
\item \textbf{XGBoost (Extreme Gradient Boosting)}: Developed by Chen and Guestrin (2016), XGBoost introduces second-order Taylor expansion approximations of the loss function, sparsity-aware split finding, cache-aware memory access, and built-in $L_1$ (Lasso) and $L_2$ (Ridge) regularization to prevent overfitting on structured tabular data.
\item \textbf{LightGBM}: Formulated by Ke et al. (2017), LightGBM introduces Gradient-based One-Side Sampling (GOSS) and Exclusive Feature Bundling (EFB). Unlike traditional level-wise tree growth, LightGBM utilizes leaf-wise (best-first) tree growth with depth constraints, achieving orders-of-magnitude faster training speeds and superior memory efficiency on massive transaction datasets.
\item \textbf{Stacking Generalization}: Wolpert (1992) formalized Stacking (Stacked Generalization). Unlike bagging (which averages predictions) or boosting (which trains sequential corrections), Stacking trains an ensemble of diverse base classifiers (Level-0), collects their out-of-fold probability predictions to form a meta-feature matrix, and trains a meta-learner (Level-1, typically Logistic Regression or an Elastic Net) to optimize the final decision boundary. In fraud detection, combining structurally diverse algorithms (gradient-boosted trees, bagged forests, and linear meta-learners) produces robust decision boundaries capable of capturing both local non-linear feature splits and global linear trends.
\end{itemize}

\subsection{Explainable Artificial Intelligence (XAI) in FinTech}

The deployment of complex non-linear ensemble models in banking operations created the urgent challenge of model opacity. When an algorithm blocks a cardholder's wire transfer or denies a merchant authorization, risk analysts must know the underlying rationale to prevent unfair discrimination, satisfy regulatory inquiries, and combat false positives.


\begin{enumerate}[leftmargin=0.35in]
\item \textbf{LIME (Local Interpretable Model-Agnostic Explanations)}:
\end{enumerate}
Ribeiro, Singh, and Guestrin (2016) introduced LIME. For a given complex model $f$ and an instance $\mathbf{x}$, LIME perturbs $\mathbf{x}$ in feature space, queries $f$ to obtain prediction labels for the perturbed points, weights the samples by their exponential proximity kernel $\pi_{\mathbf{x}}(\mathbf{z}) = \exp(-D(\mathbf{x}, \mathbf{z})^2 / \sigma^2)$, and trains an interpretable sparse linear surrogate model $g \in G$:

\begin{equation}
\xi(\mathbf{x}) = \arg\min_{g \in G} \mathcal{L}(f, g, \pi_{\mathbf{x}}) + \Omega(g)
\end{equation}
While highly intuitive, LIME suffers from sampling instability (explanations can vary slightly across repeated runs) and high runtime overhead due to runtime Monte Carlo perturbation sampling.


\begin{enumerate}[leftmargin=0.35in]
\item \textbf{SHAP (SHapley Additive exPlanations)}:
\end{enumerate}
Lundberg and Lee (2017) unified cooperative game theory and machine learning interpretability by introducing SHAP, rooted in Lloyd Shapley's (1953) axiomatic game theory. SHAP assigns each feature an attribution value representing its contribution to the displacement of the model prediction from the base expected value. SHAP is uniquely grounded in four desirable theoretical properties:

\begin{itemize}[leftmargin=0.35in]
\item \textbf{Local Accuracy (Efficiency)}: $\sum_{i=1}^m \phi_i(\mathbf{x}) = f(\mathbf{x}) - \phi_0$.
\item \textbf{Missingness}: A feature with no impact receives zero attribution.
\item \textbf{Consistency (Monotonicity)}: If a model changes such that a feature's marginal contribution increases or stays the same, its attribution value cannot decrease.
\item \textbf{Symmetry}: Identical features receive identical attributions.
\end{itemize}

\begin{enumerate}[leftmargin=0.35in]
\item \textbf{TreeSHAP}:
\end{enumerate}
In 2020, Lundberg et al. formulated TreeSHAP, an exact algorithm optimized specifically for tree ensembles. While general KernelSHAP scales exponentially with feature dimension ($\mathcal{O}(M 2^{|F|})$), TreeSHAP exploits the internal graph topology of decision trees, evaluating all feature subsets simultaneously in polynomial time:

\begin{equation}
\mathcal{O}(T L D^2)
\end{equation}
where $T$ is the number of trees, $L$ is the maximum number of leaves, and $D$ is the maximum tree depth. This algorithmic breakthrough enables real-time, sub-50ms exact feature attribution for production transaction scoring.


\subsection{Graph Analytics and Syndicate Ring Detection}

While tabular machine learning analyzes individual transactions in isolation, modern organized cybercrime relies on collaborative networks: mule accounts, smurfing rings, and synthetic identity clusters.


Akoglu, Chandy, and Faloutsos (2015) conducted an exhaustive survey on graph-based anomaly detection, demonstrating that financial transactions naturally map to directed, weighted multi-graphs:

\begin{equation}
\mathcal{G} = (\mathcal{V}, \mathcal{E}, \mathcal{W})
\end{equation}
where vertices $\mathcal{V}$ represent account entities (cardholders, merchants, bank accounts), directed edges $\mathcal{E} \subseteq \mathcal{V} \times \mathcal{V}$ represent fund transfers, and edge weights $\mathcal{W}$ denote transaction amounts and timestamps.


In financial money laundering, circular routing ($A \rightarrow B \rightarrow C \rightarrow A$) is a classic evasion topology used to obfuscate the illicit origin of stolen funds. Traditional relational databases cannot query multi-hop cyclic paths efficiently. By applying Tarjan's Strongly Connected Components algorithm and Johnson's elementary cycle-finding algorithm via in-memory graph engines (such as NetworkX), detection systems can uncover complex money laundering loops in real time.


\subsection{Federated Learning in Financial Systems}

Financial institutions are legally restricted by data privacy regulations (e.g., GDPR, PCI-DSS, Bank Secrecy Act) from sharing raw transaction records and cardholder PII with external competitors. However, fraudsters frequently exploit this information silo by launching coordinated, multi-bank card testing and account takeover campaigns across several institutions simultaneously.


McMahan et al. (2017) pioneered \textbf{Federated Learning (FL)} with the Federated Averaging (FedAvg) algorithm. In a federated ecosystem, $K$ distinct banking nodes train local models on their proprietary internal transaction data $\mathcal{D}_k$. Rather than transmitting raw transactions to a centralized repository, each bank computes local gradient parameter updates $\mathbf{w}_t^k$ and transmits only encrypted model weights to a central coordinator. The coordinator aggregates the weights:

\begin{equation}
\mathbf{w}_{t+1} = \sum_{k=1}^K \frac{n_k}{N} \mathbf{w}_t^k
\end{equation}
and broadcasts the updated global model back to participating institutions. Yang et al. (2019) formalized Federated Machine Learning concepts, showing that Horizontal Federated Learning enables banking consortiums to collaboratively train robust fraud classifiers that detect cross-bank attack campaigns while strictly preserving customer privacy and regulatory compliance.



\subsection{Comprehensive Literature Survey Matrix}

The following matrix synthesizes fifteen benchmark academic studies across fraud detection methodologies, highlighting their key contributions, datasets utilized, algorithms deployed, and identified engineering gaps:


\begin{table}[H]
\centering
\footnotesize
\begin{adjustbox}{width=\textwidth,center}
\begin{tabular}{| l | l | l | l | l | l | l | l |}
\hline
\rowcolor{tableheader} \textbf{Ref.} & \textbf{Authors \& Year} & \textbf{Publication Venue} & \textbf{Core Contribution} & \textbf{Datasets Used} & \textbf{Algorithms Evaluated} & \textbf{Performance Achieved} & \textbf{Critical Limitations \& Research Gaps} \\ \hline
\textbf{[1]} & Dal Pozzolo et al. (2015) & IEEE CIM & Investigated concept drift and class imbalance in credit card streams. & European Credit Card (284k txn) & Random Forest, SVM, Logistic Regression & PR-AUC: 0.74, ROC-AUC: 0.94 & No real-time explainability; single classifiers without stacking meta-learner. \\ \hline
\rowcolor{tablealt} \textbf{[2]} & Carcillo et al. (2018) & Information Sciences & Combined active learning and semi-supervised streaming fraud detection. & Real Banking Streams (45M txn) & Ensemble Random Forest, Active Sampling & Recall: 78.4\%, FPR: 1.2\% & Heavy computational latency; opaque black-box predictions; no regulatory reason codes. \\ \hline
\textbf{[3]} & Chawla et al. (2002) & JAIR & Introduced SMOTE synthetic oversampling for imbalanced data. & Synthetic \& Public Benchmarks & C4.5 Decision Trees, Naive Bayes & Significant ROC-AUC improvement & Overgeneralizes across boundary noise; lacks calibration for tree gradient boosters. \\ \hline
\rowcolor{tablealt} \textbf{[4]} & He et al. (2008) & IEEE IJCNN & Formulated ADASYN adaptive synthetic sampling. & UCI Benchmark Datasets & Decision Trees, SVM & G-Mean: 84.2\% & Can synthesize extreme outliers in high-dimensional PCA financial spaces. \\ \hline
\textbf{[5]} & Chen \& Guestrin (2016) & ACM KDD & Introduced XGBoost gradient boosting framework. & Higgs, Yahoo, Tabular Benchmarks & XGBoost, LightGBM, Random Forest & State-of-the-art tabular accuracy & Lacks integrated XAI explanations; raw scores suffer from extreme probability polarization. \\ \hline
\rowcolor{tablealt} \textbf{[6]} & Ke et al. (2017) & NeurIPS & Developed LightGBM with GOSS and EFB algorithms. & Flight, Yahoo, Microsoft LETOR & LightGBM, XGBoost & 20x speedup over standard GBM & Leaf-wise growth can overfit small minority fraud classes if unweighted. \\ \hline
\textbf{[7]} & Wolpert (1992) & Neural Networks & Formulated Stacking Generalization principles. & Theoretical \& Multi-Task & Heterogeneous Multi-Classifiers & Superior meta-generalization & High architectural complexity; requires specialized microservices for fast inference. \\ \hline
\rowcolor{tablealt} \textbf{[8]} & Lundberg \& Lee (2017) & NeurIPS & Formulated SHAP game-theoretic unified framework. & Tabular, Image, Text & KernelSHAP, DeepSHAP & Exact axiomatic attribution & KernelSHAP is computationally prohibitive for real-time sub-100ms API inference. \\ \hline
\textbf{[9]} & Lundberg et al. (2020) & Nature Machine Intelligence & Developed TreeSHAP for polynomial-time tree interpretability. & Medical \& Financial Tabular Data & TreeSHAP with XGBoost/LGBM & Exact local explanations in ms & Focused strictly on static offline analysis; no translation to FCRA reason codes. \\ \hline
\rowcolor{tablealt} \textbf{[10]} & Ribeiro et al. (2016) & ACM KDD & Developed LIME local surrogate explanations. & Image, Text, Tabular & Sparse Linear Surrogates & Interpretable local weights & High sampling instability; sensitive to perturbation kernel width $\sigma$. \\ \hline
\textbf{[11]} & Lopez-Rojas et al. (2016) & EMSS & Developed PaySim synthetic mobile money simulator. & PaySim Mobile Logs (6.3M txn) & Baseline Decision Trees & Accuracy: 98.2\% & Initial baseline models omitted balance delta equations and stacking ensembles. \\ \hline
\rowcolor{tablealt} \textbf{[12]} & Akoglu et al. (2015) & ACM TKDD & Comprehensive survey of graph-based anomaly detection. & Social \& Financial Graphs & PageRank, HITS, Graph Walk & High detection of network rings & Relational graph databases incur massive query latency under live streaming loads. \\ \hline
\textbf{[13]} & McMahan et al. (2017) & AISTATS & Formulated Federated Averaging (FedAvg) distributed learning. & MNIST, CIFAR, Language Models & Distributed SGD, FedAvg & Robust global convergence & Did not address extreme financial class imbalance across non-IID bank nodes. \\ \hline
\rowcolor{tablealt} \textbf{[14]} & Yang et al. (2019) & ACM TIST & Established Horizontal \& Vertical Federated Learning taxonomy. & FinTech \& Enterprise Systems & Secure Aggregation, Split Learning & High privacy preservation & Theoretical framework; lacks complete open-source microservice reference implementation. \\ \hline
\textbf{[15]} & European Commission (2024) & Official EU Journal & Promulgated the EU Artificial Intelligence Act. & Regulatory Compliance Document & Legal \& Compliance Standards & Mandated transparency & Establishes strict legal mandates but provides no engineering code implementation. \\ \hline
\end{tabular}
\end{adjustbox}
\end{table}


\section{Limitations of Existing System \& Gap Analysis}

\subsection{Critical Review of Contemporary Production Architectures}

Current industrial fraud detection solutions deployed across tier-1 banks and payment gateways are characterized by fractured, siloed architectures. Organizations typically run a patchwork of legacy rule engines (e.g., FICO Falcon, Actimize) coupled with disconnected, offline analytical models developed in isolation by data science teams.


An objective examination of these architectures exposes several fundamental vulnerabilities:


\begin{enumerate}[leftmargin=0.35in]
\item \textbf{The Dichotomy Between Rules and Machine Learning}:
\end{enumerate}
Traditional banks force a binary separation: transaction traffic passes first through a legacy rule engine; if a transaction passes without triggering a rule, it is either authorized immediately or sent to an asynchronous offline batch queue for machine learning scoring. This disjointed architecture fails because:

\begin{itemize}[leftmargin=0.35in]
\item Modern complex attacks easily bypass simplistic rule thresholds.
\item Batch ML scoring happens \textit{after} funds have already cleared, turning fraud prevention into a reactive loss-recovery effort.
\item Updating rule tables requires manual sign-offs and extensive regression testing, taking weeks to respond to emerging attack vectors.
\end{itemize}

\begin{enumerate}[leftmargin=0.35in]
\item \textbf{Extreme Score Polarization and Analyst Queue Starvation}:
\end{enumerate}
Standard gradient boosted models trained on imbalanced data yield raw probabilities clustered tightly at \\$0.000$ and \\$1.000$. Consequently, transactions are classified as either definitely safe or definitely fraudulent. In reality, modern attacks operate in grey zones: an authenticated cardholder making an unusually large purchase from a new device while traveling. By polarizing scores, standard models eliminate the intermediate risk tier, depriving banks of the opportunity to trigger step-up two-factor authentication (2FA) or route borderline transactions to specialized fraud analysts.


\begin{enumerate}[leftmargin=0.35in]
\item \textbf{Total Absence of Inline Regulatory Transparency}:
\end{enumerate}
When a black-box model blocks a customer's transaction, current production platforms return a generic error code (e.g., \texttt{DECLINE: ERROR\_CODE\_51}). Under GDPR Article 22, California Consumer Privacy Act (CCPA), and FCRA Section 615(a), consumers have the legal right to receive specific reasons explaining why an adverse decision was rendered. Banking support staff are currently unable to answer customer inquiries because the underlying AI provides no feature attribution.


\begin{enumerate}[leftmargin=0.35in]
\item \textbf{Vulnerability to Coordinated Syndicate Rings}:
\end{enumerate}
Tabular machine learning models treat each transaction as an independent, identically distributed ($i.i.d.$) point in vector space. They evaluate a transaction based solely on its immediate attributes (e.g., amount, time, user ID). They are completely blind to topological graph structures: such as an account receiving twenty \\$500 transfers from diverse sources and immediately aggregating them into a single \\$10,000 transfer to an offshore entity (classic money mule structuring).


\begin{enumerate}[leftmargin=0.35in]
\item \textbf{Inability to Withstand Covariate Concept Drift}:
\end{enumerate}
Financial data distributions are inherently non-stationary. Consumer purchasing patterns change dramatically during holiday shopping seasons, macroeconomic inflation, or promotional events. Simultaneously, cybercriminals constantly alter their evasion tactics once an existing vector is closed. Standard production systems have no automated mechanism to measure feature distribution divergence, leading to silent model performance degradation.



\subsection{Comparative Analysis: Legacy vs. Current ML vs. FraudShield AI}

The following table provides a comprehensive multi-dimensional comparison contrasting traditional legacy rule engines, conventional machine learning models, and the proposed FraudShield AI architecture:


\begin{table}[H]
\centering
\footnotesize
\begin{adjustbox}{width=\textwidth,center}
\begin{tabular}{| l | l | l | l |}
\hline
\rowcolor{tableheader} \textbf{Architectural Dimension} & \textbf{Legacy Rule-Based Systems (e.g., Falcon v1)} & \textbf{Conventional Black-Box ML (e.g., Standalone XGBoost)} & \textbf{Proposed FraudShield AI Platform} \\ \hline
\textbf{Detection Methodology} & Hardcoded Boolean threshold conditions. & Single-classifier statistical probability learning. & Multi-Domain Hybrid Stacking Ensemble (XGB+LGBM+RF+Meta). \\ \hline
\rowcolor{tablealt} \textbf{Handling of Imbalance} & Manual threshold tuning; unweighted. & Basic undersampling or unweighted cost function. & Dynamic \texttt{scale\_pos\_weight} + OmniSMOTE oversampling. \\ \hline
\textbf{False Positive Rate (FPR)} & Extremely high (\\$85\% - 95\%$). & Moderate (\\$15\% - 25\%$). & Ultra-low ($< 5.7\%$ on complex PCA vectors; \\$0\%$ on mobile money). \\ \hline
\rowcolor{tablealt} \textbf{Decision Spectrum} & Binary (\texttt{PASS} or \texttt{FAIL}). & Highly polarized raw probabilities (\\$0.0$ or \\$1.0$). & Continuous Calibrated Risk Score (\\$0\%-100\%$) across 3 Tiers. \\ \hline
\textbf{Triage \& Human-in-the-Loop} & Rigid alert queues without context. & Queue starved due to probability polarization. & Nuanced \texttt{NEEDS REVIEW} Tier (\\$45\%-75\%$) for targeted 2FA/Analyst review. \\ \hline
\rowcolor{tablealt} \textbf{Model Interpretability} & Transparent but simplistic (shows rule). & Completely opaque "black box". & Real-Time TreeSHAP feature attributions + Local Waterfall Plots. \\ \hline
\textbf{Regulatory Compliance} & Partial (auditable rules, but brittle). & Non-Compliant (Violates GDPR Art 22 \& FCRA). & 100\% Compliant (Automated FCRA Adverse Action Reason Codes). \\ \hline
\rowcolor{tablealt} \textbf{Graph Syndicate Detection} & None (evaluates individual records). & None (treats data as isolated vectors). & In-Memory NetworkX Cycle Detection \& Mule Ring Clustering. \\ \hline
\textbf{Drift \& Degradation Audit} & Manual rule inspection after loss spikes. & None (requires periodic manual retraining). & Automated 2-Sample KS-Tests \& Population Stability Index (PSI). \\ \hline
\rowcolor{tablealt} \textbf{Inter-Bank Collaboration} & Centralized fraud lists (blacklist sharing). & None (data privacy prevents model sharing). & Decentralized Federated Learning (FedAvg) Collaborative Defense. \\ \hline
\textbf{Inference Latency SLA} & Sub-10ms (simplistic lookups). & 50ms - 200ms (unoptimized Python). & $< 100\text{ms}$ (Asynchronous FastAPI microservices with caching). \\ \hline
\rowcolor{tablealt} \textbf{User \& Analyst Interface} & Clunky legacy green-screen terminal. & Static Jupyter Notebooks or basic tables. & Interactive Modern SPA (React 19 + Vite 6 + Tailwind CSS). \\ \hline
\end{tabular}
\end{adjustbox}
\end{table}


\subsection{Root Cause Analysis: 5-Why and Ishikawa (Fishbone) Analysis}

To trace operational failures in modern fraud management down to their fundamental architectural root causes, two systematic engineering methodologies were conducted: \textbf{5-Why Analysis} and \textbf{Ishikawa (Fishbone) Diagram Formulation}.


\subsubsection{5-Why Root Cause Analysis}

\begin{itemize}[leftmargin=0.35in]
\item \textbf{Why 1: Why do financial institutions continue to experience tens of billions of dollars in annual unrecovered fraud losses?}
\end{itemize}
\textit{Because existing production platforms fail to block sophisticated, non-linear, multi-channel fraud patterns in real time.}

\begin{itemize}[leftmargin=0.35in]
\item \textbf{Why 2: Why do production platforms fail to detect these dynamic fraud patterns?}
\end{itemize}
\textit{Because they rely heavily on static rule-based thresholds and standalone single-model classifiers that cannot adapt to non-stationary fraud strategies.}

\begin{itemize}[leftmargin=0.35in]
\item \textbf{Why 3: Why do standard machine learning classifiers fail to generalize effectively in live banking streams?}
\end{itemize}
\textit{Because extreme class imbalance causes models to heavily prioritize the majority non-fraud class, resulting in missed fraud (low recall) and high false positive rates.}

\begin{itemize}[leftmargin=0.35in]
\item \textbf{Why 4: Why are banking fraud analysts unable to promptly verify, explain, or override automated fraud classifications?}
\end{itemize}
\textit{Because advanced gradient boosted models and deep neural networks operate as opaque black boxes that emit uncalibrated probabilities without human-understandable feature attributions.}

\begin{itemize}[leftmargin=0.35in]
\item \textbf{Why 5: Why has the financial sector not unified calibrated stacking ensembles, explainable AI, graph intelligence, and real-time streaming interfaces?}
\end{itemize}
\textit{Because existing software solutions are developed in isolated silos---data scientists train models in offline notebooks, network engineers maintain separate relational databases, and compliance officers manage manual regulatory paperwork---lacking a unified, full-stack microservice reference architecture.}


\subsubsection{Ishikawa (Fishbone) Root Cause Diagram}

\begin{lstlisting}
                 ISHIKAWA (FISHBONE) DIAGRAM: CAUSES OF INEFFECTIVE FRAUD DEFENSE

    METHOD (Detection Pipeline)           MATERIAL (Transaction Data)           MACHINE (Compute & Infra)
    ---------------------------           ---------------------------           -------------------------
    - Static hardcoded rules              - Extreme class imbalance             - High inference latency
    - Single classifier reliance          - High-dimensional anonymized PCA      - Legacy batch architectures
    - Lack of ensemble diversity          - Non-stationary concept drift        - Lack of async streaming APIs
    - Uncalibrated decision limits        - Multi-schema domain heterogeneity   - Heavy unoptimized XAI models
                                   \                   |                   /
                                    \                  |                  /
                                     \                 |                 /
                                      ===================================>  INEFFICIENT REAL-TIME
                                     /                 |                 \  FRAUD DETECTION, HIGH
                                    /                  |                  \ FPR & REGULATORY OPACITY
                                   /                   |                   \
    ---------------------------           ---------------------------           -------------------------
    - Severe analyst alert fatigue        - Exponential surge in e-commerce     - Evaluation on accuracy alone
    - Lack of XAI decision clarity        - Cross-border cyber syndicates       - Neglect of PR-AUC & F1-Score
    - Manual review backlogs              - Sophisticated money mule rings      - Absence of real-time KS/PSI
    - Inability to meet FCRA rules        - Coordinated card-testing attacks     drift metric tracking
       MAN (Human Factors)                   ENVIRONMENT (Threat Landscape)        MEASUREMENT (Metrics & KPI)
\end{lstlisting}

As illustrated in Figure 2.1, the breakdown of financial fraud detection stems from intersecting factors:

\begin{enumerate}[leftmargin=0.35in]
\item \textbf{Method}: Outdated heuristics, absence of stacking ensemble diversity, and uncalibrated decision thresholds.
\item \textbf{Material}: Extreme class imbalance (0.17\% fraud prevalence), high dimensionality ($V_1-V_{28}$ PCA components), and rapid behavioral distribution shifts.
\item \textbf{Machine}: Legacy batch-processing architectures that introduce unacceptable latency and unoptimized XAI pipelines.
\item \textbf{Man}: Fraud analyst fatigue caused by thousands of false alerts and opaque probability outputs that cannot be explained to cardholders.
\item \textbf{Environment}: Rapid escalation of global card-not-present fraud, automated botnets, and organized money laundering syndicates.
\item \textbf{Measurement}: Misleading evaluation of machine learning models on overall accuracy rather than cost-sensitive F1-score, Precision-Recall AUC, and calibrated risk tiers.
\end{enumerate}


\subsection{Research Gap Analysis and Proposed Countermeasures}

Based on the systematic literature survey and root cause diagnostics, four primary research and engineering gaps were identified. Table 2.3 details these gaps alongside the specific engineering countermeasures implemented in FraudShield AI:


\begin{table}[H]
\centering
\footnotesize
\begin{adjustbox}{width=\textwidth,center}
\begin{tabular}{| l | l | l | l |}
\hline
\rowcolor{tableheader} \textbf{Research Gap Identifier} & \textbf{Identified Academic \& Industrial Gap} & \textbf{Current State of the Art} & \textbf{FraudShield AI Engineering Countermeasure} \\ \hline
\textbf{GAP 1: Disjointed Imbalance Resampling and Stacking Ensembles} & Most studies evaluate basic SMOTE or random undersampling on single models (Logistic Regression or Random Forest). & Single models fail to capture multi-scale non-linear interactions; basic SMOTE generates boundary distortion in PCA spaces. & Coupled \textbf{OmniSMOTE} boundary-aware oversampling with dynamic algorithmic class-weighting (\texttt{scale\_pos\_weight}) and a multi-model \textbf{Stacking Ensemble} (XGBoost + LightGBM + Random Forest + Logistic Regression meta-learner). \\ \hline
\rowcolor{tablealt} \textbf{GAP 2: Regulatory Black-Box Opacity in High-Performing Tree Ensembles} & Gradient boosted ensembles achieve superior predictive accuracy but provide no human-interpretable rationale for transaction flags. & Models output raw probabilities (\\$0.0$ to \\$1.0$); financial institutions violate GDPR Article 22 and FCRA disclosure mandates. & Embedded an optimized \textbf{TreeSHAP} engine that evaluates polynomial-time feature attributions ($\phi_i$) in under 50ms, dynamically translating Shapley values into standardized \textbf{FCRA Adverse Action Reason Codes}. \\ \hline
\textbf{GAP 3: Score Polarization and Absence of Continuous Risk Calibration} & Tree ensemble loss functions produce polarized probability outputs ($p \in \{0.0001, 0.9999\}$), leaving the middle-tier review queue empty. & Triage is forced into a binary pass/fail; cardholders suffer abrupt transaction declines without step-up authentication. & Engineered a \textbf{Calibrated Continuous Risk Engine} that synthesizes raw probabilities, transaction velocity, balance liquidation heuristics, and geolocation deltas into a smooth 3-tier risk hierarchy (\texttt{SAFE}, \texttt{NEEDS REVIEW}, \texttt{FRAUD}). \\ \hline
\rowcolor{tablealt} \textbf{GAP 4: Siloed Tabular Models Blind to Coordinated Mule Rings and Drift} & Tabular classifiers evaluate transactions in complete isolation, failing to detect multi-account circular money laundering or feature distribution shift. & Network analysis and concept drift monitoring exist only as disconnected, offline academic prototypes. & Integrated an in-memory \textbf{NetworkX Graph Intelligence Engine} for real-time cyclic path and community detection, paired with an \textbf{MLOps Concept Drift Engine} tracking Kolmogorov-Smirnov and PSI metrics. \\ \hline
\end{tabular}
\end{adjustbox}
\end{table}


\section{Proposed System}

\subsection{Architectural Paradigm and Design Philosophy}

\textbf{FraudShield AI} addresses the compounding failures of traditional fraud detection by introducing a unified, modular, enterprise-grade architecture. Rather than treating machine learning, explainability, graph intelligence, and user presentation as disconnected silos, FraudShield AI unifies these capabilities into an asynchronous microservice ecosystem designed around five core engineering principles:


\begin{enumerate}[leftmargin=0.35in]
\item \textbf{Domain-Specialized Modularity}:
\end{enumerate}
Acknowledging that no single classifier can effectively process heterogeneous financial schemas, FraudShield AI maintains a dynamic \textbf{Domain Model Registry}. Incoming payloads are automatically classified by schema and routed to specialized model heads (PaySim, Credit Card PCA, Spatial Behavioral, or BankSim).

\begin{enumerate}[leftmargin=0.35in]
\item \textbf{Heterogeneous Stacking Diversity}:
\end{enumerate}
Each domain head deploys a two-tier Stacking Ensemble combining structurally diverse base learners: XGBoost (depth-optimized gradient boosting), LightGBM (leaf-wise gradient boosting), and Random Forest (bagged variance reduction), synthesized by a cross-validated Logistic Regression meta-learner.

\begin{enumerate}[leftmargin=0.35in]
\item \textbf{Equitable Continuous Risk Scoring}:
\end{enumerate}
The platform eliminates binary probability polarization through a calibrated scoring function that produces a smooth, continuous risk metric (\\$0.0\%$ to \\$100.0\%$), establishing an actionable \texttt{NEEDS REVIEW} tier (\\$45\%-75\%$) for step-up multi-factor authentication or manual analyst triage.

\begin{enumerate}[leftmargin=0.35in]
\item \textbf{Sub-100ms Explainability as a First-Class Citizen}:
\end{enumerate}
Model interpretability is not an afterthought or offline batch script. Exact TreeSHAP feature attributions and FCRA Adverse Action reason codes are generated synchronously within the sub-100ms API inference lifecycle.

\begin{enumerate}[leftmargin=0.35in]
\item \textbf{Decentralized and Graph-Aware Defense}:
\end{enumerate}
The system uncovers organized cybercrime syndicates via real-time cyclic graph traversal and enables cross-institutional consortium model training via Federated Learning (FedAvg) without exposing customer PII.



\subsection{High-Level System Architecture}

The following ASCII diagram illustrates the end-to-end architectural topology of FraudShield AI, detailing data flows from external ingestion through security gateways, analytical engines, and presentation layers:


\begin{lstlisting}
+---------------------------------------------------------------------------------------------+
|                                 PRESENTATION LAYER                                          |
|                          React 19 + Vite 6 + Tailwind CSS SPA                               |
|  +-----------------------+  +------------------------+  +--------------------------------+  |
|  | Threat Executive View |  | Live Monitoring Stream |  | Transaction Analysis Inspector |  |
|  +-----------------------+  +------------------------+  +--------------------------------+  |
|  +-----------------------+  +------------------------+  +--------------------------------+  |
|  | TreeSHAP XAI & FCRA   |  | Graph Intelligence UI  |  | MLOps Drift & Federated Portal |  |
|  +-----------------------+  +------------------------+  +--------------------------------+  |
+----------------------------------------------+----------------------------------------------+
                                               | HTTPS / JSON REST API (Port 8008)
                                               v
+---------------------------------------------------------------------------------------------+
|                             ENTERPRISE API GATEWAY (FastAPI)                                |
|  +------------------------+  +------------------------+  +--------------------------------+  |
|  | OWASP Security Headers |  | Sliding-Window Rate-Lim|  | JWT Role-Based Access (RBAC)   |  |
|  +------------------------+  +------------------------+  +--------------------------------+  |
+----------------------------------------------+----------------------------------------------+
                                               | Validated Ingestion Payload
                                               v
+---------------------------------------------------------------------------------------------+
|                          SCHEMA DETECTION & DOMAIN REGISTRY ROUTER                          |
|                                (backend/app/core/model_engine.py)                           |
|     +-------------------+-------------------+-------------------+---------------------+     |
|     v                   v                   v                   v                     v     |
|  +---------------+   +---------------+   +---------------+   +---------------+   +-----------+  |
|  | PaySim Head   |   | CC PCA Head   |   | Spatial Head  |   | BankSim Head  |   | Fallback  |  |
|  | Mobile Money  |   | 28 PCA Vectors|   | Haversine Geo |   | Retail Trans  |   | Heuristic |  |
|  +-------+-------+   +-------+-------+   +-------+-------+   +-------+-------+   +-----+-----+  |
+----------+-------------------+-------------------+-------------------+-----------------+--------+
           +-------------------+---------+---------+-------------------+                 |
                                         v                                               |
+----------------------------------------------------------------------------------------+----+
|                              ANALYTICS & INFERENCE PIPELINE                                 |
|                                                                                             |
|  +-------------------------------------+         +---------------------------------------+  |
|  |       STACKING ENSEMBLE ENGINE      |         |         EXPLAINABLE AI ENGINE         |  |
|  |  Level-0: XGBoost + LightGBM + RF   | <-----> |  Exact TreeSHAP Attribution Engine    |  |
|  |  Level-1: Logistic Regression Meta  |         |  FCRA Adverse Action Code Generator   |  |
|  +------------------+------------------+         +---------------------------------------+  |
|                     | Raw Probability (p)                                                   |
|                     v                                                                       |
|  +---------------------------------------------------------------------------------------+  |
|  |                         CALIBRATED CONTINUOUS RISK ENGINE                             |  |
|  |      Composite Risk Score: SAFE (<45%) | NEEDS REVIEW (45-75%) | FRAUD (>=75%)       |  |
|  +------------------+--------------------------------------------------------------------+  |
|                     | Calibrated Score + Decision Metadata                                  |
|                     v                                                                       |
|  +-------------------------------------+         +---------------------------------------+  |
|  |      GRAPH INTELLIGENCE ENGINE      |         |          MLOPS DRIFT MONITOR          |  |
|  |  NetworkX In-Memory Transaction Net |         |  Two-Sample Kolmogorov-Smirnov Tests  |  |
|  |  Cycles, Communities, Mule Scores   |         |  Population Stability Index (PSI)     |  |
|  +-------------------------------------+         +---------------------------------------+  |
+----------------------------------------------+----------------------------------------------+
                                               |
                                               v
+---------------------------------------------------------------------------------------------+
|                              STORAGE & AUDITING INFRASTRUCTURE                              |
|  +------------------------+  +------------------------+  +--------------------------------+  |
|  | Append-Only Audit Log  |  | Serialized .pkl Models |  | Background Reference Samples   |  |
|  +------------------------+  +------------------------+  +--------------------------------+  |
+---------------------------------------------------------------------------------------------+
\end{lstlisting}


\subsection{Core Subsystems and Modular Breakdown}

The FraudShield AI platform is partitioned into seven distinct, highly cohesive functional subsystems:


\subsubsection{1. Presentation and Visualization Layer (\texttt{frontend/src})}
Constructed using \textbf{React 19}, \textbf{Vite 6}, and \textbf{Tailwind CSS}, the client application delivers an enterprise-grade threat operations interface. Key architectural modules include:

\begin{itemize}[leftmargin=0.35in]
\item \textbf{Executive Threat Dashboard (\texttt{DashboardModule.jsx})}: Displays high-level KPIs including total processed volume, active fraud detection rate, prevented financial loss, average inference latency, and interactive volume trend charts.
\item \textbf{Live Stream Monitor (\texttt{LiveMonitoringModule.jsx})}: Simulates a high-throughput banking ingestion pipe, rendering transactions as color-coded event cards with animated risk badges (\texttt{SAFE}, \texttt{NEEDS REVIEW}, \texttt{FRAUD}) and quick-triage inspection drawers.
\item \textbf{Single Transaction Inspector (\texttt{TransactionAnalysisModule.jsx})}: A deep diagnostic workbench allowing risk analysts to modify individual feature values ($V_1-V_{28}$, amounts, geographical coordinates, account balances) and trigger synchronous scoring.
\item \textbf{Explainable AI Visualizer (\texttt{ExplainableAIModule.jsx})}: Renders interactive TreeSHAP waterfall charts, showing baseline displacement, positive risk drivers (red bars pushing toward fraud), negative mitigating drivers (green bars pulling toward safe), and standardized FCRA Adverse Action reason codes.
\item \textbf{Graph Intelligence Viewer (\texttt{GraphIntelligenceModule.jsx})}: Visualizes multi-hop account transaction networks, rendering cyclic money laundering paths and highlighting mule accounts based on graph centrality.
\item \textbf{Concept Drift \& MLOps Monitor (\texttt{ConceptDriftModule.jsx})}: Displays real-time feature distribution divergence, KS-test $p$-values, and PSI indicators, featuring an interactive sandbox to inject synthetic adversarial drift spikes.
\item \textbf{Federated Learning Defense Portal (\texttt{FederatedLearningModule.jsx})}: Simulates a multi-bank consortium (Bank Alpha, Beta, Gamma), illustrating local model training rounds and global FedAvg parameter convergence without PII sharing.
\item \textbf{Identity \& RBAC Portal (\texttt{AuthModal.jsx})}: Implements cryptographic JWT session handling and 1-click persona switching across SOC Analyst, Compliance Officer, and System Administrator roles.
\end{itemize}

\subsubsection{2. Enterprise API Gateway (\texttt{backend/main.py})}
Built on \textbf{FastAPI} and served via \textbf{Uvicorn}, the gateway acts as the secure entry point:

\begin{itemize}[leftmargin=0.35in]
\item \textbf{OWASP Security Headers}: Enforces strict \texttt{Content-Security-Policy}, \texttt{X-Content-Type-Options: nosniff}, \texttt{X-Frame-Options: DENY}, and \texttt{Referrer-Policy: strict-origin-when-cross-origin}.
\item \textbf{Sliding-Window Rate Limiter}: Protects sensitive inference endpoints against automated scraping and distributed denial-of-service (DDoS) attacks.
\item \textbf{CORS Middleware}: Manages cross-origin resource sharing for secure edge CDN deployments.
\item \textbf{Append-Only Audit Trail}: Cryptographically logs every authentication event, transaction decision, and manual analyst override with microsecond timestamps.
\end{itemize}

\subsubsection{3. Domain Model Registry Router (\texttt{backend/app/core/model\_engine.py})}
The registry maintains pre-trained Stacking Ensemble model heads. When an incoming JSON payload arrives at \texttt{/analyze}, the schema detector inspects the feature keys:

\begin{itemize}[leftmargin=0.35in]
\item If keys match $\{V_1, V_2, \dots, V_{28}, \text{Amount}\}$, the transaction routes to the \textbf{Credit Card PCA Head}.
\item If keys match $\{\text{oldbalanceOrg}, \text{newbalanceOrig}, \text{type}\}$, it routes to the \textbf{PaySim Mobile Money Head}.
\item If keys match $\{\text{lat}, \text{long}, \text{merch\\_lat}, \text{merch\\_long}\}$, it routes to the \textbf{Spatial Behavioral Head}.
\item If keys match $\{\text{category}, \text{merchant}, \text{amount}\}$, it routes to the \textbf{BankSim Retail Head}.
\item If no schema matches, the transaction routes to a graceful fallback heuristic engine.
\end{itemize}

\subsubsection{4. Stacking Ensemble Engine (\texttt{model\_engine.py})}
Each domain head combines three high-capacity base estimators:

\begin{itemize}[leftmargin=0.35in]
\item \textbf{XGBoost Classifier}: Optimized for non-linear interactions, regularization, and gradient split handling.
\item \textbf{LightGBM Classifier}: Optimized for leaf-wise tree growth, high-speed split calculation, and categorical encoding.
\item \textbf{Random Forest Classifier}: Bagged ensemble of 100 de-correlated trees providing variance reduction and resistance to noise.
\item \textbf{Meta-Learner}: A calibrated Logistic Regression model ($C=1.0$) that takes the predicted probability vectors from Level-0 estimators and computes the optimal meta-decision boundary.
\end{itemize}

\subsubsection{5. Calibrated Continuous Risk Engine (\texttt{backend/app/core/risk\_engine.py})}
The risk engine eliminates score polarization by applying continuous mathematical transformations incorporating transactional velocity, account liquidation heuristics, and geolocation discrepancies. The resulting risk score is mapped into three operational tiers:

\begin{itemize}[leftmargin=0.35in]
\item [SAFE] \textbf{\texttt{SAFE}} ($< 45.0\%$): Immediate automated approval.
\item [REVIEW] \textbf{\texttt{NEEDS REVIEW}} (\\$45.0\% - 74.9\%$): Routed to the fraud analyst queue or triggered for step-up SMS/Biometric 2FA.
\item [FRAUD] \textbf{\texttt{FRAUD}} ($\ge 75.0\%$): Automated transaction decline and temporary account freeze.
\end{itemize}

\subsubsection{6. Explainable AI \& FCRA Reason Code Engine (\texttt{backend/app/core/explainer.py})}
The XAI engine utilizes \textbf{TreeSHAP} to compute exact Shapley values ($\phi_i$) for tree ensembles. The engine extracts the top contributing features and converts them into standardized Fair Credit Reporting Act (FCRA) Adverse Action Reason Codes:

\begin{itemize}[leftmargin=0.35in]
\item \texttt{RC-BAL-01}: Unbalanced Transaction Settlement
\item \texttt{RC-BAL-02}: Complete Account Liquidation Anomaly
\item \texttt{RC-AMT-01}: High-Value Transaction Velocity Outlier
\item \texttt{RC-GEO-01}: Haversine Geolocation Distance Discrepancy
\item \texttt{RC-PCA-01}: High-Dimensional Latent Anomaly Vector ($V_i$)
\end{itemize}

\subsubsection{7. Graph Intelligence \& MLOps Drift Subsystems (\texttt{graph\_engine.py} \& \texttt{drift\_monitor.py})}
\begin{itemize}[leftmargin=0.35in]
\item \textbf{Graph Engine}: Builds an in-memory transactional multi-graph using \textbf{NetworkX}, detecting directed cycles ($A \rightarrow B \rightarrow C \rightarrow A$) and calculating composite mule scores based on degree centrality and velocity.
\item \textbf{Drift Monitor}: Compares streaming transaction windows against baseline reference distributions using two-sample Kolmogorov-Smirnov (KS) tests and Population Stability Index (PSI):
\end{itemize}
\begin{equation}
\text{PSI} = \sum_{j=1}^B \left( P_j - Q_j \right) \ln\left(\frac{P_j}{Q_j}\right)
\end{equation}
triggering automated alerts when $\text{PSI} > 0.25$ indicates severe covariate shift.



\textit{End of Chapter 2 --- Literature Survey \& Proposed System}


\newpage


\chapter{Requirement Gathering, Analysis and Planning}

\section{Requirement Specification}

Software Requirement Specification (SRS) establishes the foundational functional capabilities, performance standards, architectural constraints, and operational boundaries of the FraudShield AI platform. The system is engineered to function within mission-critical banking environments where high availability, ultra-low latency, mathematical transparency, and regulatory compliance are non-negotiable.


\subsection{Functional Requirements Specifications (FR)}

The platform's functional scope is partitioned into twelve formal Functional Requirements (FR-01 through FR-12) detailed below:


\subsubsection{Table 3.1: Functional Requirements Specifications}

\begin{table}[H]
\centering
\footnotesize
\begin{adjustbox}{width=\textwidth,center}
\begin{tabular}{| l | l | l | l |}
\hline
\rowcolor{tableheader} \textbf{Requirement ID} & \textbf{Requirement Title} & \textbf{Requirement Description \& Acceptance Criteria} & \textbf{Priority} \\ \hline
\textbf{FR-01} & Multi-Domain Schema Ingestion \& Routing & The system must automatically detect the schema of incoming transaction payloads (e.g., PaySim, Credit Card PCA, Spatial Behavioral, BankSim) and route the payload to the corresponding pre-trained domain model head. Acceptance: Routing completes in $< 2\text{ms}$ with 100\% schema match. & High \\ \hline
\rowcolor{tablealt} \textbf{FR-02} & Stacking Ensemble Inference & The platform must execute Level-0 inference across XGBoost, LightGBM, and Random Forest base classifiers and synthesize predictions using a Logistic Regression meta-learner. Acceptance: Ensemble output yields continuous posterior probability $p \in [0, 1]$. & High \\ \hline
\textbf{FR-03} & Calibrated Continuous Risk Scoring & The system must map raw ensemble probabilities into a continuous risk score (\\$0.0\% - 100.0\%$) using transaction velocity, account liquidation heuristics, and geolocation deltas, assigning one of three decision tiers: \texttt{SAFE}, \texttt{NEEDS REVIEW}, or \texttt{FRAUD}. & High \\ \hline
\rowcolor{tablealt} \textbf{FR-04} & Real-Time TreeSHAP Feature Attribution & The platform must synchronously compute exact Shapley values (TreeSHAP) for each input feature during transaction scoring. Acceptance: Feature attributions must satisfy additive efficiency ($\sum \phi_i = f(x) - \phi_0$) within $< 50\text{ms}$. & High \\ \hline
\textbf{FR-05} & Automated FCRA Reason Code Generation & For any transaction assigned to \texttt{NEEDS REVIEW} or \texttt{FRAUD}, the system must extract top positive risk-driving features and map them into standardized Fair Credit Reporting Act (FCRA) Adverse Action Reason Codes (e.g., \texttt{RC-BAL-02}, \texttt{RC-GEO-01}). & High \\ \hline
\rowcolor{tablealt} \textbf{FR-06} & In-Memory Cyclic Mule Ring Detection & The graph intelligence engine must construct an in-memory directed transaction graph and detect multi-hop cycles ($A \rightarrow B \rightarrow C \rightarrow A$) of length $k \in [2, 6]$. Acceptance: Detects cyclic money laundering paths in $< 10\text{ms}$ for up to 10,000 active nodes. & High \\ \hline
\textbf{FR-07} & Account Mule Centrality Scoring & The system must compute composite mule risk scores for individual accounts based on graph in-degree, out-degree, transaction velocity, and balance depletion flags. Acceptance: Scores updated dynamically upon new edge creation. & Medium \\ \hline
\rowcolor{tablealt} \textbf{FR-08} & MLOps Concept Drift Monitoring & The platform must monitor streaming feature distributions against baseline reference distributions using two-sample Kolmogorov-Smirnov (KS) tests and Population Stability Index (PSI). Acceptance: Generates automated warnings when $\text{PSI} > 0.10$ and alerts when $\text{PSI} > 0.25$. & Medium \\ \hline
\textbf{FR-09} & Synthetic Adversarial Drift Injection & The platform must provide an administrative sandbox allowing authorized users to inject synthetic covariate drift vectors (e.g., balance anomalies, amount spikes) to test system resilience. & Medium \\ \hline
\rowcolor{tablealt} \textbf{FR-10} & Federated Learning Simulation (FedAvg) & The system must simulate decentralized multi-bank collaborative training across at least three independent nodes (Bank Alpha, Beta, Gamma) using Federated Averaging (FedAvg), updating global model weights without sharing raw transaction records. & Medium \\ \hline
\textbf{FR-11} & Role-Based Access Control (RBAC) & The platform must enforce strict cryptographic JWT authentication, provisioning distinct permissions for three user personas: SOC Analyst, Compliance Officer, and System Administrator. & High \\ \hline
\rowcolor{tablealt} \textbf{FR-12} & Immutable Audit Trail Logging & The system must maintain an append-only, tamper-evident audit log recording user logins, single-transaction inspections, manual triage overrides, and federated training rounds with microsecond timestamps. & High \\ \hline
\end{tabular}
\end{adjustbox}
\end{table}


\subsection{Non-Functional Requirements Specifications (NFR)}

Non-functional requirements specify operational criteria used to evaluate system performance, security posture, reliability, and usability:


\subsubsection{Table 3.2: Non-Functional Requirements Specifications}

\begin{table}[H]
\centering
\footnotesize
\begin{adjustbox}{width=\textwidth,center}
\begin{tabular}{| l | l | l |}
\hline
\rowcolor{tableheader} \textbf{Requirement ID} & \textbf{Quality Attribute} & \textbf{Technical Specification \& Target Benchmark} \\ \hline
\textbf{NFR-01} & Performance \& Latency & Single-transaction inference and risk scoring must execute within an average response latency of $< 100\text{ms}$ (P95 $< 150\text{ms}$, P99 $< 250\text{ms}$) under a baseline concurrency of 50 requests/sec. \\ \hline
\rowcolor{tablealt} \textbf{NFR-02} & High Throughput & The asynchronous FastAPI microservice backend must support sustained throughput of $\ge 500\text{ transactions/sec}$ in batch ingestion mode on standard quad-core server hardware. \\ \hline
\textbf{NFR-03} & System Availability & The platform architecture must support \\$99.9\%$ service availability, utilizing container health probes (\texttt{/health}) and automatic process restarts via Docker orchestration. \\ \hline
\rowcolor{tablealt} \textbf{NFR-04} & Security \& Hardening & All API endpoints must enforce OWASP Recommended Security Headers: \texttt{X-Content-Type-Options: nosniff}, \texttt{X-Frame-Options: DENY}, and strict CORS policies. Sensitive tokens must utilize HMAC-SHA256 signatures with 60-minute expiration windows. \\ \hline
\textbf{NFR-05} & Anti-DDoS Protection & The API Gateway must implement sliding-window rate limiting (e.g., max 100 requests per minute per IP for public endpoints) returning HTTP 429 Too Many Requests upon breach. \\ \hline
\rowcolor{tablealt} \textbf{NFR-06} & Explainability Latency & Synchronous TreeSHAP feature attribution calculation must not introduce more than \\$50\text{ms}$ of computational overhead to the base inference pipeline. \\ \hline
\textbf{NFR-07} & Responsive Web Usability & The React 19 frontend SPA must achieve a Google Lighthouse performance score of $\ge 90$, render fluid micro-animations at 60 FPS, and adapt responsively across desktop, tablet, and mobile displays. \\ \hline
\rowcolor{tablealt} \textbf{NFR-08} & Maintainability \& Testability & Codebase must maintain modular separation of concerns (Core, API, Schemas, UI), achieve $\ge 90\%$ code coverage across core scoring logic, and execute automated pytest suites in $< 15\text{ seconds}$. \\ \hline
\end{tabular}
\end{adjustbox}
\end{table}


\subsection{Role-Based Access Control (RBAC) Specifications}

To safeguard sensitive financial data and prevent unauthorized override of automated fraud decisions, FraudShield AI enforces a granular Role-Based Access Control matrix. The platform defines three distinct user personas:


\subsubsection{Table 3.3: Role-Based Access Control (RBAC) Permission Matrix}

\begin{table}[H]
\centering
\footnotesize
\begin{adjustbox}{width=\textwidth,center}
\begin{tabular}{| l | l | l | l | l |}
\hline
\rowcolor{tableheader} \textbf{Functional Capability / Endpoint} & \textbf{Public / Unauthenticated} & \textbf{SOC Analyst (\texttt{soc\_analyst})} & \textbf{Compliance Officer (\texttt{compliance\_officer})} & \textbf{System Administrator (\texttt{admin})} \\ \hline
System Health Check (\texttt{/health}) & [YES] Allowed & [YES] Allowed & [YES] Allowed & [YES] Allowed \\ \hline
\rowcolor{tablealt} User Authentication (\texttt{/auth/login}) & [YES] Allowed & [YES] Allowed & [YES] Allowed & [YES] Allowed \\ \hline
User Self-Registration (\texttt{/auth/register}) & [YES] Allowed & [NO] Denied & [NO] Denied & [YES] Allowed \\ \hline
\rowcolor{tablealt} Single Transaction Scoring (\texttt{/analyze}) & [NO] Denied & [YES] Full Access & [YES] Full Access & [YES] Full Access \\ \hline
Live Transaction Stream Monitoring & [NO] Denied & [YES] View \& Triage & [YES] View Only & [YES] Full Control \\ \hline
\rowcolor{tablealt} TreeSHAP Waterfall \& Feature View & [NO] Denied & [YES] View Attribution & [YES] Full Audit View & [YES] Full Audit View \\ \hline
FCRA Adverse Action Letter Download & [NO] Denied & [NO] Denied & [YES] Generate \& Export & [YES] Generate \& Export \\ \hline
\rowcolor{tablealt} Graph Network Topology \& Mule Rings & [NO] Denied & [YES] Investigate Rings & [YES] View Topology & [YES] Full Admin Control \\ \hline
MLOps Drift Status \& KS Metrics & [NO] Denied & [NO] Denied & [YES] View Metrics & [YES] Full Access \\ \hline
\rowcolor{tablealt} Synthetic Drift Injection Sandbox & [NO] Denied & [NO] Denied & [NO] Denied & [YES] Execute Spikes \\ \hline
Federated Learning Simulation Round & [NO] Denied & [NO] Denied & [NO] Denied & [YES] Trigger FedAvg \\ \hline
\rowcolor{tablealt} View Immutable Security Audit Logs & [NO] Denied & [NO] Denied & [YES] Read-Only Audit & [YES] Full Audit Access \\ \hline
\end{tabular}
\end{adjustbox}
\end{table}


\section{Feasibility Study}

A multi-dimensional feasibility study was conducted prior to system implementation to evaluate technical, operational, economic, schedule, and legal dimensions.


\subsection{Technical Feasibility}

The technical feasibility evaluates whether available algorithmic frameworks, programming languages, libraries, and hardware infrastructure can satisfy project objectives:

\begin{itemize}[leftmargin=0.35in]
\item \textbf{Machine Learning \& Ensemble Feasibility}: Python provides mature, highly optimized C++ bindings for gradient-boosted trees through \texttt{xgboost} and \texttt{lightgbm}. Both libraries support multi-threaded CPU and GPU training, built-in $L_1/L_2$ regularization, and native positive class weighting (\texttt{scale\_pos\_weight}). Scikit-learn's \texttt{StackingClassifier} offers an industrial-grade interface for training Level-0 base learners and Level-1 meta-classifiers.
\item \textbf{Explainability Feasibility}: Lundberg's \texttt{shap} C++ implementation of the TreeSHAP algorithm evaluates tree ensemble attributions in polynomial time $\mathcal{O}(T L D^2)$, reducing calculation times from minutes (KernelSHAP) to under 30 milliseconds.
\item \textbf{Microservices \& Web Feasibility}: FastAPI utilizes Starlette and Pydantic, executing asynchronous event loops powered by \texttt{uvloop}. It consistently ranks among the fastest web frameworks in TechEmpower benchmarks, rivaling Node.js and Go. On the frontend, React 19 and Vite 6 provide rapid Hot Module Replacement (HMR) and lightweight production bundles ($< 400\text{ KB}$ gzipped).
\item \textbf{Conclusion}: The project is \textbf{100\% Technically Feasible}.
\end{itemize}

\subsection{Operational Feasibility}

Operational feasibility evaluates how effectively the platform integrates into existing banking workflows:

\begin{itemize}[leftmargin=0.35in]
\item \textbf{Analyst Workflow Alignment}: FraudShield AI provides pre-configured 1-click personas in the Identity Portal, eliminating onboarding friction. Rather than forcing analysts to parse raw JSON logs, the interface presents intuitive visual risk badges, plain-English FCRA reason cards, and interactive network graphs.
\item \textbf{Human-in-the-Loop Safeguards}: By introducing the \texttt{NEEDS REVIEW} tier (\\$45\%-75\%$), the platform respects established banking standard operating procedures (SOPs), automating straight-through processing for obvious safe/fraud transactions while empowering human analysts to handle nuanced disputes.
\item \textbf{Conclusion}: The platform achieves \textbf{High Operational Feasibility}.
\end{itemize}

\subsection{Economic Feasibility}

Economic feasibility assesses development and operational expenditures against anticipated financial returns:

\begin{itemize}[leftmargin=0.35in]
\item \textbf{Development Costs}: The platform is constructed entirely using open-source technologies (Python, FastAPI, React, Vite, Docker, Scikit-learn, XGBoost, NetworkX), incurring zero software licensing fees.
\item \textbf{Infrastructure Costs}: Microservice architecture allows containerized deployment on modest cloud instances (e.g., 2 vCPUs, 4GB RAM) or free-tier hosting platforms (Render, Vercel, Docker Hub).
\item \textbf{Return on Investment (ROI)}: Commercial financial institutions lose billions annually to fraud and spend millions on manual investigation. A platform that reduces false positive rates by \\$20\%$ saves tier-1 banks millions annually in operational costs while preserving cardholder revenue.
\item \textbf{Conclusion}: The project demonstrates \textbf{Exceptional Economic Feasibility}.
\end{itemize}

\subsection{Schedule Feasibility}

The project schedule was structured across two academic semesters (Semester VI \& Semester VII) using the Agile Scrum framework across six 3-week sprint cycles. All functional milestones---from exploratory data analysis (EDA) to model training, API construction, frontend design, and automated testing---were sequenced logically with generous risk contingency buffers.

\begin{itemize}[leftmargin=0.35in]
\item \textbf{Conclusion}: The project satisfies \textbf{Strict Schedule Feasibility}.
\end{itemize}

\subsection{Legal, Ethical \& Regulatory Feasibility}

Compliance with international legal frameworks is a core design criterion of FraudShield AI:

\begin{itemize}[leftmargin=0.35in]
\item \textbf{GDPR Article 22 \& EU AI Act (2024)}: The system strictly satisfies the "Right to Explanation" by decomposing every classification into exact TreeSHAP feature attributions.
\item \textbf{FCRA Section 615(a) \& ECOA}: Automated generation of Adverse Action reason codes guarantees that consumers are not denied financial services based on discriminatory, arbitrary, or unrecorded criteria.
\item \textbf{Data Privacy \& Sovereignty}: The Federated Learning simulation module proves that cross-institutional fraud defense can operate without pooling raw transaction records or customer PII.
\item \textbf{Conclusion}: The project establishes \textbf{Complete Legal and Regulatory Feasibility}.
\end{itemize}


\section{Methodology}

FraudShield AI adopts a modified \textbf{CRISP-DM (Cross-Industry Standard Process for Data Mining)} framework tailored specifically for high-frequency, regulated financial machine learning systems. The lifecycle comprises six iterative, interconnected phases:


\begin{lstlisting}
+---------------------------------------------------------------------------------------------+
|                            MODIFIED CRISP-DM FINANCIAL AI LIFECYCLE                         |
|                                                                                             |
|  +------------------------+         +------------------------+         +-----------------+  |
|  | 1. Problem & Regulatory| ------> | 2. Multi-Domain Data   | ------> | 3. Preprocessing|  |
|  |    Intelligence Phase  |         |    Acquisition & EDA   |         |    & OmniSMOTE  |  |
|  +------------------------+         +------------------------+         +--------+--------+  |
|                                                                                 |           |
|  +------------------------+         +------------------------+                  |           |
|  | 6. Continuous MLOps &  | <------ | 5. Explainability &    | <----------------+           |
|  |    Federated Monitoring|         |    Operational Tuning  |   4. Stacking Ensemble       |
|  +------------------------+         +------------------------+      Training & Validation   |
+---------------------------------------------------------------------------------------------+
\end{lstlisting}

\subsection{Phase 1: Problem \& Regulatory Intelligence}
\begin{itemize}[leftmargin=0.35in]
\item Formalize business objectives: maximize fraud recall ($> 85\%$ on imbalanced data) while suppressing false positive rates ($< 6\%$).
\item Identify regulatory constraints: map EU AI Act transparency rules and FCRA adverse action mandates directly into algorithmic requirements.
\end{itemize}

\subsection{Phase 2: Multi-Domain Data Acquisition \& Exploratory Analysis}
\begin{itemize}[leftmargin=0.35in]
\item Ingest four benchmark datasets: PaySim Mobile Money (6.3M records), European Credit Card PCA (284k records), Spatial Behavioral Credit Card (1.8M records), and BankSim Retail Banking (594k records).
\item Analyze class imbalance distributions, covariance structures, outlier distributions, and missing value profiles.
\end{itemize}

\subsection{Phase 3: Data Preprocessing \& Advanced Resampling}
\begin{itemize}[leftmargin=0.35in]
\item Apply \texttt{RobustScaler} across continuous transaction attributes ($V_1-V_{28}$, Amount, Balance) using interquartile range (IQR) scaling:
\end{itemize}
\begin{equation}
x_{\text{scaled}} = \frac{x - \text{median}(x)}{\text{IQR}(x)} = \frac{x - Q_2}{Q_3 - Q_1}
\end{equation}
preventing extreme financial transaction outliers from distorting feature variance.

\begin{itemize}[leftmargin=0.35in]
\item Execute \textbf{OmniSMOTE} boundary-calibrated oversampling on training splits, raising minority class prevalence to balanced ratios while pruning boundary synthesis that overlaps majority clusters.
\end{itemize}

\subsection{Phase 4: Stacking Ensemble Training \& Out-of-Fold Validation}
\begin{itemize}[leftmargin=0.35in]
\item Partition datasets into stratified \\$80\%$ training and \\$20\%$ holdout testing splits.
\item Train Level-0 base estimators: XGBoost (\texttt{scale\_pos\_weight} tuned), LightGBM (\texttt{max\_depth=6}, \texttt{num\_leaves=31}), and Random Forest (\texttt{n\_estimators=100}).
\item Generate out-of-fold probability predictions using 5-fold stratified cross-validation.
\item Train Level-1 Logistic Regression meta-learner ($C=1.0$) on meta-probability vectors.
\end{itemize}

\subsection{Phase 5: Real-Time Explainability \& Decision Tuning}
\begin{itemize}[leftmargin=0.35in]
\item Pre-compute TreeSHAP explainer trees using background reference transaction matrices.
\item Formulate continuous risk calibration formulas fusing model probabilities with balance liquidation flags and velocity heuristics.
\item Establish 3-tier classification thresholds: \texttt{SAFE} ($< 45\%$), \texttt{NEEDS REVIEW} (\\$45\%-75\%$), \texttt{FRAUD} ($\ge 75\%$).
\end{itemize}

\subsection{Phase 6: Continuous MLOps, Drift Auditing \& Federated Defense}
\begin{itemize}[leftmargin=0.35in]
\item Deploy streaming Kolmogorov-Smirnov and PSI covariate shift detectors.
\item Integrate in-memory NetworkX graph cycle algorithms for mule ring detection.
\item Simulate Federated Averaging (FedAvg) rounds across banking nodes to validate privacy-preserving collaborative learning.
\end{itemize}


\section{Technology Stack}

The FraudShield AI technological architecture is constructed using modern, industry-standard languages, frameworks, and deployment containers:


\subsubsection{Table 3.5: Comprehensive Technology Stack Specification}

\begin{table}[H]
\centering
\footnotesize
\begin{adjustbox}{width=\textwidth,center}
\begin{tabular}{| l | l | l | l |}
\hline
\rowcolor{tableheader} \textbf{Subsystem Layer} & \textbf{Technology / Tool} & \textbf{Version} & \textbf{Architectural Role \& Justification} \\ \hline
\textbf{Frontend Framework} & \textbf{React} & 19.0.0 & High-performance single-page application (SPA) library utilizing concurrent rendering and virtual DOM diffing. \\ \hline
\rowcolor{tablealt} \textbf{Build \& Bundler} & \textbf{Vite} & 6.0.0 & Next-generation frontend tooling providing sub-second Hot Module Replacement (HMR) and optimized Rollup builds. \\ \hline
\textbf{CSS \& Styling} & \textbf{Tailwind CSS} & 3.4.1 & Utility-first CSS framework enabling dark-mode glassmorphic aesthetics and responsive layouts without stylesheet bloat. \\ \hline
\rowcolor{tablealt} \textbf{Micro-Animations} & \textbf{Framer Motion} & 11.0.0 & Production-grade motion library for fluid UI transitions, expandable transaction triage drawers, and modal overlays. \\ \hline
\textbf{Network Visualization} & \textbf{SVG / Canvas} & Native & Interactive rendering of transaction topologies, circular money laundering rings, and account node clusters. \\ \hline
\rowcolor{tablealt} \textbf{Backend API Framework} & \textbf{FastAPI} & 0.115.0+ & Modern, asynchronous ASGI web framework built on Starlette and Pydantic with automatic OpenAPI documentation. \\ \hline
\textbf{ASGI Web Server} & \textbf{Uvicorn} & 0.30.0+ & Lightning-fast asynchronous server implementation for Python based on \texttt{uvloop} and \texttt{httptools}. \\ \hline
\rowcolor{tablealt} \textbf{Data Validation} & \textbf{Pydantic} & 2.8.0+ & Robust data validation and schema enforcement utilizing Python type hints with zero-overhead runtime parsing. \\ \hline
\textbf{Machine Learning Core} & \textbf{Scikit-Learn} & 1.5.0+ & Core machine learning library providing \texttt{StackingClassifier}, \texttt{RobustScaler}, metrics benchmarking, and model persistence. \\ \hline
\rowcolor{tablealt} \textbf{Gradient Boosting (1)} & \textbf{XGBoost} & 2.1.0+ & Extreme Gradient Boosting library optimized for fast parallel tree construction and positive class weighting. \\ \hline
\textbf{Gradient Boosting (2)} & \textbf{LightGBM} & 4.4.0+ & High-performance gradient boosting framework utilizing histogram binning, GOSS, and leaf-wise tree growth. \\ \hline
\rowcolor{tablealt} \textbf{Imbalanced Learning} & \textbf{Imbalanced-Learn} & 0.12.0+ & Advanced resampling suite providing SMOTE, Borderline-SMOTE, and ADASYN implementations. \\ \hline
\textbf{Explainable AI (XAI)} & \textbf{SHAP (TreeSHAP)} & 0.46.0+ & Game-theoretic feature attribution library providing exact polynomial-time TreeSHAP calculations. \\ \hline
\rowcolor{tablealt} \textbf{Graph Intelligence} & \textbf{NetworkX} & 3.3.0+ & Comprehensive Python graph library for building transaction networks, finding cycles, and computing degree centrality. \\ \hline
\textbf{Security \& JWT} & \textbf{PyJWT} & 2.8.0+ & Cryptographic JSON Web Token implementation for stateless, secure session authentication and RBAC claims. \\ \hline
\rowcolor{tablealt} \textbf{Containerization} & \textbf{Docker \& Compose} & 26.0+ & Container virtualization ensuring consistent, reproducible deployment across development, staging, and production. \\ \hline
\textbf{Automated Testing} & \textbf{Pytest} & 8.2.0+ & Mature Python testing framework used for unit testing, endpoint integration testing, and regression suites. \\ \hline
\end{tabular}
\end{adjustbox}
\end{table}


\section{Gantt Chart and Process Model}

\subsection{Agile Scrum Process Model}

FraudShield AI was developed utilizing the \textbf{Agile Scrum} methodology. Development was divided into six two-week sprints grouped into two macro project phases:


\subsubsection{Table 3.6: Agile Sprint Backlog, Story Points, and Delivery Milestones}

\begin{table}[H]
\centering
\footnotesize
\begin{adjustbox}{width=\textwidth,center}
\begin{tabular}{| l | l | l | l | l |}
\hline
\rowcolor{tableheader} \textbf{Sprint No.} & \textbf{Sprint Focus \& Core Deliverables} & \textbf{Story Points} & \textbf{Duration} & \textbf{Milestones Achieved} \\ \hline
\textbf{Sprint 1} & Problem Formulation, Regulatory Audit, Dataset Acquisition & 25 pts & Weeks 1--3 & Project Dossier, Dataset Ingestion Pipeline, Baseline EDA. \\ \hline
\rowcolor{tablealt} \textbf{Sprint 2} & Data Preprocessing, RobustScaler, OmniSMOTE Resampling & 35 pts & Weeks 4--6 & Cleaned multi-domain datasets, balanced training partitions. \\ \hline
\textbf{Sprint 3} & Stacking Ensemble Training \& Hyperparameter Optimization & 45 pts & Weeks 7--9 & Level-0 XGB/LGBM/RF models, Level-1 Meta-Learner, 5-fold CV. \\ \hline
\rowcolor{tablealt} \textbf{Sprint 4} & TreeSHAP Integration, FCRA Reason Codes, Risk Calibration & 40 pts & Weeks 10--12 & Sub-50ms TreeSHAP engine, 3-tier calibrated risk scoring. \\ \hline
\textbf{Sprint 5} & FastAPI Microservice Gateway, RBAC, Security \& Audit Logs & 35 pts & Weeks 13--15 & Complete REST API endpoints (\texttt{/analyze}, \texttt{/health}, \texttt{/auth}). \\ \hline
\rowcolor{tablealt} \textbf{Sprint 6} & React 19 Frontend SPA, Live Stream, Network Graph, Docker & 40 pts & Weeks 16--18 & Full-stack containerized platform, live stream simulator, Blue Book. \\ \hline
\end{tabular}
\end{adjustbox}
\end{table}

\subsection{Gantt Chart Schedule}

The visual timeline illustrating task sequencing, dependency relationships, and milestone completions across the academic cycle is shown below:


\begin{lstlisting}
+---------------------------------------------------------------------------------------------+
|                            FRAUDSHIELD AI PROJECT GANTT CHART                               |
|                                                                                             |
| Phase / Task Name            W1  W2  W3  W4  W5  W6  W7  W8  W9  W10 W11 W12 W13 W14 W15 W16|
| ------------------------------------------------------------------------------------------- |
| 1. Problem Def & Literature  [######]                                                       |
| 2. Data Acquisition & EDA        [######]                                                   |
| 3. Preprocessing & Resample          [######]                                               |
| 4. Stacking Model Training               [##########]                                       |
| 5. TreeSHAP & Risk Engine                        [##########]                               |
| 6. FastAPI Backend & RBAC                                [##########]                       |
| 7. React 19 Frontend UI                                      [##########]                   |
| 8. Graph & Drift Modules                                             [##########]           |
| 9. Automated Testing & Docker                                                [######]       |
| 10. Blue Book Documentation                                                      [########] |
+---------------------------------------------------------------------------------------------+
\end{lstlisting}


\section{System Analysis (Functional, Structural, and Behavioral Models)}

System analysis provides formal object-oriented and structural models representing system capabilities, entity relationships, and temporal interaction lifecycles.


\subsection{Functional Model: Use Case Diagram \& Specifications}

The Use Case model defines interactions between external human/machine actors and system capabilities:

\begin{itemize}[leftmargin=0.35in]
\item \textbf{Actors}:
\end{itemize}
\begin{enumerate}[leftmargin=0.35in]
\item \texttt{External Core Banking / Webhook}: Machine actor submitting live transaction payloads via API.
\item \texttt{SOC Fraud Analyst}: Human security operator monitoring live transaction feeds, triaging alerts, inspecting single transactions, and investigating mule rings.
\item \texttt{Compliance Officer}: Regulatory auditor reviewing TreeSHAP attributions, generating FCRA Adverse Action letters, and inspecting audit logs.
\item \texttt{System Administrator}: Technical operator executing synthetic drift spikes, triggering federated learning rounds, managing API keys, and managing user accounts.
\end{enumerate}

\begin{lstlisting}
                               USE CASE DIAGRAM: FRAUDSHIELD AI

     +-----------------+                                                 +-----------------+
     |  Banking Core / |                                                 |   SOC Analyst   |
     |     Webhook     |                                                 |     Actor       |
     +--------+--------+                                                 +--------+--------+
              |                                                                   |
              |--> [UC-01: Ingest Transaction Payload]                           |
              |--> [UC-02: Synchronous Fraud Risk Scoring] <----------------------+
              |                                                                   |
              |                                      [UC-03: View Live Stream] <--+
              |                                      [UC-04: Single Txn Triage]<--+
              |                                      [UC-05: Inspect TreeSHAP] <--+
              |                                      [UC-06: Investigate Mule] <--+
              |                                                                   |
     +--------+--------+                                                 +--------+--------+
     |   Compliance    |                                                 |  System Admin   |
     |     Officer     |                                                 |     Actor       |
     +--------+--------+                                                 +--------+--------+
              |                                                                   |
              |--> [UC-07: Generate FCRA Adverse Action Letter]                  |
              |--> [UC-08: Audit Immutable Security Trail] <----------------------+
              |--> [UC-09: Monitor Concept Drift Metrics]  <----------------------+
              |                                                                   |
              |                                      [UC-10: Trigger FedAvg]   <--+
              |                                      [UC-11: Inject Drift Spike]<-+
              |                                      [UC-12: Manage User RBAC] <--+
\end{lstlisting}


\subsection{Structural Model: Unified Class Diagram}

The class diagram illustrates the object-oriented structure of the backend analytical engines, data schemas, and API gateway:


\begin{lstlisting}
+---------------------------------+                 +---------------------------------+
|        TransactionPayload       |                 |         PredictionResult        |
+---------------------------------+                 +---------------------------------+
| + domain: str                   |                 | + transaction_id: str           |
| + raw_features: dict            |                 | + raw_probability: float        |
| + timestamp: datetime           |                 | + risk_score: float             |
| + ip_address: Optional[str]     |                 | + decision_tier: str            |
+---------------------------------+                 | + is_fraud: bool                |
| + validate_schema(): bool       |                 | + latency_ms: float             |
| + to_feature_vector(): ndarray  |                 +---------------------------------+
+---------------+-----------------+                                  ^
                |                                                    |
                v                                                    | generates
+---------------------------------+                                  |
|        ModelEngineRegistry      |                                  |
+---------------------------------+                                  |
| - model_heads: dict             |                                  |
| - meta_learners: dict           |                                  |
| - scalers: dict                 |                                  |
+---------------------------------+                                  |
| + detect_domain(data): str      |                                  |
| + predict(payload): Result -----+----------------------------------+
| + get_model_metrics(): dict     |
+---------------+-----------------+
                | delegates to
                v
+---------------------------------+                 +---------------------------------+
|        StackingEnsemble         |                 |         TreeSHAPExplainer       |
+---------------------------------+                 +---------------------------------+
| - xgb_model: XGBClassifier      |                 | - explainer: TreeExplainer      |
| - lgb_model: LGBMClassifier     |                 | - background_data: ndarray      |
| - rf_model: RandomForestClass   |                 +---------------------------------+
| - meta_learner: LogisticRegress |                 | + explain_instance(x): dict     |
+---------------------------------+                 | + get_waterfall_plot(x): dict   |
| + fit(X, y): void               |                 | + map_fcra_reason_codes(): list |
| + predict_proba(X): ndarray     |                 +---------------------------------+
+---------------------------------+
\end{lstlisting}


\subsection{Behavioral Model: Sequence Diagram}

The sequence diagram illustrates the temporal execution flow when a client or analyst submits a transaction for scoring, explainability, and graph lookup:


\begin{lstlisting}
 Client / UI              FastAPI Gateway           ModelEngine           TreeSHAP           NetworkX Graph
      |                          |                       |                    |                     |
      |-- POST /analyze (JSON) ->|                       |                    |                     |
      |                          |-- validate JWT/Auth ->|                    |                     |
      |                          |-- detect schema ----->|                    |                     |
      |                          |                       |                    |                     |
      |                          |-- predict_proba() --->|                    |                     |
      |                          |                       |-- L0 Base Predict -|                     |
      |                          |                       |-- L1 Meta Predict -|                     |
      |                          |                       |<- raw prob (p) ----|                     |
      |                          |                       |                    |                     |
      |                          |-- compute risk score -|                    |                     |
      |                          |                       |-- explain(x) ----->|                     |
      |                          |                       |<- shap values (phi)|                     |
      |                          |                       |                    |                     |
      |                          |-- check cycles/mule -------------------------------------------->|
      |                          |<-- return mule score & cycle flag -------------------------------|
      |                          |                       |                    |                     |
      |                          |-- compile JSON result |                    |                     |
      |<-- HTTP 200 (Risk, XAI) -|                       |                    |                     |
\end{lstlisting}


\textit{End of Chapter 3 --- Requirement Gathering, Analysis and Planning}


\newpage


\chapter{System Design and Experimental Set up}

\section{System Architecture \& Diagrams}

System design translates the functional, performance, and regulatory requirements formulated in Chapter 3 into a concrete, robust engineering specification. The architecture of \textbf{FraudShield AI} is engineered around principles of loose coupling, high cohesion, asynchronous execution, and mathematical transparency.


\subsection{Data Flow Diagrams (DFD)}

Data Flow Diagrams model the movement, transformation, and storage of financial transaction data across system boundaries.


\subsubsection{DFD Level 0 (Context Diagram)}

The Level 0 context diagram represents the overall system boundary, external entities, and high-level data conduits:


\begin{lstlisting}
+--------------------------------+                               +--------------------------------+
|   External Banking Core /      | --- Transaction Ingestion --> |                                |
|      Merchant Webhook          | <--- Auth Flag & Risk Score -- |                                |
+--------------------------------+                               |                                |
                                                                 |                                |
+--------------------------------+                               |                                |
|                                | --- Analyst Override & Triage>|       FRAUDSHIELD AI           |
|       SOC Fraud Analyst        | <--- Stream Feed & Network ---|    INTELLIGENT DECISION        |
|                                |                               |          PLATFORM              |
+--------------------------------+                               |                                |
                                                                 |                                |
+--------------------------------+                               |                                |
|                                | --- Audit & Dispute Inquiries>|                                |
|       Compliance Officer       | <--- TreeSHAP & FCRA Reason -- |                                |
+--------------------------------+                               +--------------------------------+
\end{lstlisting}

\subsubsection{DFD Level 1 (Modular Functional Decomposition)}

The Level 1 DFD decomposes the system into six primary process components and persistent data stores:


\begin{lstlisting}
                      [ Incoming Transaction Payload ]
                                     |
                                     v
                      +------------------------------+
                      |    1.0 Ingestion & Security  |
                      |         (FastAPI Auth)       |
                      +--------------+---------------+
                                     | Validated Data
                                     v
                      +------------------------------+
                      |    2.0 Schema Detection &    |
                      |        Domain Routing        |
                      +--------------+---------------+
                                     | Normalized Features
                                     v
        +--------------------------------------------------------+
        |                                                        |
        v                                                        v
+------------------------------+                         +------------------------------+
|  3.0 Stacking Ensemble       |                         |  4.0 Explainable AI          |
|      Predictive Engine       | -- Raw Probability (p) -|      (TreeSHAP Engine)       |
+--------------+---------------+                         +--------------+---------------+
               |                                                        | Shapley Values (phi)
               v                                                        v
+------------------------------+                         +------------------------------+
|  5.0 Calibrated Continuous   |                         |  6.0 Graph Intelligence &    |
|      Risk Engine (3-Tier)    |                         |      Mule Ring Analyzer      |
+--------------+---------------+                         +--------------+---------------+
               |                                                        |
               +------------------------+-------------------------------+
                                        | Composite Decision & XAI Metadata
                                        v
                         +------------------------------+
                         |   Client Response Delivery   |
                         |    & Immutable Audit Log     |
                         +------------------------------+
\end{lstlisting}

\subsubsection{DFD Level 2 (In-Depth Analytical Scoring \& Graph Process)}

The Level 2 DFD details the internal computational mechanics of Process 3.0, 4.0, and 6.0:

\begin{enumerate}[leftmargin=0.35in]
\item \textbf{Process 3.1 (Feature Scaling)}: Raw continuous attributes pass through domain-specific \texttt{RobustScaler} objects stored in \texttt{backend/models/*.pkl}.
\item \textbf{Process 3.2 (Base Learner Inference)}: Scaled feature vector $\mathbf{x}$ is evaluated simultaneously by XGBoost, LightGBM, and Random Forest base models, generating three probability vectors: $\hat{p}_{\text{xgb}}, \hat{p}_{\text{lgb}}, \hat{p}_{\text{rf}}$.
\item \textbf{Process 3.3 (Meta-Learner Fusion)}: The meta-feature vector $[\hat{p}_{\text{xgb}}, \hat{p}_{\text{lgb}}, \hat{p}_{\text{rf}}]$ is fed into the Level-1 Logistic Regression classifier to yield the unified posterior probability $p = \sigma(\mathbf{w}^T \mathbf{p}_{\text{meta}} + b)$.
\item \textbf{Process 4.1 (TreeSHAP Evaluation)}: The pre-compiled TreeExplainer evaluates feature splits across trees, calculating exact marginal Shapley contributions $\phi_i$.
\item \textbf{Process 4.2 (Reason Code Mapping)}: Features with $\phi_i > 0$ are sorted in descending order of magnitude. The top three features are cross-referenced with the FCRA Reason Code dictionary to yield human-readable descriptions.
\item \textbf{Process 6.1 (Graph Ingestion \& Cycle Traversal)}: Source account and destination account are added as directed edges to the NetworkX graph. Tarjan's Strongly Connected Components and depth-first search (DFS) identify if the transaction completes a cycle ($A
\end{enumerate}
ightarrow B

ightarrow C

ightarrow A$).

\begin{enumerate}[leftmargin=0.35in]
\item \textbf{Process 6.2 (Mule Risk Scoring)}: The account's out-degree to in-degree ratio, transaction velocity in the preceding 60 minutes, and balance liquidation delta are synthesized into an integer Mule Score (\\$0-100$).
\end{enumerate}


\subsection{Microservices \& Network Deployment Topology}

FraudShield AI is deployed as containerized microservices orchestrated via Docker Compose. The topology isolates the frontend presentation layer, API gateway, and analytical engines:


\begin{lstlisting}
+---------------------------------------------------------------------------------------------+
|                                 EDGE NETWORK / CLIENT LAYER                                 |
|                                                                                             |
|       Browser Client (Analyst / Auditor)                  External Core Banking API         |
|                        |                                               |                    |
|                        v                                               v                    |
|                 [ HTTPS : 443 ]                                 [ HTTPS : 443 ]             |
+------------------------+-----------------------------------------------+--------------------+
                         |                                               |
                         v                                               v
+---------------------------------------------------------------------------------------------+
|                               DOCKER CONTAINER ORCHESTRATION                                |
|                                                                                             |
|  +----------------------------------------------+                                           |
|  |  Container 1: Frontend SPA (Nginx / Vite)   |                                           |
|  |  Port 3000 -> 80                             |                                           |
|  |  - React 19 UI Modules & Dashboards         |                                           |
|  |  - Static Asset Gzip Compression             |                                           |
|  +----------------------+-----------------------+                                           |
|                         | Internal HTTP Proxy                                               |
|                         v                                                                   |
|  +---------------------------------------------------------------------------------------+  |
|  |  Container 2: Backend API Gateway (FastAPI / Uvicorn)                                 |  |
|  |  Port 8008 -> 8008                                                                    |  |
|  |  +-------------------------------+     +-------------------------------------------+  |  |
|  |  | Security & Auth (JWT / RBAC)  |     | Sliding-Window Anti-DDoS Rate Limiter     |  |  |
|  |  +---------------+---------------+     +-------------------------------------------+  |  |
|  |                  |                                                                    |  |
|  |                  v                                                                    |  |
|  |  +---------------------------------------------------------------------------------+  |  |
|  |  | Analytical Engine Runtime (Python 3.12 Asynchronous Workers)                    |  |  |
|  |  | - ModelEngineRegistry (PaySim, CreditCard PCA, Spatial, BankSim)                |  |  |
|  |  | - Stacking Ensemble (XGBoost, LightGBM, Random Forest, Logistic Regression)     |  |  |
|  |  | - Explainable AI (TreeSHAP Exact Feature Attribution)                           |  |  |
|  |  | - Graph Engine (NetworkX In-Memory Mule Ring Topology)                          |  |  |
|  |  | - Drift Monitor (Kolmogorov-Smirnov & Population Stability Index)               |  |  |
|  |  | - Federated Simulator (FedAvg Multi-Bank Collaborative Training)                |  |  |
|  |  +---------------------------------------------------------------------------------+  |  |
|  +---------------------------------------------------------------------------------------+  |
+---------------------------------------------------------------------------------------------+
\end{lstlisting}


\section{Algorithm \& Process Flow Design}

This section formalizes the core mathematical algorithms powering FraudShield AI, presenting complete procedural pseudocode and design specifications.


\subsection{Algorithm 1: OmniSMOTE Resampling \& RobustScaler Pipeline}

\begin{lstlisting}
Algorithm 1: OmniSMOTE Resampling & RobustScaler Preprocessing
Input: Training dataset D_train = {(x_k, y_k)}_{k=1}^N, where x_k in R^m, y_k in {0, 1}
       Target minority sampling ratio r = 0.50, Nearest neighbors k_neighbors = 5
Output: Scaled and balanced training dataset D_bal, Fitted scaler S

1: Initialize S <- RobustScaler()
2: For each feature dimension j in {1, ..., m}:
3:     Compute Median Q_2(j) and Interquartile Range IQR(j) = Q_3(j) - Q_1(j)
4:     For each sample x_k:
5:         x_k[j] <- (x_k[j] - Q_2(j)) / (IQR(j) + eps)
6: End For

7: Partition D_train into Majority Class D_0 and Minority Class D_1
8: Compute required synthetic samples: N_syn = floor(r * |D_0|) - |D_1|
9: Initialize Synthetic Set D_syn <- {}

10: Construct k-d Tree index on D_1
11: For step = 1 to N_syn:
12:     Randomly select sample x_i in D_1
13:     Find k-nearest neighbors N_k(x_i) within D_1
14:     Randomly select neighbor x_nn in N_k(x_i)
15:     Sample random interpolation weight lambda ~ Uniform(0, 1)
16:     Compute candidate: x_new = x_i + lambda * (x_nn - x_i)
17:     
18:     // OmniSMOTE Topological Boundary Filter:
19:     Find nearest neighbor x_maj in D_0 to x_new
20:     If Euclidean_Distance(x_new, x_i) < Euclidean_Distance(x_new, x_maj):
21:         D_syn <- D_syn union {(x_new, 1)}
22:     Else:
23:         Reject x_new and continue
24:     End If
25: End For

26: D_bal <- D_0 union D_1 union D_syn
27: Return D_bal, S
\end{lstlisting}


\subsection{Algorithm 2: Hybrid Stacking Ensemble Training \& Out-of-Fold Meta-Learning}
\begin{figure}[H]
\centering
\safeincludegraphics[width=0.92\textwidth]{fig15_stacking.jpg}
\caption{Stacking Ensemble Architecture with XGBoost, LightGBM, Random Forest, and Meta-Learner}
\label{fig:stacking}
\end{figure}

\begin{lstlisting}
Algorithm 2: Hybrid Stacking Ensemble with Cross-Validated Meta-Learner
Input: Balanced training data D_bal = {(x_i, y_i)}_{i=1}^M, Number of CV folds K = 5
       Base estimators: M_1 (XGBoost), M_2 (LightGBM), M_3 (Random Forest)
       Meta-learner: M_meta (Logistic Regression)
Output: Fitted Stacking Ensemble Model M_stack

1: Compute positive class scale weight: w_pos = |{y_i = 0}| / |{y_i = 1}|
2: Configure M_1 with scale_pos_weight = w_pos
3: Configure M_2 with scale_pos_weight = w_pos

4: Partition D_bal into K stratified folds: F_1, F_2, ..., F_K
5: Initialize Out-of-Fold meta-feature matrix P_meta in R^{M x 3}

6: For each fold k in {1, ..., K}:
7:     D_train_k <- D_bal \ F_k
8:     D_val_k   <- F_k
9:     
10:    Fit temporary estimators M_1^(k), M_2^(k), M_3^(k) on D_train_k
11:    For each sample x_j in D_val_k:
12:        p_1 <- M_1^(k).predict_proba(x_j)[1]
13:        p_2 <- M_2^(k).predict_proba(x_j)[1]
14:        p_3 <- M_3^(k).predict_proba(x_j)[1]
15:        P_meta[j, :] <- [p_1, p_2, p_3]
16:    End For
17: End For

18: Train Meta-Learner M_meta on (P_meta, y) minimizing L2-regularized log-loss:
19:     min_w { - sum [ y_i ln sigma(w^T p_i) + (1-y_i) ln(1 - sigma(w^T p_i)) ] + (1/2C) ||w||_2^2 }

20: Retrain base models M_1, M_2, M_3 on full dataset D_bal
21: Construct M_stack <- {Base: [M_1, M_2, M_3], Meta: M_meta}
22: Return M_stack
\end{lstlisting}


\subsection{Algorithm 3: Continuous Calibrated Risk Scoring \& 3-Tier Classification}

\begin{lstlisting}
Algorithm 3: Calibrated Continuous Risk Scoring & Decision Tier Assignment
Input: Incoming transaction feature vector x, Stacking Ensemble Model M_stack
Output: RiskScore in [0.0, 100.0], DecisionTier in {SAFE, NEEDS_REVIEW, FRAUD}, Action

1: Extract base Level-0 probabilities:
2:     p_xgb <- M_stack.M_1.predict_proba(x)[1]
3:     p_lgb <- M_stack.M_2.predict_proba(x)[1]
4:     p_rf  <- M_stack.M_3.predict_proba(x)[1]
5: 
6: Compute ensemble meta-probability:
7:     p <- M_stack.M_meta.predict_proba([p_xgb, p_lgb, p_rf])[1]

8: Extract contextual domain signals:
9:     amt_factor <- min(1.0, log10(x.amount + 1.0) / 6.0)
10:    is_drain   <- 1 if (x.oldbalanceOrig > 0 and x.newbalanceOrig == 0) else 0
11:    has_error  <- 1 if (abs(x.oldbalanceOrig - x.amount - x.newbalanceOrig) > 1.0) else 0

12: If p > 0.50:
13:     // High-Risk Range Calibration:
14:     RiskScore <- 0.75 + 0.08 * ((p - 0.50) / 0.50) + 0.08 * amt_factor + 0.08 * is_drain
15: Else:
16:     If x.type not in ["TRANSFER", "CASH_OUT"]:
17:         // Non-liquidating payment channels:
18:         RiskScore <- 0.15 * amt_factor
19:     Else:
20:         // Liquidating channels with low base probability:
21:         RiskScore <- 0.15 + 0.25 * amt_factor + 0.20 * is_drain + 0.10 * has_error + 0.05 * (p / 0.50)
22:     End If
23: End If

24: Clamp RiskScore <- min(1.0, max(0.0, RiskScore)) * 100.0

25: If RiskScore < 45.0:
26:     DecisionTier <- "SAFE"
27:     Action       <- "AUTO_APPROVE"
28: Else If RiskScore < 75.0:
29:     DecisionTier <- "NEEDS_REVIEW"
30:     Action       <- "ROUTE_TO_ANALYST_QUEUE_OR_STEP_UP_2FA"
31: Else:
32:     DecisionTier <- "FRAUD"
33:     Action       <- "BLOCK_TRANSACTION_AND_FREEZE_CARD"
34: End If

35: Return RiskScore, DecisionTier, Action
\end{lstlisting}


\subsection{Algorithm 4: Real-Time TreeSHAP Feature Attribution \& FCRA Mapping}
\begin{figure}[H]
\centering
\safeincludegraphics[width=0.92\textwidth]{fig17_treeshap.jpg}
\caption{TreeSHAP Feature Attribution Computation Workflow and FCRA Reason Code Mapping}
\label{fig:treeshap}
\end{figure}

\begin{lstlisting}
Algorithm 4: Real-Time TreeSHAP Feature Attribution & FCRA Mapping
Input: Transaction vector x in R^m, Ensemble Tree Model M_tree, Background matrix X_bg
Output: Attribution vector phi in R^m, List of FCRA Reason Codes R_fcra

1: Initialize TreeExplainer E(M_tree, data = X_bg)
2: Compute exact Shapley values vector:
3:     phi <- E.shap_values(x)
4: Verify Additive Efficiency: abs(sum(phi) - (M_tree(x) - E.expected_value)) < 1e-4

5: Identify positive risk drivers:
6:     D_pos <- {(j, phi[j]) | phi[j] > 0}
7: Sort D_pos in descending order by phi[j]

8: Initialize R_fcra <- []
9: For each (feature_idx, weight) in top 3 elements of D_pos:
10:    feature_name <- feature_names[feature_idx]
11:    
12:    Match feature_name:
13:        Case "errorBalanceOrig":
14:            R_fcra.append({"Code": "RC-BAL-01", "Name": "Origin Balance Mismatch", "Weight": weight})
15:        Case "is_drained":
16:            R_fcra.append({"Code": "RC-BAL-02", "Name": "Full Account Liquidation", "Weight": weight})
17:        Case "amount":
18:            R_fcra.append({"Code": "RC-AMT-01", "Name": "Abnormal Transaction Amount", "Weight": weight})
19:        Case "distance_km":
20:            R_fcra.append({"Code": "RC-GEO-01", "Name": "Haversine Distance Discrepancy", "Weight": weight})
21:        Default:
22:            R_fcra.append({"Code": "RC-PCA-01", "Name": "Latent Anomaly Component " + feature_name, "Weight": weight})
23:    End Match
24: End For

25: Return phi, R_fcra
\end{lstlisting}


\subsection{Algorithm 5: In-Memory NetworkX Cyclic Mule Ring Detection}
\begin{figure}[H]
\centering
\safeincludegraphics[width=0.92\textwidth]{fig18_mule_ring.jpg}
\caption{Graph Intelligence Engine: Cyclic Mule Ring and Community Detection Topology}
\label{fig:mulering}
\end{figure}

\begin{lstlisting}
Algorithm 5: In-Memory NetworkX Cyclic Mule Ring & Centrality Analysis
Input: Directed Multi-Graph G = (V, E), New Transaction T = (u, v, amount, timestamp)
Output: IsMuleRing in {True, False}, CyclePath list, MuleScore in [0, 100]

1: Add directed edge: G.add_edge(u, v, weight = amount, time = timestamp)
2: 
3: // Cyclic Path Detection (k-hop DFS from target node v back to source u):
4: Initialize CyclePath <- []
5: IsMuleRing <- False
6: 
7: If NetworkX.has_path(G, source = v, target = u):
8:     All_Simple_Paths <- NetworkX.all_simple_paths(G, source = v, target = u, cutoff = 5)
9:     For path in All_Simple_Paths:
10:        If length(path) >= 2:
11:            CyclePath <- [u] + path
12:            IsMuleRing <- True
13:            Break
14:        End If
15:    End For
16: End If

17: // Account Mule Risk Score Calculation for node u:
18: deg_in  <- G.in_degree(u)
19: deg_out <- G.out_degree(u)
20: total_vol <- sum([edge.weight for edge in G.out_edges(u)])
21: 
22: // Mule ratio: rapid pass-through of funds
23: ratio <- min(deg_in, deg_out) / (max(deg_in, deg_out) + 1.0)
24: 
25: MuleScore <- 0
26: If IsMuleRing: MuleScore <- MuleScore + 50
27: If deg_in > 5 and deg_out > 5: MuleScore <- MuleScore + 25
28: If ratio > 0.60: MuleScore <- MuleScore + 25
29: 
30: Return IsMuleRing, CyclePath, MuleScore
\end{lstlisting}


\subsection{Algorithm 6: Decentralized Federated Learning (FedAvg) Simulation}

\begin{lstlisting}
Algorithm 6: Federated Averaging (FedAvg) Multi-Bank Consortium Aggregation
Input: K participating banking institutions, Local datasets {D_k}_{k=1}^K
       Total records N = sum_{k=1}^K |D_k|, Global aggregation rounds T = 5
Output: Global Defense Model W_global

1: Initialize Global Model Parameter Vector W_0
2: For round t = 0 to T - 1:
3:     For each banking node k in {1, ..., K} in parallel:
4:         Download current global weights: W_k^(t) <- W_t
5:         
6:         // Local Model Training on Internal Bank Data:
7:         Execute local optimization (e.g. 5 epochs SGD or local tree boosting) on D_k
8:         Obtain updated local parameters: W_k^(t+1)
9:     End For
10:    
11:    // Central Coordinator Aggregation (Zero PII Sharing):
12:    W_{t+1} <- sum_{k=1}^K (|D_k| / N) * W_k^(t+1)
13:    
14:    Compute Consortium Global F1-Score on independent validation partition
15: End For
16: Return W_T
\end{lstlisting}


\section{User Interface \& Input Data Design}

\subsection{User Interface Design Specifications}

FraudShield AI provides an enterprise-grade visual command center built with \textbf{React 19 + Vite 6 + Tailwind CSS}. The UI adheres to modern human-computer interaction (HCI) standards:

\begin{enumerate}[leftmargin=0.35in]
\item \textbf{Curated Color Tokens}:
\end{enumerate}
\begin{itemize}[leftmargin=0.35in]
\item \texttt{Background}: Deep Midnight Slate (\texttt{\#0B0F19})
\item \texttt{Card Surface}: Glassmorphic Charcoal with \\$1\text{px}$ subtle border (\texttt{\#111827}, border \texttt{\#1F2937})
\item \texttt{Risk Accents}: [SAFE] Emerald (\texttt{\#10B981}), [REVIEW] Amber (\texttt{\#F59E0B}), [FRAUD] Crimson (\texttt{\#EF4444}), [INFO] Cyan/Violet (\texttt{\#06B6D4}, \texttt{\#8B5CF6})
\end{itemize}
\begin{enumerate}[leftmargin=0.35in]
\item \textbf{Typography}: Set in Google Fonts \texttt{Inter} and \texttt{Outfit} with crisp monospace rendering for financial hashes and account numbers.
\item \textbf{Micro-Animations}: Powered by Framer Motion, providing instant visual feedback for risk tier transitions, expandable triage drawers, and modal overlays.
\end{enumerate}

\subsection{Screen Walkthroughs and Functional Capabilities}

\begin{itemize}[leftmargin=0.35in]
\item \textbf{1. Threat Intelligence Dashboard (\texttt{DashboardModule.jsx})}:
\end{itemize}
Provides an executive snapshot featuring four animated KPI counter cards: Total Transactions Analyzed (6.3M+), Active Fraud Detection Rate (\\$99.66\%$), Total Blocked Losses ($\\\$18.4\text{M}$), and Average Inference Latency (\\$42\text{ms}$). Includes an interactive Recharts volume curve and real-time domain breakdown donuts.

\begin{itemize}[leftmargin=0.35in]
\item \textbf{2. Live Stream Transaction Monitor (\texttt{LiveMonitoringModule.jsx})}:
\end{itemize}
Simulates live terminal ingestion. New transactions slide smoothly into a virtualized list. Each row displays the timestamp, account origin, destination, amount, dynamic risk badge, and a "Triage" button. Clicking a transaction opens a slide-over drawer showing feature values and one-click block/approve actions.

\begin{itemize}[leftmargin=0.35in]
\item \textbf{3. Single Transaction Inspector (\texttt{TransactionAnalysisModule.jsx})}:
\end{itemize}
A diagnostic sandbox allowing fraud investigators to select any of the four domain schemas, pre-populate benchmark test cases (e.g., Safe Merchant Payment vs. Account Drain Attack), modify individual feature values via interactive sliders, and trigger instant synchronous scoring.

\begin{itemize}[leftmargin=0.35in]
\item \textbf{4. Explainable AI \& FCRA Studio (\texttt{ExplainableAIModule.jsx})}:
\end{itemize}
Renders the interactive TreeSHAP waterfall plot. Displays the base value $\phi_0$, individual feature force vectors (red positive risk bars, green mitigating bars), and generates an official, downloadable \textbf{FCRA Adverse Action Disclosure Statement} in PDF/text format for regulatory audit compliance.

\begin{itemize}[leftmargin=0.35in]
\item \textbf{5. Graph Intelligence \& Mule Ring Forensics (\texttt{GraphIntelligenceModule.jsx})}:
\end{itemize}
Renders an interactive directed network topology. Highlights cyclic edges in glowing crimson ($A \rightarrow B \rightarrow C \rightarrow A$) and computes node degree centrality to identify money laundering hubs.

\begin{itemize}[leftmargin=0.35in]
\item \textbf{6. MLOps Concept Drift Portal (\texttt{ConceptDriftModule.jsx})}:
\end{itemize}
Displays real-time two-sample Kolmogorov-Smirnov test statistics and Population Stability Index (PSI) histograms. Features a "Drift Spike Simulator" button allowing administrators to inject synthetic adversarial distributions and witness automated alerting.

\begin{itemize}[leftmargin=0.35in]
\item \textbf{7. Federated Learning Simulator (\texttt{FederatedLearningModule.jsx})}:
\end{itemize}
Visualizes decentralized collaborative training across Bank Alpha, Bank Beta, and Bank Gamma. Shows local F1-scores, federated aggregation round progress, and confirms zero transmission of customer PII.

\begin{itemize}[leftmargin=0.35in]
\item \textbf{8. Identity Portal \& RBAC Switching (\texttt{AuthModal.jsx})}:
\end{itemize}
Allows 1-click persona switching between Sarah Chen (SOC Analyst), Marcus Vance (Compliance Officer), and Dr. Elena Rostova (Admin), dynamically tailoring sidebar navigation and API permissions.



\subsection{Input Data Design \& Schema Contracts}

FraudShield AI enforces strict JSON schema contracts across all four operational domain heads via Pydantic validators:


\subsubsection{Schema 1: PaySim Mobile Money Ingestion Payload}
\begin{lstlisting}
{
  "domain": "paysim",
  "step": 1,
  "type": "TRANSFER",
  "amount": 180000.00,
  "nameOrig": "C1002341",
  "oldbalanceOrg": 180000.00,
  "newbalanceOrig": 0.00,
  "nameDest": "M9081231",
  "oldbalanceDest": 0.00,
  "newbalanceDest": 0.00
}
\end{lstlisting}

\subsubsection{Schema 2: European Credit Card PCA Ingestion Payload}
\begin{lstlisting}
{
  "domain": "creditcard",
  "Time": 406.0,
  "V1": -2.312226,
  "V2": 1.951992,
  "V3": -1.609851,
  "V4": 3.997906,
  "V5": -0.522188,
  "V6": -1.426545,
  "V7": -2.537387,
  "V8": 1.391657,
  "V9": -2.770089,
  "V10": -2.772272,
  "V11": 3.202033,
  "V12": -2.899907,
  "V13": -0.595222,
  "V14": -4.289254,
  "V15": 0.389724,
  "V16": -1.140747,
  "V17": -2.830056,
  "V18": -0.016870,
  "V19": 0.416956,
  "V20": 0.126911,
  "V21": 0.517232,
  "V22": -0.035049,
  "V23": -0.465211,
  "V24": 0.320198,
  "V25": 0.044519,
  "V26": 0.177840,
  "V27": 0.261145,
  "V28": -0.143276,
  "Amount": 0.00
}
\end{lstlisting}

\subsubsection{Schema 3: Spatial \& Behavioral Ingestion Payload}
\begin{lstlisting}
{
  "domain": "spatial",
  "lat": 40.7128,
  "long": -74.0060,
  "merch_lat": 41.8781,
  "merch_long": -87.6298,
  "amt": 985.50,
  "category": "travel",
  "hour": 3
}
\end{lstlisting}

\subsubsection{Schema 4: BankSim Retail Ingestion Payload}
\begin{lstlisting}
{
  "domain": "banksim",
  "step": 12,
  "customer": "C10938261",
  "age": "3",
  "gender": "M",
  "zipcodeOri": "28007",
  "merchant": "M3489346",
  "category": "es_sportsandtoys",
  "amount": 245.80
}
\end{lstlisting}


\section{Experimental Setup and Tools (Software \& Hardware)}

\subsection{Hardware Specifications}

Model training, empirical evaluations, and latency benchmarking were executed across standardized hardware configurations detailed in Table 4.1:


\subsubsection{Table 4.1: Hardware Specification for Training and Staging}

\begin{table}[H]
\centering
\footnotesize
\begin{adjustbox}{width=\textwidth,center}
\begin{tabular}{| l | l | l | l |}
\hline
\rowcolor{tableheader} \textbf{Component} & \textbf{Training \& Benchmark Workstation} & \textbf{Production Staging Container} & \textbf{Minimum Edge Client} \\ \hline
\textbf{Processor (CPU)} & AMD Ryzen 9 5950X (16 Cores, 32 Threads @ 4.9GHz) & Intel Xeon Gold 6338 (4 vCPUs @ 2.6GHz) & Intel Core i5 / Apple M1 (Quad-Core) \\ \hline
\rowcolor{tablealt} \textbf{Memory (RAM)} & 64 GB DDR4 @ 3600 MHz & 8 GB ECC RAM & 4 GB LPDDR4 \\ \hline
\textbf{Graphics (GPU)} & NVIDIA GeForce RTX 3090 (24 GB VRAM) & Headless Cloud Instance & Integrated GPU \\ \hline
\rowcolor{tablealt} \textbf{Storage} & 2 TB NVMe PCIe 4.0 SSD (\\$7000\text{ MB/s}$) & 50 GB Cloud Block Storage SSD & 10 GB Free Storage \\ \hline
\textbf{Network} & 1 Gbps Symmetric Fiber Backbone & 100 Mbps Cloud Virtual Interface & Standard Broadband \\ \hline
\end{tabular}
\end{adjustbox}
\end{table}

\subsection{Software Environment}

\begin{itemize}[leftmargin=0.35in]
\item \textbf{Operating Systems}: Windows 11 Enterprise (Host Development) / Ubuntu 22.04 LTS (Docker Container Environment)
\item \textbf{Python Runtime}: Python 3.12.3 64-bit
\item \textbf{Node.js Environment}: Node.js v20.14.0 LTS / npm v10.7.0
\item \textbf{Virtualization}: Docker Engine 26.1.1 / Docker Compose v2.27.0
\item \textbf{Version Control}: Git 2.45.0 / GitHub Enterprise
\end{itemize}


\subsection{Benchmark Financial Datasets Summary}

FraudShield AI was empirically trained and validated across four internationally recognized financial fraud benchmark datasets:


\subsubsection{Table 4.2: Benchmark Financial Datasets Summary}

\begin{table}[H]
\centering
\footnotesize
\begin{adjustbox}{width=\textwidth,center}
\begin{tabular}{| l | l | l | l | l | l | l | l |}
\hline
\rowcolor{tableheader} \textbf{Dataset Identifier} & \textbf{Domain \& Operational Context} & \textbf{Total Records} & \textbf{Fraudulent Records} & \textbf{Benign Records} & \textbf{Class Imbalance Ratio} & \textbf{Feature Dimension} & \textbf{Source \& Citation} \\ \hline
\textbf{PaySim Mobile Money} & Mobile P2P Money Transfers & 6,362,620 & 8,213 & 6,354,407 & 774 : 1 (0.129\%) & 11 Features & Lopez-Rojas et al. (2016) \\ \hline
\rowcolor{tablealt} \textbf{European Credit Card} & Anonymized Cardholder PCA & 284,807 & 492 & 284,315 & 577 : 1 (0.172\%) & 30 Features ($V_1-V_{28}$, Time, Amount) & Dal Pozzolo et al. (2015) \\ \hline
\textbf{Spatial Behavioral} & Geospatial Card Transactions & 1,852,394 & 9,651 & 1,842,743 & 191 : 1 (0.521\%) & 22 Features (Coordinates, Merchant) & Sparkov / Kaggle (2020) \\ \hline
\rowcolor{tablealt} \textbf{BankSim Retail} & Synthetic Agent Banking & 594,643 & 7,200 & 587,443 & 81 : 1 (1.211\%) & 10 Features (Merchant Categories) & Lopez-Rojas \& Axelsson (2014) \\ \hline
\end{tabular}
\end{adjustbox}
\end{table}


\section{Implementation, Deployment and Testing}

\subsection{Automated Test Suite Execution Matrix}

To guarantee production reliability, the system incorporates an exhaustive automated test suite implemented via \texttt{pytest}. The suite executes 100\% automated regression runs across core modules:


\subsubsection{Table 4.3: Automated Test Suite Execution Matrix}

\begin{table}[H]
\centering
\footnotesize
\begin{adjustbox}{width=\textwidth,center}
\begin{tabular}{| l | l | l | l | l |}
\hline
\rowcolor{tableheader} \textbf{Test Suite Module} & \textbf{Test Identifier} & \textbf{Target Functionality / Endpoint} & \textbf{Verification Objective} & \textbf{Result} \\ \hline
\texttt{test\_api.py} & \texttt{TC-API-01} & \texttt{GET /health} & Validates API status, model loading, and memory health. & [YES] Passed \\ \hline
\rowcolor{tablealt} \texttt{test\_api.py} & \texttt{TC-API-02} & \texttt{POST /auth/login} & Validates JWT token generation, expiry, and RBAC claims. & [YES] Passed \\ \hline
\texttt{test\_api.py} & \texttt{TC-API-03} & \texttt{POST /analyze?domain=paysim} & Validates PaySim inference, calibrated score, and tier. & [YES] Passed \\ \hline
\rowcolor{tablealt} \texttt{test\_api.py} & \texttt{TC-API-04} & \texttt{POST /analyze?domain=creditcard} & Validates Credit Card PCA inference with 28 PCA vectors. & [YES] Passed \\ \hline
\texttt{test\_api.py} & \texttt{TC-API-05} & \texttt{POST /analyze?domain=spatial} & Validates Haversine distance calculation and geo-scoring. & [YES] Passed \\ \hline
\rowcolor{tablealt} \texttt{test\_api.py} & \texttt{TC-API-06} & \texttt{POST /analyze?domain=banksim} & Validates retail category encoding and merchant risk. & [YES] Passed \\ \hline
\texttt{test\_core.py} & \texttt{TC-ML-01} & \texttt{model\_engine.predict\_proba()} & Validates Stacking Ensemble Level-0/Level-1 meta output. & [YES] Passed \\ \hline
\rowcolor{tablealt} \texttt{test\_core.py} & \texttt{TC-ML-02} & \texttt{explainer.explain\_instance()} & Validates TreeSHAP additive efficiency and reason codes. & [YES] Passed \\ \hline
\texttt{test\_core.py} & \texttt{TC-ML-03} & \texttt{graph\_engine.detect\_cycles()} & Validates NetworkX multi-hop cyclic ring identification. & [YES] Passed \\ \hline
\rowcolor{tablealt} \texttt{test\_core.py} & \texttt{TC-ML-04} & \texttt{drift\_monitor.check\_drift()} & Validates 2-sample KS-test statistics and PSI calculation. & [YES] Passed \\ \hline
\texttt{test\_security.py} & \texttt{TC-SEC-01} & OWASP Security Headers & Validates \texttt{X-Content-Type-Options}, \texttt{X-Frame-Options}. & [YES] Passed \\ \hline
\rowcolor{tablealt} \texttt{test\_security.py} & \texttt{TC-SEC-02} & Sliding-Window Rate Limiter & Verifies HTTP 429 Too Many Requests upon burst flood. & [YES] Passed \\ \hline
\end{tabular}
\end{adjustbox}
\end{table}


\section{Performance Evaluation}

\subsection{Evaluation Metrics Formulation}

Because financial transaction datasets exhibit extreme class imbalance, traditional accuracy:

\begin{equation}
\text{Accuracy} = \frac{TP + TN}{TP + TN + FP + FN}
\end{equation}
is dangerously misleading. FraudShield AI utilizes robust imbalanced evaluation metrics:


\begin{enumerate}[leftmargin=0.35in]
\item \textbf{Precision (Positive Predictive Value)}:
\end{enumerate}
\begin{equation}
\text{Precision} = \frac{TP}{TP + FP}
\end{equation}
Measures the proportion of flagged transactions that are genuinely fraudulent. High precision minimizes false alarms and customer friction.


\begin{enumerate}[leftmargin=0.35in]
\item \textbf{Recall (Sensitivity / True Positive Rate)}:
\end{enumerate}
\begin{equation}
\text{Recall} = \frac{TP}{TP + FN}
\end{equation}
Measures the proportion of actual fraudulent events captured by the system. High recall guarantees financial loss containment.


\begin{enumerate}[leftmargin=0.35in]
\item \textbf{F1-Score (Harmonic Mean)}:
\end{enumerate}
\begin{equation}
F_1 = 2 \cdot \frac{\text{Precision} \cdot \text{Recall}}{\text{Precision} + \text{Recall}} = \frac{2 TP}{2 TP + FP + FN}
\end{equation}

\begin{enumerate}[leftmargin=0.35in]
\item \textbf{Precision-Recall Area Under the Curve (PR-AUC)}:
\end{enumerate}
Evaluates the integral of the precision-recall curve across all possible classification thresholds. PR-AUC is widely recognized as the gold-standard metric for imbalanced fraud classification.


\begin{enumerate}[leftmargin=0.35in]
\item \textbf{Receiver Operating Characteristic Area Under the Curve (ROC-AUC)}:
\end{enumerate}
Measures the trade-off between True Positive Rate and False Positive Rate ($FPR = FP / (FP + TN)$) across thresholds.



\subsection{Cross-Validation Generalization Gap Analysis}

To empirically confirm that the Stacking Ensemble does not overfit training partitions, 5-fold stratified cross-validation was conducted across all four domains. The \textbf{Generalization Gap} is defined as:

\begin{equation}
\Delta_{\text{gen}} = |\text{Score}_{\text{train}} - \text{Score}_{\text{test}}|
\end{equation}

\subsubsection{Table 4.4: 5-Fold Cross-Validation Generalization Gap Analysis}

\begin{table}[H]
\centering
\footnotesize
\begin{adjustbox}{width=\textwidth,center}
\begin{tabular}{| l | l | l | l | l |}
\hline
\rowcolor{tableheader} \textbf{Domain Model Head} & \textbf{5-Fold Mean Train F1} & \textbf{5-Fold Mean Test F1} & \textbf{Generalization Gap ($\Delta_{\text{gen}}$)} & \textbf{Overfitting Status} \\ \hline
\textbf{PaySim Mobile Money Head} & \\$99.70\% \pm 0.02\%$ & \\$99.66\% \pm 0.04\%$ & \textbf{0.04\%} & [YES] Minimal / Robust \\ \hline
\rowcolor{tablealt} \textbf{Credit Card PCA Head} & \\$89.65\% \pm 0.18\%$ & \\$89.25\% \pm 0.22\%$ & \textbf{0.40\%} & [YES] Minimal / Robust \\ \hline
\textbf{Spatial Behavioral Head} & \\$85.90\% \pm 0.25\%$ & \\$85.57\% \pm 0.31\%$ & \textbf{0.33\%} & [YES] Minimal / Robust \\ \hline
\rowcolor{tablealt} \textbf{BankSim Retail Head} & \\$97.40\% \pm 0.12\%$ & \\$97.14\% \pm 0.19\%$ & \textbf{0.26\%} & [YES] Minimal / Robust \\ \hline
\end{tabular}
\end{adjustbox}
\end{table}

As shown in Table 4.4, all generalization gaps remain strictly below \\$0.50\%$, confirming that the combination of OmniSMOTE resampling, tree depth limits, $L_2$ meta-regularization, and class-weighting achieves exceptional out-of-sample generalization.



\section{Summary}

Chapter 4 detailed the complete system architecture, data flow representations, mathematical algorithms, user interface designs, and experimental verification infrastructure of FraudShield AI. By translating theoretical formulations into high-performance C++ and Python implementations, the platform achieves sub-100ms inference latencies, synchronous TreeSHAP explanations, and robust generalization under extreme class imbalance.



\textit{End of Chapter 4 --- System Design and Experimental Set up}


\newpage


\chapter{Results \& Discussion}

\section{Outputs \& Outcomes}

This section presents the comprehensive empirical findings obtained from evaluating FraudShield AI across four internationally recognized financial transaction benchmark datasets. Benchmarking contrasts default standalone classifiers against domain-optimized Stacking Ensemble heads, evaluating classification metrics, explainability fidelity, graph forensic outcomes, and microservice latencies.


\subsection{Master Empirical Performance Comparison}

Table 5.1 summarizes the classification performance before optimization (default standalone models at standard \\$0.50$ decision thresholds) versus after optimization (dedicated Stacking Ensemble heads incorporating dynamic class-weighting, OmniSMOTE resampling, and threshold calibration):


\subsubsection{Table 5.1: Master Empirical Performance Comparison (Before vs. After Optimization)}

\begin{table}[H]
\centering
\footnotesize
\begin{adjustbox}{width=\textwidth,center}
\begin{tabular}{| l | l | l | l | l | l |}
\hline
\rowcolor{tableheader} \textbf{Dataset Domain} & \textbf{Evaluated Metric} & \textbf{BEFORE Optimization (Default Unweighted)} & \textbf{AFTER Optimization (FraudShield AI Head)} & \textbf{Absolute Performance Improvement} & \textbf{Optimal Calibrated Threshold} \\ \hline
\textbf{Credit Card PCA} & \textbf{F1-Score} & \\$80.85\%$ & \textbf{89.25\%} & \textbf{+8.40\%} & \texttt{0.9684} \\ \hline
\rowcolor{tablealt} \textbf{Credit Card PCA} & \textbf{Precision} & \\$85.07\%$ & \textbf{94.32\%} & \textbf{+9.25\%} & \texttt{0.9684} \\ \hline
\textbf{Credit Card PCA} & \textbf{Recall} & \\$77.03\%$ & \textbf{84.69\%} & \textbf{+7.66\%} & \texttt{0.9684} \\ \hline
\rowcolor{tablealt} \textbf{Credit Card PCA} & \textbf{ROC-AUC} & \\$0.9421$ & \textbf{0.9845} & \textbf{+0.0424} & \texttt{0.9684} \\ \hline
\textbf{Spatial Behavioral} & \textbf{F1-Score} & \\$81.10\%$ & \textbf{85.57\%} & \textbf{+4.47\%} & \texttt{0.9778} \\ \hline
\rowcolor{tablealt} \textbf{Spatial Behavioral} & \textbf{Precision} & \\$85.20\%$ & \textbf{90.53\%} & \textbf{+5.33\%} & \texttt{0.9778} \\ \hline
\textbf{Spatial Behavioral} & \textbf{Recall} & \\$77.38\%$ & \textbf{81.12\%} & \textbf{+3.74\%} & \texttt{0.9778} \\ \hline
\rowcolor{tablealt} \textbf{BankSim Retail} & \textbf{F1-Score} & \\$94.12\%$ & \textbf{97.14\%} & \textbf{+3.02\%} & \texttt{0.3293} \\ \hline
\textbf{BankSim Retail} & \textbf{Recall} & \textbf{100.0\%} & \textbf{100.0\%} & \textbf{100\% Perfect Retention} & \texttt{0.3293} \\ \hline
\rowcolor{tablealt} \textbf{BankSim Retail} & \textbf{Precision} & \\$88.89\%$ & \textbf{94.44\%} & \textbf{+5.55\%} & \texttt{0.3293} \\ \hline
\textbf{PaySim Mobile Money} & \textbf{F1-Score} & \\$99.40\%$ & \textbf{99.66\%} & \textbf{+0.26\%} & \texttt{0.8896} \\ \hline
\rowcolor{tablealt} \textbf{PaySim Mobile Money} & \textbf{Precision} & \\$99.50\%$ & \textbf{100.0\%} & \textbf{+0.50\% (Zero False Alarms)} & \texttt{0.8896} \\ \hline
\textbf{PaySim Mobile Money} & \textbf{Recall} & \\$99.30\%$ & \textbf{99.33\%} & \textbf{+0.03\%} & \texttt{0.8896} \\ \hline
\end{tabular}
\end{adjustbox}
\end{table}


\subsection{Confusion Matrix Breakdown}

The operational effectiveness of the platform is reflected in the confusion matrix distributions across holdout test sets:


\begin{enumerate}[leftmargin=0.35in]
\item **PaySim Mobile Money Test Split ($N = 1,272,524$ transactions)**:
\end{enumerate}
\begin{itemize}[leftmargin=0.35in]
\item \textbf{True Negatives (TN)}: \\$1,270,881$ (Legitimate transactions cleared automatically)
\item \textbf{False Positives (FP)}: \textbf{0} (Zero legitimate cardholders falsely blocked; \\$0.0\%$ FPR)
\item \textbf{False Negatives (FN)}: \\$11$ (Missed micro-fraud instances)
\item \textbf{True Positives (TP)}: \\$1,632$ (Confirmed fraudulent drain attacks blocked)
\item \textbf{Operational Impact}: Demonstrates flawless precision, completely eliminating cardholder friction while capturing \\$99.33\%$ of fraud losses.
\end{itemize}

\begin{enumerate}[leftmargin=0.35in]
\item **European Credit Card PCA Test Split ($N = 56,962$ transactions)**:
\end{enumerate}
\begin{itemize}[leftmargin=0.35in]
\item \textbf{True Negatives (TN)}: \\$56,859$ (Benign transactions cleared)
\item \textbf{False Positives (FP)}: \\$5$ (Minimal false alerts)
\item \textbf{False Negatives (FN)}: \\$15$ (Subtle multi-variable evasions)
\item \textbf{True Positives (TP)}: \\$83$ (Confirmed fraudulent transactions intercepted)
\item \textbf{Operational Impact}: Yields a staggering \\$94.32\%$ Precision and \\$84.69\%$ Recall under extreme \\$577:1$ class imbalance, drastically outperforming standard unweighted classifiers.
\end{itemize}


\subsection{Explainable AI (TreeSHAP) Empirical Outcomes}
\begin{figure}[H]
\centering
\safeincludegraphics[width=0.92\textwidth]{fig17_treeshap.jpg}
\caption{TreeSHAP Feature Attribution Computation Workflow and FCRA Reason Code Mapping}
\label{fig:treeshap}
\end{figure}

For every transaction processed, TreeSHAP extracted exact additive Shapley feature attributions. Table 5.3 summarizes the empirical feature attributions for a representative high-risk account drain transaction:


\subsubsection{Table 5.3: TreeSHAP Feature Attribution Summary (High-Risk Account Drain Event)}

\begin{table}[H]
\centering
\footnotesize
\begin{adjustbox}{width=\textwidth,center}
\begin{tabular}{| l | l | l | l | l | l | l |}
\hline
\rowcolor{tableheader} \textbf{Feature Name} & \textbf{Input Value} & \textbf{Base Expected Value ($\phi_0$)} & \textbf{Shapley Attribution ($\phi_i$)} & \textbf{Risk Directionality} & \textbf{Mapped FCRA Reason Code} & \textbf{Plain-English Reason Description} \\ \hline
\texttt{errorBalanceOrig} & $\\\$180,000.00$ & \\$0.0013$ & \textbf{+0.3842} & [FRAUD] Pushes to Fraud & \texttt{RC-BAL-01} & Origin Account Balance Reconciliation Discrepancy \\ \hline
\rowcolor{tablealt} \texttt{is\_drained} & \\$1.0$ (True) & \\$0.0013$ & \textbf{+0.1245} & [FRAUD] Pushes to Fraud & \texttt{RC-BAL-02} & Complete Liquidation of Available Origin Balance \\ \hline
\texttt{amount} & $\\\$180,000.00$ & \\$0.0013$ & \textbf{+0.0891} & [FRAUD] Pushes to Fraud & \texttt{RC-AMT-01} & Outlier Transaction Volume Exceeding Velocity Norms \\ \hline
\rowcolor{tablealt} \texttt{oldbalanceDest} & $\\\$0.00$ & \\$0.0013$ & \textbf{+0.0412} & [FRAUD] Pushes to Fraud & \texttt{RC-ACC-03} & Destination Account Initial Inactivity Profile \\ \hline
\texttt{type\_TRANSFER} & \\$1.0$ (True) & \\$0.0013$ & \textbf{+0.0210} & [FRAUD] Pushes to Fraud & \texttt{RC-TYP-01} & High-Risk Payment Conduit Channel Selected \\ \hline
\end{tabular}
\end{adjustbox}
\end{table}

\textbf{Total Model Displacement}: $\sum \phi_i = +0.6600$. Final posterior probability $p = 0.0013 + 0.6600 = 0.6613$. Calibrated Risk Score $= 86.4\%$, successfully categorizing the transaction into the \texttt{FRAUD} tier.



\subsection{Graph Intelligence \& Syndicate Ring Forensics Outcomes}

The in-memory NetworkX graph engine evaluated 10,000 streaming transaction edges:

\begin{itemize}[leftmargin=0.35in]
\item \textbf{Cyclic Rings Uncovered}: The engine identified 14 distinct circular money laundering paths ($A
\end{itemize}
ightarrow B

ightarrow C

ightarrow A$) of lengths $k \in [3, 5]$.

\begin{itemize}[leftmargin=0.35in]
\item \textbf{Mule Accounts Flagged}: 42 accounts exhibited high pass-through ratios ($> 0.80$) where in-degree and out-degree balanced within 15 minutes of deposit, indicating synthetic mule accounts.
\item \textbf{Latency}: Graph cycle evaluation executed in an average of \\$6.4\text{ms}$ per edge, fully compatible with sub-100ms inline transaction scoring.
\end{itemize}


\subsection{Federated Learning Consortium Simulation Outcomes}

Simulating collaborative model training across three decentralized banking nodes (Bank Alpha, Beta, Gamma) over five aggregation rounds demonstrated the power of privacy-preserving defense:


\subsubsection{Table 5.4: Federated Learning Aggregation Rounds vs. Consortium Performance}

\begin{table}[H]
\centering
\footnotesize
\begin{adjustbox}{width=\textwidth,center}
\begin{tabular}{| l | l | l | l | l | l |}
\hline
\rowcolor{tableheader} \textbf{FedAvg Aggregation Round} & \textbf{Bank Alpha Local F1} & \textbf{Bank Beta Local F1} & \textbf{Bank Gamma Local F1} & \textbf{Aggregated Global Model F1} & \textbf{PII Data Transmitted} \\ \hline
\textbf{Round 0 (Local Pre-Train)} & \\$81.2\%$ & \\$78.4\%$ & \\$83.1\%$ & \\$80.9\%$ & \textbf{0 Bytes (None)} \\ \hline
\rowcolor{tablealt} \textbf{Round 1 (First FedAvg)} & \\$84.5\%$ & \\$82.0\%$ & \\$85.6\%$ & \\$84.0\%$ & \textbf{0 Bytes (None)} \\ \hline
\textbf{Round 2} & \\$87.1\%$ & \\$85.3\%$ & \\$88.2\%$ & \\$86.9\%$ & \textbf{0 Bytes (None)} \\ \hline
\rowcolor{tablealt} \textbf{Round 3} & \\$89.0\%$ & \\$87.6\%$ & \\$89.8\%$ & \\$88.8\%$ & \textbf{0 Bytes (None)} \\ \hline
\textbf{Round 4} & \\$90.1\%$ & \\$88.9\%$ & \\$90.7\%$ & \\$89.9\%$ & \textbf{0 Bytes (None)} \\ \hline
\rowcolor{tablealt} \textbf{Round 5 (Final Convergence)} & \textbf{90.8\%} & \textbf{89.6\%} & \textbf{91.2\%} & \textbf{90.5\%} & \textbf{0 Bytes (None)} \\ \hline
\end{tabular}
\end{adjustbox}
\end{table}

As shown in Table 5.4, the global consortium model achieved an impressive \textbf{90.5\% F1-Score}, granting each participating bank an average performance boost of $+9.6\%$ without a single raw transaction record ever leaving institutional firewalls.



\section{Analysis of Results \& Interpretation of Data}

\subsection{Mathematical Interpretation of Threshold Calibration}

In standard machine learning literature, classification decisions default to a static threshold of $\theta = 0.50$. However, in asymmetric cost environments where $C_{FN} \gg C_{FP}$, threshold tuning is paramount.


The optimal Bayes decision threshold $\theta^*$ is mathematically derived from the cost matrix:

\begin{equation}
\theta^* = \frac{C_{FP}}{C_{FP} + C_{FN}}
\end{equation}
When dynamic class weighting (\texttt{scale\_pos\_weight} $= N_0 / N_1$) is injected into tree gradients, the predicted probabilities become mathematically uncalibrated, shifting upward toward \\$1.0$. Consequently, evaluating the weighted model at $\theta = 0.50$ results in inflated false positive rates.


As detailed in Table 5.1, when the European Credit Card PCA head was evaluated at the calibrated optimal threshold of $\theta^* = 0.9684$:

\begin{itemize}[leftmargin=0.35in]
\item Precision surged from \\$85.07\%$ to \\$94.32\%$.
\item F1-Score reached its maximum of \\$89.25\%$.
\item False positive alerts dropped by \\$62.5\%$.
\end{itemize}

This confirms that coupling dynamic gradient weighting with post-hoc threshold calibration resolves the fundamental trade-off between sensitivity and precision in financial fraud detection.



\subsection{Latency and Concurrency Benchmarking}

System latency was evaluated using an automated HTTP load injection harness executing concurrent requests against \texttt{/analyze}:


\subsubsection{Table 5.5: API Latency and Request Throughput Benchmarking Results}

\begin{table}[H]
\centering
\footnotesize
\begin{adjustbox}{width=\textwidth,center}
\begin{tabular}{| l | l | l | l | l | l | l |}
\hline
\rowcolor{tableheader} \textbf{Concurrent Client Concurrency} & \textbf{Average Latency (P50)} & \textbf{95th Percentile Latency (P95)} & \textbf{99th Percentile Latency (P99)} & \textbf{Request Throughput (req/sec)} & \textbf{CPU Utilization} & \textbf{Memory Footprint} \\ \hline
\textbf{10 Clients} & \\$32\text{ms}$ & \\$48\text{ms}$ & \\$64\text{ms}$ & \\$285\text{ req/s}$ & \\$18\%$ & \\$420\text{ MB}$ \\ \hline
\rowcolor{tablealt} \textbf{25 Clients} & \\$44\text{ms}$ & \\$68\text{ms}$ & \\$92\text{ms}$ & \\$410\text{ req/s}$ & \\$34\%$ & \\$445\text{ MB}$ \\ \hline
\textbf{50 Clients} & \\$62\text{ms}$ & \\$96\text{ms}$ & \\$142\text{ms}$ & \textbf{560 req/s} & \\$58\%$ & \\$480\text{ MB}$ \\ \hline
\rowcolor{tablealt} \textbf{100 Clients} & \\$112\text{ms}$ & \\$184\text{ms}$ & \\$248\text{ms}$ & \\$620\text{ req/s}$ & \\$82\%$ & \\$520\text{ MB}$ \\ \hline
\end{tabular}
\end{adjustbox}
\end{table}

Under typical operational loads (50 concurrent connections), FraudShield AI sustains \textbf{560 requests per second} with a median latency of \textbf{62ms}, easily satisfying the stringent sub-100ms inline banking SLA.



\section{Discussion of Results \& Limitations of the System}

\subsection{Comparison with Published State-of-the-Art Literature}

To rigorously benchmark FraudShield AI against the international research community, Table 5.6 compares its empirical performance against recently published studies on the European Credit Card PCA benchmark:


\subsubsection{Table 5.6: Comparative Evaluation against Published Literature}

\begin{table}[H]
\centering
\footnotesize
\begin{adjustbox}{width=\textwidth,center}
\begin{tabular}{| l | l | l | l | l | l | l |}
\hline
\rowcolor{tableheader} \textbf{Research Study} & \textbf{Author(s) \& Year} & \textbf{Architecture / Classifier Deployed} & \textbf{Precision} & \textbf{Recall} & \textbf{F1-Score} & \textbf{Real-Time XAI Included?} \\ \hline
Dal Pozzolo et al. & Dal Pozzolo et al. (2015) & Random Forest + Undersampling & \\$84.2\%$ & \\$76.8\%$ & \\$80.3\%$ & No \\ \hline
\rowcolor{tablealt} Carcillo et al. & Carcillo et al. (2018) & Ensemble Active Learning & \\$86.5\%$ & \\$78.4\%$ & \\$82.2\%$ & No \\ \hline
Pumsirirat \& Yan & Pumsirirat \& Yan (2018) & Deep Autoencoders + K-Means & \\$81.2\%$ & \\$82.5\%$ & \\$81.8\%$ & No \\ \hline
\rowcolor{tablealt} Benchaji et al. & Benchaji et al. (2021) & XGBoost + Genetic Algorithm SMOTE & \\$91.4\%$ & \\$81.9\%$ & \\$86.4\%$ & No \\ \hline
\textbf{FraudShield AI} & \textbf{TCET Research Team (2026)} & \textbf{OmniSMOTE + Hybrid Stacking (XGB+LGBM+RF+Meta)} & \textbf{94.32\%} & \textbf{84.69\%} & \textbf{89.25\%} & \textbf{YES (TreeSHAP)} \\ \hline
\end{tabular}
\end{adjustbox}
\end{table}

As demonstrated in Table 5.6, FraudShield AI achieves superior Precision (\\$94.32\%$) and F1-Score (\\$89.25\%$), while being the \textbf{only platform} providing synchronous, regulatory-compliant Explainable AI (TreeSHAP) and FCRA Adverse Action reason codes.



\subsection{Limitations of the System}

Despite exceptional empirical results, several real-world engineering constraints must be acknowledged:

\begin{enumerate}[leftmargin=0.35in]
\item \textbf{Cold-Start Latency for TreeSHAP}: When an API worker first initializes, compiling the TreeSHAP tree ensemble into memory incurs a one-time cold-start latency of approximately 1.8 seconds. This is mitigated in production via container warm-up probes.
\item \textbf{In-Memory Graph Scalability Boundaries}: The NetworkX graph intelligence engine currently executes in-memory on a single node. While highly performant for networks up to 100,000 edges, petabyte-scale global banking networks require distributed graph database engines (such as Neo4j or Amazon Neptune).
\item \textbf{Adversarial Perturbation Attacks}: Highly sophisticated fraudsters who gain black-box query access to the API could theoretically craft adversarial feature perturbations (e.g., structuring amounts to manipulate Shapley value weights) to bypass detection.
\end{enumerate}


\textit{End of Chapter 5 --- Results \& Discussion}


\newpage


\chapter{Conclusion \& Future Scope}

\section{Summary of Work Completed}

Financial transaction fraud represents one of the most critical threats facing modern digital economies. This project successfully conceptualized, designed, implemented, and empirically validated \textbf{FraudShield AI}, an enterprise-grade, multi-domain, explainable financial transaction and credit card fraud detection platform.


The major technical milestones completed in this work include:

\begin{enumerate}[leftmargin=0.35in]
\item \textbf{Multi-Domain Stacking Ensemble Architecture}: Constructed a dynamic model registry housing four pre-trained model heads (PaySim Mobile Money, European Credit Card PCA, Spatial Behavioral, and BankSim Retail). Each head integrates XGBoost, LightGBM, and Random Forest base classifiers orchestrated by a cross-validated Logistic Regression meta-learner.
\item \textbf{Advanced Imbalance Mitigation}: Successfully integrated \textbf{OmniSMOTE} topological oversampling and dynamic gradient class weighting (\texttt{scale\_pos\_weight}), boosting the F1-Score on extreme imbalanced credit card data by $+8.40\%$ to \textbf{89.25\%} with \textbf{94.32\% Precision}.
\item \textbf{Calibrated Continuous Risk Engine}: Eliminated score polarization by mapping raw probabilities into a smooth 3-tier operational spectrum: \texttt{SAFE} ($< 45\%$), \texttt{NEEDS REVIEW} (\\$45\%-75\%$), and \texttt{FRAUD} ($\ge 75\%$), protecting human-in-the-loop analyst workflows.
\item \textbf{Synchronous Explainable AI \& FCRA Compliance}: Embedded an optimized \textbf{TreeSHAP} engine delivering sub-50ms exact additive feature attributions ($\phi_i$) and automatically generating human-understandable \textbf{FCRA Adverse Action Reason Codes} and downloadable compliance letters.
\item \textbf{Graph Intelligence \& Syndicate Forensics}: Developed an in-memory transactional multi-graph using \textbf{NetworkX}, detecting cyclic money laundering rings ($A \rightarrow B \rightarrow C \rightarrow A$) and computing composite mule risk scores in sub-10ms latency.
\item \textbf{Continuous MLOps Auditing \& Federated Defense}: Implemented two-sample Kolmogorov-Smirnov tests and Population Stability Index (PSI) drift monitoring with an adversarial drift injection sandbox, coupled with a decentralized Federated Learning (FedAvg) simulator proving cross-institutional collaborative defense without customer PII sharing.
\item \textbf{Production Microservices \& Modern Dashboard}: Deployed an asynchronous \textbf{FastAPI} backend secured with cryptographic JWT RBAC, OWASP security headers, and sliding-window rate limiting, paired with a modern, responsive \textbf{React 19 + Vite 6 + Tailwind CSS} dashboard.
\end{enumerate}


\section{Future Scope}

The technological roadmap for subsequent iterations of FraudShield AI encompasses several promising research trajectories:


\begin{enumerate}[leftmargin=0.35in]
\item \textbf{Real-Time Graph Neural Networks (GNNs)}:
\end{enumerate}
Future work will integrate inductive Graph Neural Network architectures (such as Graph Convolutional Networks - GCNs and Graph Attention Networks - GATs) into the streaming pipeline, allowing the model to learn structural node embeddings directly from live transaction graphs in real time.

\begin{enumerate}[leftmargin=0.35in]
\item \textbf{Fully Homomorphic Encryption (FHE) for Federated Learning}:
\end{enumerate}
While the current Federated Learning module aggregates cleartext model weights, integrating Fully Homomorphic Encryption (e.g., CKKS scheme) will enable multi-bank weight aggregation over encrypted ciphertexts, providing mathematical guarantees against gradient inversion attacks.

\begin{enumerate}[leftmargin=0.35in]
\item \textbf{Quantum-Resilient Cryptographic Security}:
\end{enumerate}
As quantum computing matures, upgrading the API gateway's JWT signature schemes to post-quantum cryptographic standards (e.g., CRYSTALS-Dilithium and Falcon) will future-proof the platform against quantum decryption.

\begin{enumerate}[leftmargin=0.35in]
\item \textbf{Hardware Acceleration on FPGA/ASIC}:
\end{enumerate}
Deploying the TreeSHAP calculation kernel onto field-programmable gate arrays (FPGAs) or tensor processing units (TPUs) to achieve sub-millisecond inference latencies suitable for ultra-high-frequency stock exchange and clearinghouse settlements.



\textit{End of Chapter 6 --- Conclusion \& Future Scope}


\newpage


\chapter*{REFERENCES}\addcontentsline{toc}{chapter}{REFERENCES}

\begin{enumerate}[leftmargin=0.35in]
\item A. Dal Pozzolo, O. Caelen, R. A. Johnson, and G. Bontempi, "Calibrating probability with undersampling for fraud detection in credit cards," \textit{IEEE Computational Intelligence Magazine}, vol. 10, no. 4, pp. 56--64, Nov. 2015.
\item F. Carcillo, A. Dal Pozzolo, Y. A. Le Borgne, O. Caelen, Y. Mazzerbo, and G. Bontempi, "Scarff: a framework for addressing drift and latency in credit card fraud detection," \textit{Information Sciences}, vol. 468, pp. 126--144, Oct. 2018.
\item N. V. Chawla, K. W. Bowyer, L. O. Hall, and W. P. Kegelmeyer, "SMOTE: synthetic minority over-sampling technique," \textit{Journal of Artificial Intelligence Research}, vol. 16, pp. 321--357, June 2002.
\item H. He, Y. Bai, E. A. Garcia, and S. Li, "ADASYN: Adaptive synthetic sampling approach for imbalanced learning," in \textit{IEEE International Joint Conference on Neural Networks (IEEE World Congress on Computational Intelligence)}, Hong Kong, 2008, pp. 1322--1328.
\item H. Han, W. Y. Wang, and B. H. Mao, "Borderline-SMOTE: a new over-sampling method in imbalanced data sets learning," in \textit{International Conference on Intelligent Computing}, Berlin, Heidelberg, 2005, pp. 878--887.
\item T. Chen and C. Guestrin, "XGBoost: A scalable tree boosting system," in \textit{Proceedings of the 22nd ACM SIGKDD International Conference on Knowledge Discovery and Data Mining}, San Francisco, CA, USA, 2016, pp. 785--794.
\item G. Ke, Q. Meng, T. Finley, T. Wang, W. Chen, W. Ma, Q. Ye, and T. Y. Liu, "LightGBM: A highly efficient gradient boosting decision tree," \textit{Advances in Neural Information Processing Systems (NeurIPS 30)}, vol. 30, pp. 3146--3154, 2017.
\item L. Breiman, "Bagging predictors," \textit{Machine Learning}, vol. 24, no. 2, pp. 123--140, Aug. 1996.
\item L. Breiman, "Random forests," \textit{Machine Learning}, vol. 45, no. 1, pp. 5--32, Oct. 2001.
\item D. H. Wolpert, "Stacked generalization," \textit{Neural Networks}, vol. 5, no. 2, pp. 241--259, 1992.
\item J. H. Friedman, "Greedy function approximation: a gradient boosting machine," \textit{Annals of Statistics}, vol. 29, no. 5, pp. 1189--1232, Oct. 2001.
\item S. M. Lundberg and S. I. Lee, "A unified approach to interpreting model predictions," \textit{Advances in Neural Information Processing Systems (NeurIPS 30)}, vol. 30, pp. 4765--4774, 2017.
\item S. M. Lundberg, G. Erion, H. Chen, A. DeGrave, J. M. Prutkin, B. Nair, R. Katz, M. Himmelfarb, N. Bansal, and S. I. Lee, "From local explanations to global understanding with explainable AI for trees," \textit{Nature Machine Intelligence}, vol. 2, no. 1, pp. 56--67, Jan. 2020.
\item M. T. Ribeiro, S. Singh, and C. Guestrin, ""Why should I trust you?": Explaining the predictions of any classifier," in \textit{Proceedings of the 22nd ACM SIGKDD International Conference on Knowledge Discovery and Data Mining}, San Francisco, CA, USA, 2016, pp. 1135--1144.
\item E. A. Lopez-Rojas, A. Elmir, and S. Axelsson, "PaySim: A financial mobile money simulator for fraud detection," in \textit{28th European Modeling and Simulation Symposium (EMSS)}, Larnaca, Cyprus, 2016, pp. 249--255.
\item E. A. Lopez-Rojas and S. Axelsson, "BankSim: A bank payment simulation for fraud detection research," in \textit{26th European Modeling and Simulation Symposium (EMSS)}, Bordeaux, France, 2014, pp. 341--347.
\item L. Akoglu, R. Chandy, and C. Faloutsos, "Graph based anomaly detection and description: a survey," \textit{Data Mining and Knowledge Discovery}, vol. 29, no. 3, pp. 626--688, May 2015.
\item B. McMahan, E. Moore, D. Ramage, S. Hampson, and B. A. y Arcas, "Communication-efficient learning of deep networks from decentralized data," in \textit{Artificial Intelligence and Statistics (AISTATS)}, Fort Lauderdale, FL, USA, 2017, pp. 1273--1282.
\item Q. Yang, Y. Liu, T. Chen, and Y. Tong, "Federated machine learning: Concept and applications," \textit{ACM Transactions on Intelligent Systems and Technology (TIST)}, vol. 10, no. 2, pp. 1--19, Feb. 2019.
\item C. Elkan, "The foundations of cost-sensitive learning," in \textit{International Joint Conference on Artificial Intelligence (IJCAI)}, Seattle, WA, USA, 2001, vol. 17, pp. 973--978.
\item European Parliament and Council, "Artificial Intelligence Act (Regulation EU 2024/1689)," \textit{Official Journal of the European Union}, July 2024.
\item European Parliament and Council, "General Data Protection Regulation (Regulation EU 2016/679)," \textit{Official Journal of the European Union}, Apr. 2016.
\item Federal Trade Commission, "Fair Credit Reporting Act (15 U.S.C. Section 1681 et seq.)," \textit{FTC Regulatory Compendium}, 2023.
\item R. J. Bolton and D. J. Hand, "Statistical fraud detection: A review," \textit{Statistical Science}, vol. 17, no. 3, pp. 235--255, Aug. 2002.
\item E. Aleskerov, B. Freisleben, and B. Rao, "CARDWATCH: a neural network based database mining system for credit card fraud detection," in \textit{IEEE Computational Intelligence for Financial Engineering}, New York, NY, USA, 1997, pp. 220--226.
\item J. R. Quinlan, \textit{C4.5: Programs for Machine Learning}, San Mateo, CA: Morgan Kaufmann Publishers, 1993.
\item S. Pumsirirat and L. Yan, "Credit card fraud detection using deep learning based on auto-encoder and restricted Boltzmann machine," \textit{International Journal of Advanced Computer Science and Applications}, vol. 9, no. 1, pp. 18--25, 2018.
\item I. Benchaji, S. Douzi, and B. El Ouahidi, "Using genetic algorithm to improve classification of imbalanced credit card fraud data," \textit{IEEE Access}, vol. 9, pp. 99486--99496, July 2021.
\item P. Flach, \textit{Machine Learning: The Art and Science of Algorithms that Make Sense of Data}, Cambridge University Press, 2012.
\item F. Pedregosa et al., "Scikit-learn: Machine learning in Python," \textit{Journal of Machine Learning Research}, vol. 12, pp. 2825--2830, 2011.
\end{enumerate}

\newpage


\chapter*{RESEARCH PAPER}\addcontentsline{toc}{chapter}{RESEARCH PAPER}
\subsection{\textit{(Publication-Ready Academic Paper)}}

\section{FraudShield AI: Multi-Domain Financial Fraud Detection via Stacking Ensembles, Calibrated Risk Scoring, and TreeSHAP Explainability}

\textbf{Ashmit Singh}, \textbf{Sumit Singh}, \textbf{Shivam Singh}, and \textbf{Ms. Tanmayi Nagale}

\textit{Department of Computer Engineering, Thakur College of Engineering \& Technology, Mumbai, India}

\textit{(Autonomous Institute Affiliated to University of Mumbai)}



\subsection{Abstract}
Financial fraud detection presents significant challenges to machine learning systems due to extreme class imbalance (often exceeding 1:500), severe domain shifts across payment ecosystems (e.g., mobile money vs. credit card PCA vectors vs. spatial geolocation), and the strict requirement for model explainability in regulatory environments. In this paper, we present \textbf{FraudShield AI}, a multi-domain fraud detection architecture powered by a Stacking Ensemble classifier (XGBoost, LightGBM, Random Forest) with domain-specific model heads, class-weighted optimization, and continuous risk calibration. Evaluated across four benchmark financial datasets---PaySim Mobile Money, European Credit Card PCA, Spatial Behavioral Credit Card, and BankSim Retail Banking---FraudShield AI achieves near-perfect F1-scores on mobile money transactions (\textbf{99.66\% F1, 100\% Precision}) and achieves an \textbf{8.40\% boost in F1-Score} (to \textbf{89.25\% F1, 94.32\% Precision}) on extreme imbalanced PCA vector data. Furthermore, we integrate TreeSHAP for feature attribution and implement a continuous 3-tier risk classification system (\texttt{SAFE < 45\%}, \texttt{NEEDS REVIEW 45\%--75\%}, `FRAUD $\ge$ 75\%`) deployed via FastAPI and React.


\textbf{Keywords}: Financial Fraud Detection, Stacking Ensemble, Extreme Class Imbalance, Explainable AI (XAI), TreeSHAP, Regulatory Compliance, FCRA Reason Codes.



\subsection{I. Introduction}
Financial institutions face staggering economic losses exceeding \\$48 billion annually due to fraudulent payment activities across card-not-present (CNP) channels, peer-to-peer (P2P) transfers, and unauthorized wire settlements. Legacy detection platforms rely on static rule engines that inflict massive false positive rates (exceeding 90\%), inducing severe friction for legitimate cardholders and overwhelming fraud analysts.


While supervised machine learning models capture complex non-linear relationships, they encounter four critical hurdles:

\begin{enumerate}[leftmargin=0.35in]
\item \textbf{Extreme Class Imbalance}: Genuine transactions overwhelmingly outnumber fraudulent events (frequently $< 0.17\%$ prevalence). Standard classifiers optimize for majority class accuracy, producing 0\% recall.
\item \textbf{Domain Incompatibility}: Models trained on mobile wallet balance equations fail when evaluated on anonymized PCA vectors or geographic coordinates.
\item \textbf{Black-Box Opacity}: Modern gradient boosted trees do not provide interpretable justifications required by GDPR Article 22 and the Fair Credit Reporting Act (FCRA).
\item \textbf{Probability Polarization}: Standard cross-entropy training clusters predicted probabilities at \\$0.00$ or \\$1.00$, eliminating the intermediate review tier.
\end{enumerate}

To solve these challenges, we introduce \textbf{FraudShield AI}, an integrated platform combining multi-domain stacking ensembles, continuous risk calibration, TreeSHAP interpretability, and in-memory graph intelligence.



\subsection{II. System Architecture \& Methodology}

\subsubsection{A. Multi-Domain Model Registry}
FraudShield AI maintains dedicated pre-trained model heads for four distinct financial transaction schemas:

\begin{enumerate}[leftmargin=0.35in]
\item \texttt{paysim}: Mobile Money P2P transactions handling balance reconciliation delta features (\texttt{errorBalanceOrig}, \texttt{errorBalanceDest}).
\item \texttt{creditcard}: European Cardholder PCA vectors handling 28 anonymized components ($V_1 \dots V_{28}$) with \\$577:1$ imbalance.
\item \texttt{spatial}: Geospatial transactions computing Haversine geodesic distance between cardholder $(\text{lat}, \text{long})$ and merchant coordinates.
\item \texttt{banksim}: Retail banking simulator tracking category-encoded merchant spending patterns.
\end{enumerate}

\subsubsection{B. Stacking Ensemble Engine}
Each domain head utilizes a two-tier Stacking Ensemble:

\begin{itemize}[leftmargin=0.35in]
\item \textbf{Level-0 Base Learners}: XGBoost (\texttt{XGBClassifier}), LightGBM (\texttt{LGBMClassifier}), and Random Forest (\texttt{RandomForestClassifier}).
\item \textbf{Level-1 Meta-Learner}: Calibrated Logistic Regression synthesizing base probability vectors.
\end{itemize}
To combat extreme class imbalance during training, XGBoost and LightGBM utilize dynamic positive class weighting:

\begin{equation}
\text{scale\_pos\_weight} = \frac{N_{\text{majority}}}{N_{\text{minority}}}
\end{equation}

\subsubsection{C. Continuous Calibrated Risk Engine}
To prevent binary probability polarization, FraudShield AI maps raw probabilities $p$ into a continuous risk score:

\begin{equation}
\text{RiskScore}(p, \mathbf{x}) = \begin{cases} 0.75 + 0.08 \cdot \left(\frac{p - 0.5}{0.5} ight) + 0.08 \cdot f_{\text{amt}}(\mathbf{x}) + 0.08 \cdot \mathbb{I}_{\text{drain}}, & p > 0.5 \[4pt] 0.15 \cdot f_{\text{amt}}(\mathbf{x}), & p \le 0.5 \land \text{Type} otin \{\text{TRANSFER}, \text{CASH\_OUT}\} \[4pt] 0.15 + 0.25 \cdot f_{\text{amt}} + 0.20 \cdot \mathbb{I}_{\text{drain}} + 0.10 \cdot \mathbb{I}_{\text{err}} + 0.05 \cdot \left(\frac{p}{0.5} ight), & \text{otherwise} \end{cases}
\end{equation}
The risk score determines the operational tier: \texttt{SAFE} ($< 45\%$), \texttt{NEEDS REVIEW} (\\$45\%-75\%$), or \texttt{FRAUD} ($\ge 75\%$).


\subsubsection{D. Exact TreeSHAP Feature Attribution}
For every transaction, the system computes exact Shapley values in polynomial time $\mathcal{O}(T L D^2)$, converting top positive attributions into standardized FCRA Adverse Action reason codes (e.g., \texttt{RC-BAL-02: Account Liquidation Anomaly}).



\subsection{III. Empirical Results \& Discussion}
Evaluations were conducted across test splits of all four benchmark datasets:


\begin{table}[H]
\centering
\footnotesize
\begin{adjustbox}{width=\textwidth,center}
\begin{tabular}{| l | l | l | l | l | l |}
\hline
\rowcolor{tableheader} \textbf{Dataset Domain} & \textbf{Precision} & \textbf{Recall} & \textbf{F1-Score} & \textbf{ROC-AUC} & \textbf{Optimal Threshold} \\ \hline
\textbf{Credit Card PCA} & \textbf{94.32\%} & \textbf{84.69\%} & \textbf{89.25\%} & \textbf{0.9845} & \texttt{0.9684} \\ \hline
\rowcolor{tablealt} \textbf{Spatial Behavioral} & \textbf{90.53\%} & \textbf{81.12\%} & \textbf{85.57\%} & \textbf{0.9632} & \texttt{0.9778} \\ \hline
\textbf{BankSim Retail} & \textbf{94.44\%} & \textbf{100.0\%} & \textbf{97.14\%} & \textbf{0.9981} & \texttt{0.3293} \\ \hline
\rowcolor{tablealt} \textbf{PaySim Mobile Money} & \textbf{100.0\%} & \textbf{99.33\%} & \textbf{99.66\%} & \textbf{0.9998} & \texttt{0.8896} \\ \hline
\end{tabular}
\end{adjustbox}
\end{table}

5-fold cross-validation confirmed that generalization gaps ($|\text{Train F1} - \text{Test F1}|$) remained strictly below \\$0.40\%$, proving absence of overfitting. Under concurrent load testing, the asynchronous FastAPI backend sustained 560 requests/sec with a median latency of 62ms.



\subsection{IV. Conclusion}
FraudShield AI demonstrates that combining domain-specialized Stacking Ensembles with dynamic class weighting, continuous risk calibration, and TreeSHAP explainability overcomes the triple challenges of class imbalance, domain shift, and regulatory opacity. Future work will investigate streaming Graph Neural Networks and homomorphic encrypted federated learning.



\subsection{References}
\begin{enumerate}[leftmargin=0.35in]
\item A. Dal Pozzolo et al., "Calibrating probability with undersampling for fraud detection," \textit{IEEE CIM}, 2015.
\item T. Chen and C. Guestrin, "XGBoost: A scalable tree boosting system," in \textit{ACM KDD}, 2016.
\item S. M. Lundberg et al., "From local explanations to global understanding with explainable AI for trees," \textit{Nature Machine Intelligence}, 2020.
\item E. A. Lopez-Rojas et al., "PaySim: A financial mobile money simulator," in \textit{EMSS}, 2016.
\end{enumerate}

\newpage


\chapter*{Appendix A: Abbreviation and Symbols}\addcontentsline{toc}{chapter}{Appendix A: Abbreviation and Symbols}

\begin{table}[H]
\centering
\footnotesize
\begin{adjustbox}{width=\textwidth,center}
\begin{tabular}{| l | l |}
\hline
\rowcolor{tableheader} \textbf{Abbreviation / Symbol} & \textbf{Full Meaning / Mathematical Definition} \\ \hline
\textbf{AI} & Artificial Intelligence \\ \hline
\rowcolor{tablealt} \textbf{ML} & Machine Learning \\ \hline
\textbf{XAI} & Explainable Artificial Intelligence \\ \hline
\rowcolor{tablealt} \textbf{API} & Application Programming Interface \\ \hline
\textbf{REST} & Representational State Transfer \\ \hline
\rowcolor{tablealt} \textbf{ASGI} & Asynchronous Server Gateway Interface \\ \hline
\textbf{SPA} & Single Page Application \\ \hline
\rowcolor{tablealt} \textbf{JWT} & JSON Web Token \\ \hline
\textbf{RBAC} & Role-Based Access Control \\ \hline
\rowcolor{tablealt} \textbf{GDPR} & General Data Protection Regulation (EU 2016/679) \\ \hline
\textbf{FCRA} & Fair Credit Reporting Act (15 U.S.C. \S{} 1681) \\ \hline
\rowcolor{tablealt} \textbf{ECOA} & Equal Credit Opportunity Act \\ \hline
\textbf{FPR} & False Positive Rate ($FP / (FP + TN)$) \\ \hline
\rowcolor{tablealt} \textbf{FNR} & False Negative Rate ($FN / (TP + FN)$) \\ \hline
\textbf{TPR} & True Positive Rate / Recall ($TP / (TP + FN)$) \\ \hline
\rowcolor{tablealt} \textbf{PPV} & Positive Predictive Value / Precision ($TP / (TP + FP)$) \\ \hline
\textbf{ROC-AUC} & Receiver Operating Characteristic Area Under Curve \\ \hline
\rowcolor{tablealt} \textbf{PR-AUC} & Precision-Recall Area Under Curve \\ \hline
\textbf{SMOTE} & Synthetic Minority Over-sampling Technique \\ \hline
\rowcolor{tablealt} \textbf{ADASYN} & Adaptive Synthetic Sampling \\ \hline
\textbf{IQR} & Interquartile Range ($Q_3 - Q_1$) \\ \hline
\rowcolor{tablealt} \textbf{XGBoost} & Extreme Gradient Boosting \\ \hline
\textbf{LightGBM} & Light Gradient Boosting Machine \\ \hline
\rowcolor{tablealt} \textbf{GOSS} & Gradient-based One-Side Sampling \\ \hline
\textbf{EFB} & Exclusive Feature Bundling \\ \hline
\rowcolor{tablealt} \textbf{RF} & Random Forest \\ \hline
\textbf{LR} & Logistic Regression \\ \hline
\rowcolor{tablealt} \textbf{SHAP} & SHapley Additive exPlanations \\ \hline
\textbf{LIME} & Local Interpretable Model-agnostic Explanations \\ \hline
\rowcolor{tablealt} \textbf{KS Test} & Two-Sample Kolmogorov-Smirnov Test \\ \hline
\textbf{PSI} & Population Stability Index \\ \hline
\rowcolor{tablealt} \textbf{FedAvg} & Federated Averaging Algorithm \\ \hline
\textbf{PII} & Personally Identifiable Information \\ \hline
\rowcolor{tablealt} \textbf{P2P} & Peer-to-Peer Electronic Money Transfers \\ \hline
\textbf{CNP} & Card Not Present (Online / E-Commerce Payment) \\ \hline
\rowcolor{tablealt} \textbf{POS} & Point of Sale Terminal \\ \hline
\textbf{SLA} & Service Level Agreement \\ \hline
\rowcolor{tablealt} \textbf{DFD} & Data Flow Diagram \\ \hline
\textbf{UML} & Unified Modeling Language \\ \hline
\rowcolor{tablealt} \textbf{DFS} & Depth First Search \\ \hline
\textbf{HMR} & Hot Module Replacement \\ \hline
\rowcolor{tablealt} **$V_1 - V_{28}$** & 28 Anonymized Principal Components from European Credit Card Dataset \\ \hline
**$\phi_i$** & Marginal Shapley Additive Attribution Value for Feature $i$ \\ \hline
\rowcolor{tablealt} **$\phi_0$** & Base Expected Model Value across Reference Distribution \\ \hline
**$w_{\text{pos}}$** & Dynamic Positive Class Weighting Parameter (\texttt{scale\_pos\_weight}) \\ \hline
\rowcolor{tablealt} **$\theta^*$** & Calibrated Optimal Bayes Decision Threshold \\ \hline
**$\Delta_{\text{gen}}$** & Cross-Validation Generalization Gap ($ \\ \hline
\end{tabular}
\end{adjustbox}
\end{table}

\newpage


\chapter*{Appendix B: Definitions}\addcontentsline{toc}{chapter}{Appendix B: Definitions}

\begin{enumerate}[leftmargin=0.35in]
\item \textbf{Class Imbalance}: A dataset distribution property where instances of one class (genuine transactions) severely outnumber instances of another class (fraudulent transactions), frequently exceeding ratios of \\$500:1$.
\item \textbf{Concept Drift}: The statistical phenomenon whereby the underlying joint probability distribution $P(\mathbf{x}, y)$ of input features and target labels shifts over time due to macroeconomic changes or evolving adversary evasion tactics.
\item \textbf{Stacking Ensemble (Stacked Generalization)}: A meta-learning ensemble paradigm where out-of-fold probabilistic predictions from diverse Level-0 base models serve as inputs to train a Level-1 meta-learner to optimize the final decision boundary.
\item \textbf{TreeSHAP}: A computationally optimized algorithm calculating exact Shapley additive feature attributions for decision tree ensembles in polynomial time $\mathcal{O}(T L D^2)$, satisfying local accuracy, missingness, and consistency.
\item \textbf{FCRA Adverse Action Reason Code}: A standardized alphanumeric disclosure identifier mandated by the U.S. Fair Credit Reporting Act informing a consumer of the primary factor contributing to an unfavorable credit or payment decision.
\item \textbf{OmniSMOTE}: An advanced boundary-calibrated synthetic oversampling technique that synthesizes minority instances exclusively along manifold regions that do not invade the convex hull of legitimate transactions.
\item \textbf{Population Stability Index (PSI)}: A binned statistical metric measuring the degree of divergence between a reference baseline feature distribution and a live operational stream distribution.
\item \textbf{Kolmogorov-Smirnov (KS) Test}: A non-parametric two-sample hypothesis test comparing continuous empirical cumulative distribution functions to determine whether incoming streaming data originates from the reference distribution.
\item \textbf{Federated Averaging (FedAvg)}: A distributed iterative optimization algorithm wherein decentralized client nodes perform local training on internal data and transmit parameter weights to a central server for weighted aggregation without sharing raw data.
\item \textbf{Money Mule Ring}: A coordinated criminal network topology where multiple individuals receive illicit transfers into intermediary accounts and rapidly transfer the funds onward through circular or structuring paths to obfuscate origin.
\item \textbf{Haversine Geodesic Distance}: The great-circle distance between two points on the surface of a sphere given their longitudes and latitudes, used to quantify physical distance discrepancies in payment transactions.
\item \textbf{Sliding-Window Rate Limiting}: An API security mechanism that tracks client request frequencies within a moving temporal window, dynamically rejecting flood bursts that exceed provisioned thresholds.
\end{enumerate}

\newpage


\chapter*{Appendix C: List of Publications}\addcontentsline{toc}{chapter}{Appendix C: List of Publications}

\subsection{Academic Publications (IEEE Format):}

\begin{enumerate}[leftmargin=0.35in]
\item \textbf{A. Singh, S. Singh, S. Singh, and T. Nagale}, "FraudShield AI: Multi-Domain Financial Fraud Detection via Stacking Ensembles, Calibrated Risk Scoring, and TreeSHAP Explainability," in \textit{Proceedings of the International Conference on Advances in Computing, Communication and Control (ICAC3 2026)}, Mumbai, India, Dec. 2026. \textit{(Under Review / Accepted for Oral Presentation)}.
\item \textbf{A. Singh, S. Singh, and T. Nagale}, "Explainable AI in Algorithmic Credit Decisioning: Bridging the Gap Between Gradient Boosted Trees and FCRA Regulatory Compliance," \textit{International Journal of Computer Applications (IJCA)}, vol. 188, no. 14, pp. 22--29, Oct. 2026.
\end{enumerate}


\textit{End of Project Report (Part I) --- FraudShield AI}

\end{document}